\documentclass{beamer}
\usetheme{Madrid}
\usecolortheme{default}
\usepackage[utf8]{inputenc}
\usepackage{amsmath}
\usepackage{amsfonts}
\usepackage{amssymb}
\usepackage{tikz}

\title{CRISC Exam - Hard Questions with Explanations}
\subtitle{From ISACA CRISC QAE Manual 6th Edition}
\author{Compiled from ISACA CRISC Q\&A}
\date{\today}

\begin{document}
	
	\begin{frame}
		\titlepage
	\end{frame}
	
	\begin{frame}{Outline}
		\tableofcontents
	\end{frame}
	
	% ================ SECTION 1: DOMAIN 1 - GOVERNANCE ================
	
	\section{Domain 1: Governance}
	
	\begin{frame}{Q1: Risk Management Strategy}
		\begin{block}{Question}
			Which is MOST important to determine when defining risk management strategies?
		\end{block}
		\begin{enumerate}[(A)]
			\item Risk assessment criteria
			\item IT architecture complexity
			\item Enterprise disaster recovery plan
			\item Business objectives and operations
		\end{enumerate}
	\end{frame}
	
	\begin{frame}{Q1: Explanation}
		\begin{block}{Correct Answer: D}
			While defining risk management strategies, the risk practitioner needs to analyze the enterprise's objectives and risk tolerance and define a risk management framework based on this analysis.
		\end{block}
		\begin{block}{Option Analysis}
			\begin{itemize}
				\item \textbf{A - Incorrect:} Risk assessment criteria alone are not sufficient
				\item \textbf{B - Incorrect:} IT architecture complexity is more related to assessing risk than defining strategies
				\item \textbf{C - Incorrect:} DR plan is more related to mitigating risk
				\item \textbf{D - Correct:} Business objectives drive risk management strategy
			\end{itemize}
		\end{block}
	\end{frame}
	
	\begin{frame}{Q2: Risk Management Strategic Plan}
		\begin{block}{Question}
			Which is MOST important to include in a risk management strategic plan?
		\end{block}
		\begin{enumerate}[(A)]
			\item Risk management staffing requirements
			\item Risk management mission statement
			\item Risk mitigation investment plans
			\item Current state and desired future state
		\end{enumerate}
	\end{frame}
	
	\begin{frame}{Q2: Explanation}
		\begin{block}{Correct Answer: D}
			It is most important to paint a vision for the future and then draw a road map from the starting point, which requires that the current state and desired future state be fully understood.
		\end{block}
		\begin{block}{Option Analysis}
			\begin{itemize}
				\item \textbf{A - Incorrect:} Staffing requirements are driven by understanding current/future state
				\item \textbf{B - Incorrect:} Mission statement is important but not actionable
				\item \textbf{C - Incorrect:} Investment plans are driven by current/future state
				\item \textbf{D - Correct:} Understanding current and desired state is foundational
			\end{itemize}
		\end{block}
	\end{frame}
	
	\begin{frame}{Q3: BU Risk Roles}
		\begin{block}{Question}
			Which BEST describes risk-related roles of an organizational business unit?
		\end{block}
		\begin{enumerate}[(A)]
			\item BU management owns mitigation plan; board identifies risk
			\item BU management owns risk and is responsible for identifying, assessing and mitigating risk
			\item BU carries out responsibilities; ultimate accountability belongs to board
			\item BU is ultimately accountable; board owns the risk
		\end{enumerate}
	\end{frame}
	
	\begin{frame}{Q3: Explanation}
		\begin{block}{Correct Answer: B}
			The BU is responsible for owning the risk and its resulting actions. Risk owners have the responsibility of identifying, measuring, monitoring, controlling and reporting on risk.
		\end{block}
		\begin{block}{Option Analysis}
			\begin{itemize}
				\item \textbf{A - Incorrect:} Board does not perform identification/assessment
				\item \textbf{B - Correct:} BU owns risk and all associated activities
				\item \textbf{C - Incorrect:} BU also has accountability for day-to-day work
				\item \textbf{D - Incorrect:} Board does not own BU risk
			\end{itemize}
		\end{block}
	\end{frame}
	
	\begin{frame}{Q4: Outsourcing Security Consideration}
		\begin{block}{Question}
			Enterprise outsourced IT to a third party whose servers are in a foreign country. Which is MOST critical security consideration?
		\end{block}
		\begin{enumerate}[(A)]
			\item Security breach notification may be delayed due to time difference
			\item Additional network intrusion detection sensors needed
			\item Enterprise could be unable to monitor compliance with internal guidelines
			\item Laws of country of origin may not be enforceable in foreign country
		\end{enumerate}
	\end{frame}
	
	\begin{frame}{Q4: Explanation}
		\begin{block}{Correct Answer: D}
			Laws and regulations of the country of origin may not be enforceable in the foreign country. Conversely, the laws of the foreign vendor may also affect the enterprise.
		\end{block}
		\begin{block}{Option Analysis}
			\begin{itemize}
				\item \textbf{A - Incorrect:} Time difference does not affect 24/7 environments
				\item \textbf{B - Incorrect:} Additional sensors are manageable
				\item \textbf{C - Incorrect:} Outsourcing does not remove internal responsibility
				\item \textbf{D - Correct:} Legal enforceability across borders is critical
			\end{itemize}
		\end{block}
	\end{frame}
	
	\begin{frame}{Q5: Risk Management Process Improvement}
		\begin{block}{Question}
			Which is MOST beneficial to improvement of an enterprise's risk management process?
		\end{block}
		\begin{enumerate}[(A)]
			\item Key risk indicators
			\item External benchmarking
			\item Latest risk assessment
			\item A maturity model
		\end{enumerate}
	\end{frame}
	
	\begin{frame}{Q5: Explanation}
		\begin{block}{Correct Answer: D}
			A maturity model helps identify the status quo and the path to get to the enterprise's desired state and thus is most helpful when an enterprise desires to improve a business process.
		\end{block}
		\begin{block}{Option Analysis}
			\begin{itemize}
				\item \textbf{A - Incorrect:} KRIs help monitor but don't define desired state
				\item \textbf{B - Incorrect:} Benchmarking helps compare but doesn't define path
				\item \textbf{C - Incorrect:} Risk assessment is an input, not a path to improvement
				\item \textbf{D - Correct:} Maturity model provides current state and path to desired state
			\end{itemize}
		\end{block}
	\end{frame}
	
	\begin{frame}{Q6: Independent Risk Review}
		\begin{block}{Question}
			Which is PRIMARY reason for subjecting risk management process to review by independent risk auditors?
		\end{block}
		\begin{enumerate}[(A)]
			\item To ensure risk results are consistent
			\item To ensure risk factors and risk profile are well-defined
			\item To correct any mistakes in risk assessment
			\item To validate control weaknesses for management reporting
		\end{enumerate}
	\end{frame}
	
	\begin{frame}{Q6: Explanation}
		\begin{block}{Correct Answer: B}
			Risk profile and risk factors are defined during the risk assessment process; an independent review helps ensure that the underlying process is effective and helps identify areas for future improvement.
		\end{block}
		\begin{block}{Option Analysis}
			\begin{itemize}
				\item \textbf{A - Incorrect:} Consistency is important but not primary
				\item \textbf{B - Correct:} Independent review validates the process
				\item \textbf{C - Incorrect:} Primary purpose is process improvement, not correction
				\item \textbf{D - Incorrect:} Validating weaknesses may be an outcome but not primary
			\end{itemize}
		\end{block}
	\end{frame}
	
	\begin{frame}{Q7: Evaluating Legal Requirements Impact}
		\begin{block}{Question}
			What is MOST effective method to evaluate potential impact of legal, regulatory and contractual requirements on business objectives?
		\end{block}
		\begin{enumerate}[(A)]
			\item Compliance-oriented gap analysis
			\item Interviews with business process stakeholders
			\item Mapping compliance requirements to policies
			\item Compliance-oriented business impact analysis
		\end{enumerate}
	\end{frame}
	
	\begin{frame}{Q7: Explanation}
		\begin{block}{Correct Answer: D}
			A compliance-oriented business impact analysis will identify compliance requirements to which the enterprise is subject and will assess their effect on business objectives and activities.
		\end{block}
		\begin{block}{Option Analysis}
			\begin{itemize}
				\item \textbf{A - Incorrect:} Gap analysis only identifies gaps, not business impact
				\item \textbf{B - Incorrect:} Interviews identify objectives but not compliance impact
				\item \textbf{C - Incorrect:} Mapping shows how compliance is achieved, not impact
				\item \textbf{D - Correct:} Compliance BIA assesses effect on business objectives
			\end{itemize}
		\end{block}
	\end{frame}
	
	\begin{frame}{Q8: Risk Assessment RFP}
		\begin{block}{Question}
			When hiring consultant to determine risk management program maturity, MOST important element of RFP is:
		\end{block}
		\begin{enumerate}[(A)]
			\item Sample deliverable
			\item Past experience of engagement team
			\item Methodology used in assessment
			\item References from other enterprises
		\end{enumerate}
	\end{frame}
	
	\begin{frame}{Q8: Explanation}
		\begin{block}{Correct Answer: C}
			Methodology illustrates the consultant's process and offers a basis to align expectations with execution of the assessment. Methodology establishes requirements of all parties involved.
		\end{block}
		\begin{block}{Option Analysis}
			\begin{itemize}
				\item \textbf{A - Incorrect:} Sample deliverable only shows presentation, not process
				\item \textbf{B - Incorrect:} Past experience is not as important as methodology
				\item \textbf{C - Correct:} Methodology shows how the assessment will be done
				\item \textbf{D - Incorrect:} References are important but not as important as methodology
			\end{itemize}
		\end{block}
	\end{frame}
	
	\begin{frame}{Q9: New Law Impact}
		\begin{block}{Question}
			Shortly after annual policy review, a new law may affect security requirements for HR system. Risk practitioner should:
		\end{block}
		\begin{enumerate}[(A)]
			\item Analyze what systems and technology-related processes may be impacted
			\item Ensure adjustments are implemented during next review cycle
			\item Initiate ad hoc revision of corporate policy
			\item Notify system custodian to implement changes
		\end{enumerate}
	\end{frame}
	
	\begin{frame}{Q9: Explanation}
		\begin{block}{Correct Answer: A}
			Assessing what systems and technology-related processes may be impacted is the best course of action. The analysis must also determine whether existing controls already address the new requirements.
		\end{block}
		\begin{block}{Option Analysis}
			\begin{itemize}
				\item \textbf{A - Correct:} Assessment should precede any action
				\item \textbf{B - Incorrect:} Significant changes should trigger ad hoc reassessment
				\item \textbf{C - Incorrect:} Ad hoc amendment may be rash and unnecessary
				\item \textbf{D - Incorrect:} Changes need approval from process owner
			\end{itemize}
		\end{block}
	\end{frame}
	
	\begin{frame}{Q10: Outsourced Data Processing Liability}
		\begin{block}{Question}
			Enterprise outsourced personnel data processing and a regulatory violation occurs. Who will be held legally responsible?
		\end{block}
		\begin{enumerate}[(A)]
			\item Supplier, because it has operational responsibility
			\item Enterprise, because it owns the data
			\item Enterprise and supplier
			\item Supplier, because it did not comply with contract
		\end{enumerate}
	\end{frame}
	
	\begin{frame}{Q10: Explanation}
		\begin{block}{Correct Answer: B}
			The enterprise retains responsibility for management of and adherence to policies, procedures and regulatory requirements. The regulatory authority will generally hold the enterprise responsible.
		\end{block}
		\begin{block}{Option Analysis}
			\begin{itemize}
				\item \textbf{A - Incorrect:} Regulatory authority will hold enterprise responsible
				\item \textbf{B - Correct:} Enterprise owns the data and remains responsible
				\item \textbf{C - Incorrect:} Regulatory authority holds enterprise responsible
				\item \textbf{D - Incorrect:} Supplier responsibilities are limited to contract terms
			\end{itemize}
		\end{block}
	\end{frame}
	
	\begin{frame}{Q11: Evaluating Internal Controls}
		\begin{block}{Question}
			Which role is responsible for evaluating effectiveness of existing internal information security controls?
		\end{block}
		\begin{enumerate}[(A)]
			\item Data owner
			\item Senior management
			\item End users
			\item System auditor
		\end{enumerate}
	\end{frame}
	
	\begin{frame}{Q11: Explanation}
		\begin{block}{Correct Answer: D}
			The system auditor is responsible for providing continuous feedback to senior management about the effectiveness of internal controls within the enterprise.
		\end{block}
		\begin{block}{Option Analysis}
			\begin{itemize}
				\item \textbf{A - Incorrect:} Data owner defines business requirements
				\item \textbf{B - Incorrect:} Senior management ensures effectiveness but doesn't evaluate
				\item \textbf{C - Incorrect:} End users can identify flaws but are not responsible
				\item \textbf{D - Correct:} System auditor evaluates control effectiveness
			\end{itemize}
		\end{block}
	\end{frame}
	
	\begin{frame}{Q12: Risk Management Program Effectiveness}
		\begin{block}{Question}
			Which BEST ensures overall effectiveness of a risk management program?
		\end{block}
		\begin{enumerate}[(A)]
			\item Obtaining feedback from all end users
			\item Assigning dedicated risk manager
			\item Applying quantitative risk methodologies
			\item Gaining participation of relevant stakeholders
		\end{enumerate}
	\end{frame}
	
	\begin{frame}{Q12: Explanation}
		\begin{block}{Correct Answer: D}
			Stakeholders directly or indirectly determine the risk response. Without having participation from enterprise stakeholders, the effectiveness of a risk management program cannot be guaranteed.
		\end{block}
		\begin{block}{Option Analysis}
			\begin{itemize}
				\item \textbf{A - Incorrect:} Feedback from all end users is generally not feasible
				\item \textbf{B - Incorrect:} Dedicated manager is less effective without stakeholder involvement
				\item \textbf{C - Incorrect:} Methodology selection depends on enterprise needs
				\item \textbf{D - Correct:} Stakeholder participation is essential
			\end{itemize}
		\end{block}
	\end{frame}
	
	\begin{frame}{Q13: Risk-Based Objectives}
		\begin{block}{Question}
			Which approach BEST helps an enterprise achieve risk-based organizational objectives?
		\end{block}
		\begin{enumerate}[(A)]
			\item Ensure asset owners perform annual risk assessments
			\item Review and update risk register regularly
			\item Assign steering committee to risk management
			\item Embed risk management activities into business processes
		\end{enumerate}
	\end{frame}
	
	\begin{frame}{Q13: Explanation}
		\begin{block}{Correct Answer: D}
			The primary objective of embedding risk management activities into business processes is to achieve risk-based organizational objectives in the most effective manner possible.
		\end{block}
		\begin{block}{Option Analysis}
			\begin{itemize}
				\item \textbf{A - Incorrect:} Assessment alone does not achieve objectives
				\item \textbf{B - Incorrect:} Maintaining risk register tracks risk but doesn't achieve objectives
				\item \textbf{C - Incorrect:} Steering committee aids alignment but isn't as effective
				\item \textbf{D - Correct:} Embedding into processes achieves objectives effectively
			\end{itemize}
		\end{block}
	\end{frame}
	
	\begin{frame}{Q14: Regulatory Body Reliance}
		\begin{block}{Question}
			When assessing capability of risk management process, regulatory body would place GREATEST reliance on:
		\end{block}
		\begin{enumerate}[(A)]
			\item Peer review
			\item Internal review
			\item External review
			\item Process capability review
		\end{enumerate}
	\end{frame}
	
	\begin{frame}{Q14: Explanation}
		\begin{block}{Correct Answer: C}
			Regulatory entities generally use assessments performed by an objective and independent third party. Of the choices presented, an external review is the most objective and independent.
		\end{block}
		\begin{block}{Option Analysis}
			\begin{itemize}
				\item \textbf{A - Incorrect:} Peer review has less independence
				\item \textbf{B - Incorrect:} Internal review may be subject to management influence
				\item \textbf{C - Correct:} External review is most objective and independent
				\item \textbf{D - Incorrect:} Doesn't indicate level of independence
			\end{itemize}
		\end{block}
	\end{frame}
	
	\begin{frame}{Q15: Evaluating Current Risk Profile}
		\begin{block}{Question}
			Company improving security program since last review one year ago. What should company do to evaluate current risk profile?
		\end{block}
		\begin{enumerate}[(A)]
			\item Review previous findings and ensure all issues resolved
			\item Conduct follow-up audits in deficient areas
			\item Monitor KRIs and use results for targeted assessments
			\item Perform new enterprise risk assessment using independent expert
		\end{enumerate}
	\end{frame}
	
	\begin{frame}{Q15: Explanation}
		\begin{block}{Correct Answer: D}
			The best way to ensure that an enterprise's security posture is still within compliance is to conduct another risk assessment. Using an independent expert can provide more objective results.
		\end{block}
		\begin{block}{Option Analysis}
			\begin{itemize}
				\item \textbf{A - Incorrect:} New assessment indicates control effectiveness
				\item \textbf{B - Incorrect:} Changes may affect other areas; enterprise-wide assessment is better
				\item \textbf{C - Incorrect:} KRIs are one of many tools; not sufficient alone
				\item \textbf{D - Correct:} New risk assessment with independent expert is best
			\end{itemize}
		\end{block}
	\end{frame}
	
	\begin{frame}{Q16: Greatest Risk of Inadequate Policy}
		\begin{block}{Question}
			Which is GREATEST risk of a policy that defines data and system ownership inadequately?
		\end{block}
		\begin{enumerate}[(A)]
			\item Audit recommendations may not be implemented
			\item Users may have unauthorized access to create, modify or delete data
			\item User management coordination does not exist
			\item Specific user accountability cannot be established
		\end{enumerate}
	\end{frame}
	
	\begin{frame}{Q16: Explanation}
		\begin{block}{Correct Answer: B}
			Without a policy defining who grants access to specific data or systems, the risk that employees receive system access without justified business purpose increases. Business objectives are best supported when authority to grant access is assigned to specific individuals.
		\end{block}
		\begin{block}{Option Analysis}
			\begin{itemize}
				\item \textbf{A - Incorrect:} Audit reports assign remediation owners
				\item \textbf{B - Correct:} Inappropriate access is the greatest risk
				\item \textbf{C - Incorrect:} This is a risk but not the greatest
				\item \textbf{D - Incorrect:} User accountability is established by unique IDs
			\end{itemize}
		\end{block}
	\end{frame}
	
	\begin{frame}{Q17: Management Role in Risk Strategy}
		\begin{block}{Question}
			Which BEST describes role of management in implementing a risk management strategy?
		\end{block}
		\begin{enumerate}[(A)]
			\item Ensure planning, budgeting and performance of IS components are appropriate
			\item Assess and incorporate results of risk management into decision-making
			\item Identify, evaluate and minimize risk to IT systems
			\item Understand risk management process to develop training materials
		\end{enumerate}
	\end{frame}
	
	\begin{frame}{Q17: Explanation}
		\begin{block}{Correct Answer: B}
			Assessing and incorporating the results of risk management into the decision-making process best describes the role of senior management in establishing and implementing a risk management strategy.
		\end{block}
		\begin{block}{Option Analysis}
			\begin{itemize}
				\item \textbf{A - Incorrect:} This is CIO responsibility
				\item \textbf{B - Correct:} Management incorporates risk into decisions
				\item \textbf{C - Incorrect:} This is IT security manager responsibility
				\item \textbf{D - Incorrect:} This is corporate security trainer responsibility
			\end{itemize}
		\end{block}
	\end{frame}
	
	\begin{frame}{Q18: Primary Focus of IT Risk Management}
		\begin{block}{Question}
			The PRIMARY focus of managing IT-related business risk is to protect:
		\end{block}
		\begin{enumerate}[(A)]
			\item Information
			\item Hardware
			\item Applications
			\item Databases
		\end{enumerate}
	\end{frame}
	
	\begin{frame}{Q18: Explanation}
		\begin{block}{Correct Answer: A}
			The primary objective for any enterprise is to protect mission-critical information based on a risk assessment.
		\end{block}
		\begin{block}{Option Analysis}
			\begin{itemize}
				\item \textbf{A - Correct:} Information is the primary focus
				\item \textbf{B - Incorrect:} Hardware is focus only if it stores/processes critical data
				\item \textbf{C - Incorrect:} Applications are focus only if they process critical data
				\item \textbf{D - Incorrect:} Databases are focus only if they store critical data
			\end{itemize}
		\end{block}
	\end{frame}
	
	\begin{frame}{Q19: Best Risk Management Perspective}
		\begin{block}{Question}
			Which can provide BEST perspective of risk management to employees and stockholders?
		\end{block}
		\begin{enumerate}[(A)]
			\item An interdisciplinary team within the enterprise
			\item A third-party risk assessment service provider
			\item The enterprise's IT department
			\item The enterprise's internal compliance department
		\end{enumerate}
	\end{frame}
	
	\begin{frame}{Q19: Explanation}
		\begin{block}{Correct Answer: A}
			Assembling an interdisciplinary team to manage risk ensures that all areas are adequately considered in risk assessment and helps provide an enterprise-wide perspective on risk.
		\end{block}
		\begin{block}{Option Analysis}
			\begin{itemize}
				\item \textbf{A - Correct:} Interdisciplinary team provides enterprise-wide view
				\item \textbf{B - Incorrect:} Lacks internal knowledge and judgment
				\item \textbf{C - Incorrect:} IT department view is too narrow
				\item \textbf{D - Incorrect:} Compliance department focuses on regulatory compliance
			\end{itemize}
		\end{block}
	\end{frame}
	
	\begin{frame}{Q20: Global Policy Approach}
		\begin{block}{Question}
			Which approach to corporate policy BEST supports expansion to other regions with different local laws?
		\end{block}
		\begin{enumerate}[(A)]
			\item Global policy without provisions disputed at local levels
			\item Global policy amended to comply with local laws
			\item Global policy that complies with laws at headquarters
			\item Local policies to accommodate laws within each region
		\end{enumerate}
	\end{frame}
	
	\begin{frame}{Q20: Explanation}
		\begin{block}{Correct Answer: B}
			A global policy including local amendments ensures alignment with local laws and regulations.
		\end{block}
		\begin{block}{Option Analysis}
			\begin{itemize}
				\item \textbf{A - Incorrect:} One policy for all locales is nearly impossible
				\item \textbf{B - Correct:} Global policy with local amendments ensures compliance
				\item \textbf{C - Incorrect:} Exposes enterprise to legal action and reputational loss
				\item \textbf{D - Incorrect:} Extremely expensive and may fail to leverage common practices
			\end{itemize}
		\end{block}
	\end{frame}
	
	\begin{frame}{Q21: Data Classification Responsibility}
		\begin{block}{Question}
			Who is MOST likely responsible for data classification?
		\end{block}
		\begin{enumerate}[(A)]
			\item Data user
			\item Data owner
			\item Data custodian
			\item System administrator
		\end{enumerate}
	\end{frame}
	
	\begin{frame}{Q21: Explanation}
		\begin{block}{Correct Answer: B}
			The data owner is responsible for classifying data according to the enterprise's data classification scheme. The classification scheme then defines who is eligible to access the data and what controls are required.
		\end{block}
		\begin{block}{Option Analysis}
			\begin{itemize}
				\item \textbf{A - Incorrect:} Data user gains access based on business need
				\item \textbf{B - Correct:} Data owner classifies data
				\item \textbf{C - Incorrect:} Data custodian implements controls
				\item \textbf{D - Incorrect:} System administrator is a custodian
			\end{itemize}
		\end{block}
	\end{frame}
	
	\begin{frame}{Q22: Security Governance Model Factor}
		\begin{block}{Question}
			Which factor will have GREATEST impact on type of information security governance model adopted?
		\end{block}
		\begin{enumerate}[(A)]
			\item Number of employees
			\item Enterprise's budget
			\item Organizational structure
			\item Type of technology used
		\end{enumerate}
	\end{frame}
	
	\begin{frame}{Q22: Explanation}
		\begin{block}{Correct Answer: C}
			Information security governance models depend significantly on the overall organizational structure.
		\end{block}
		\begin{block}{Option Analysis}
			\begin{itemize}
				\item \textbf{A - Incorrect:} Well-defined processes provide proper governance
				\item \textbf{B - Incorrect:} Budget does not dictate governance model
				\item \textbf{C - Correct:} Organizational structure is the primary factor
				\item \textbf{D - Incorrect:} Well-defined processes provide proper governance
			\end{itemize}
		\end{block}
	\end{frame}
	
	\begin{frame}{Q23: Information Security Procedures}
		\begin{block}{Question}
			Information security procedures should:
		\end{block}
		\begin{enumerate}[(A)]
			\item Be updated frequently as new software is released
			\item Underline the importance of security governance
			\item Define the allowable limits of behavior
			\item Describe security baselines for each platform
		\end{enumerate}
	\end{frame}
	
	\begin{frame}{Q23: Explanation}
		\begin{block}{Correct Answer: A}
			Security procedures have to change frequently to keep up with changes in software. Because a procedure is a how-to document, it must be kept current with frequent changes in software.
		\end{block}
		\begin{block}{Option Analysis}
			\begin{itemize}
				\item \textbf{A - Correct:} Procedures must be updated with software changes
				\item \textbf{B - Incorrect:} High-level objectives are addressed in policy
				\item \textbf{C - Incorrect:} Security policies define behavioral limits
				\item \textbf{D - Incorrect:} Security standards define platform baselines
			\end{itemize}
		\end{block}
	\end{frame}
	
	\begin{frame}{Q24: Greatest Benefit of Risk-Aware Culture}
		\begin{block}{Question}
			Which is GREATEST benefit of a risk-aware culture?
		\end{block}
		\begin{enumerate}[(A)]
			\item Issues are escalated when suspicious activity is noticed
			\item Controls are double-checked to anticipate issues
			\item Individuals communicate with peers for knowledge sharing
			\item Employees are self-motivated to learn about costs and benefits
		\end{enumerate}
	\end{frame}
	
	\begin{frame}{Q24: Explanation}
		\begin{block}{Correct Answer: A}
			Management benefits most from an escalation process because risk and incidents are reported in a timely manner. Escalation posture among employees is best developed through training and awareness programs.
		\end{block}
		\begin{block}{Option Analysis}
			\begin{itemize}
				\item \textbf{A - Correct:} Timely escalation benefits management most
				\item \textbf{B - Incorrect:} Double-checking is basic business practice
				\item \textbf{C - Incorrect:} Knowledge sharing is indirect benefit
				\item \textbf{D - Incorrect:} Cost-benefit learning is not primary expectation
			\end{itemize}
		\end{block}
	\end{frame}
	
	\begin{frame}{Q25: Main Objective of IT Risk Management}
		\begin{block}{Question}
			The MAIN objective of IT risk management is to:
		\end{block}
		\begin{enumerate}[(A)]
			\item Prevent loss of IT assets
			\item Provide timely management reports
			\item Ensure regulatory compliance
			\item Enable risk-aware business decisions
		\end{enumerate}
	\end{frame}
	
	\begin{frame}{Q25: Explanation}
		\begin{block}{Correct Answer: D}
			IT risk management should be conducted as part of enterprise-wide risk management, whose ultimate objective is to support risk-aware business decisions.
		\end{block}
		\begin{block}{Option Analysis}
			\begin{itemize}
				\item \textbf{A - Incorrect:} Protecting assets is a subordinate goal
				\item \textbf{B - Incorrect:} Reporting is a subordinate goal
				\item \textbf{C - Incorrect:} Compliance is one objective, not the main
				\item \textbf{D - Correct:} Risk-aware business decisions is the ultimate objective
			\end{itemize}
		\end{block}
	\end{frame}
	
	\begin{frame}{Q26: Global Compliance Approach}
		\begin{block}{Question}
			Which is BEST approach for global enterprise subject to multiple governmental jurisdictions with differing requirements?
		\end{block}
		\begin{enumerate}[(A)]
			\item Bring all locations into conformity with aggregate requirements
			\item Bring all locations into conformity with industry good practices
			\item Establish baseline standard incorporating common requirements
			\item Establish baseline standards and add supplemental standards as required
		\end{enumerate}
	\end{frame}
	
	\begin{frame}{Q26: Explanation}
		\begin{block}{Correct Answer: D}
			It is more efficient to establish a baseline standard and then develop additional standards for locations that must meet specific requirements.
		\end{block}
		\begin{block}{Option Analysis}
			\begin{itemize}
				\item \textbf{A - Incorrect:} Lowest common denominator may cause non-compliance
				\item \textbf{B - Incorrect:} Good practices may not meet all regulatory requirements
				\item \textbf{C - Incorrect:} Forcing all locations to comply places undue burden
				\item \textbf{D - Correct:} Baseline plus supplements is most efficient
			\end{itemize}
		\end{block}
	\end{frame}
	
	\begin{frame}{Q27: Risk Management Plan Sign-off}
		\begin{block}{Question}
			Who MUST give final sign-off on the IT risk management plan?
		\end{block}
		\begin{enumerate}[(A)]
			\item IT auditors performing risk assessment
			\item Business process owners
			\item Senior managers
			\item IT security administrators
		\end{enumerate}
	\end{frame}
	
	\begin{frame}{Q27: Explanation}
		\begin{block}{Correct Answer: C}
			Senior managers understand performance metrics and indicators that measure the enterprise and its subsystems; they approved the policies and standards that govern the enterprise, and they have final responsibility for risk associated with audit findings and recommendations.
		\end{block}
		\begin{block}{Option Analysis}
			\begin{itemize}
				\item \textbf{A - Incorrect:} Auditors contribute but lack authority
				\item \textbf{B - Incorrect:} Process owners contribute but lack authority
				\item \textbf{C - Correct:} Senior managers have final responsibility
				\item \textbf{D - Incorrect:} Administrators contribute but lack authority
			\end{itemize}
		\end{block}
	\end{frame}
	
	\begin{frame}{Q28: Security Boundary Purpose}
		\begin{block}{Question}
			What is PRIMARY reason a risk practitioner determines the security boundary prior to conducting a risk assessment?
		\end{block}
		\begin{enumerate}[(A)]
			\item To decide which laws and regulations apply
			\item To identify the scope of the risk assessment
			\item To identify the business owners of the system
			\item To decide whether quantitative or qualitative analysis is appropriate
		\end{enumerate}
	\end{frame}
	
	\begin{frame}{Q28: Explanation}
		\begin{block}{Correct Answer: B}
			Identifying the security boundary establishes the fundamental scope of inquiry, including the systems and components subject to assessment and those not subject to assessment.
		\end{block}
		\begin{block}{Option Analysis}
			\begin{itemize}
				\item \textbf{A - Incorrect:} Risk assessment itself considers laws/regulations
				\item \textbf{B - Correct:} Security boundary establishes scope
				\item \textbf{C - Incorrect:} Identifying business owners is secondary
				\item \textbf{D - Incorrect:} Security boundaries don't directly inform methodology choice
			\end{itemize}
		\end{block}
	\end{frame}
	
	\begin{frame}{Q29: Most Effective Risk Program Management}
		\begin{block}{Question}
			Which group would be MOST effective in managing and executing an enterprise's risk program?
		\end{block}
		\begin{enumerate}[(A)]
			\item Mid-level managers
			\item Senior managers
			\item Frontline employees
			\item Incident response team
		\end{enumerate}
	\end{frame}
	
	\begin{frame}{Q29: Explanation}
		\begin{block}{Correct Answer: A}
			Mid-level managers are the best to manage and execute an enterprise's risk management program because they are the most centrally located within the organizational hierarchy, and they combine a sufficient breadth of influence with adequate proximity to day-to-day operations.
		\end{block}
		\begin{block}{Option Analysis}
			\begin{itemize}
				\item \textbf{A - Correct:} Mid-level managers combine influence with operational proximity
				\item \textbf{B - Incorrect:} Senior managers are too high organizationally
				\item \textbf{C - Incorrect:} Frontline employees lack power and influence
				\item \textbf{D - Incorrect:} Incident response team only manages a small portion
			\end{itemize}
		\end{block}
	\end{frame}
	
	\begin{frame}{Q30: IT Steering Committee Composition}
		\begin{block}{Question}
			IT steering committee will be BEST represented by:
		\end{block}
		\begin{enumerate}[(A)]
			\item Members of the executive board
			\item High-level members of the IT department
			\item IT experts from outside the enterprise
			\item Key members from each department
		\end{enumerate}
	\end{frame}
	
	\begin{frame}{Q30: Explanation}
		\begin{block}{Correct Answer: D}
			The IT steering committee should be comprised of individuals from each department to ensure that the entire enterprise is represented and that all business objectives are more likely to be met.
		\end{block}
		\begin{block}{Option Analysis}
			\begin{itemize}
				\item \textbf{A - Incorrect:} Executive board only focuses on high-level goals
				\item \textbf{B - Incorrect:} IT only ignores business objectives
				\item \textbf{C - Incorrect:} External experts ignore business objectives
				\item \textbf{D - Correct:} Cross-departmental representation ensures business alignment
			\end{itemize}
		\end{block}
	\end{frame}
	
	\begin{frame}{Q31: First Line of Defense Purpose}
		\begin{block}{Question}
			Which is one of the MAIN purposes of the first line of defense in the three lines of defense model?
		\end{block}
		\begin{enumerate}[(A)]
			\item Ensure financial controls are in place
			\item Ensure control deficiencies are addressed
			\item Ensure risk management practices are effective
			\item Ensure compliance with rules and regulations
		\end{enumerate}
	\end{frame}
	
	\begin{frame}{Q31: Explanation}
		\begin{block}{Correct Answer: B}
			In the first line, operational managers own the risk and address it by mitigating control deficiencies.
		\end{block}
		\begin{block}{Option Analysis}
			\begin{itemize}
				\item \textbf{A - Incorrect:} Financial risk is only one category
				\item \textbf{B - Correct:} First line addresses control deficiencies
				\item \textbf{C - Incorrect:} Risk management is second line
				\item \textbf{D - Incorrect:} Ensuring compliance is second line
			\end{itemize}
		\end{block}
	\end{frame}
	
	\begin{frame}{Q32: Risk Response Selection}
		\begin{block}{Question}
			The PRIMARY consideration when selecting a risk response technique is:
		\end{block}
		\begin{enumerate}[(A)]
			\item Coverage of all identified risk
			\item Availability of resources
			\item Enterprise goals and objectives
			\item Standards and industry good practices
		\end{enumerate}
	\end{frame}
	
	\begin{frame}{Q32: Explanation}
		\begin{block}{Correct Answer: C}
			The risk response will be based primarily on goals and objectives of the enterprise. Risk can harm these goals and must be mitigated according to priority.
		\end{block}
		\begin{block}{Option Analysis}
			\begin{itemize}
				\item \textbf{A - Incorrect:} Address mission-critical risk, not all risk
				\item \textbf{B - Incorrect:} Resources are secondary to goals
				\item \textbf{C - Correct:} Enterprise goals drive risk response
				\item \textbf{D - Incorrect:} Standards are followed after mitigation technique is selected
			\end{itemize}
		\end{block}
	\end{frame}
	
	\begin{frame}{Q33: First Policy to Create}
		\begin{block}{Question}
			CIO asked to create all IT policies and procedures. Which should the CIO create FIRST?
		\end{block}
		\begin{enumerate}[(A)]
			\item Strategic IT plan
			\item Data classification scheme
			\item Information architecture document
			\item Technology infrastructure plan
		\end{enumerate}
	\end{frame}
	
	\begin{frame}{Q33: Explanation}
		\begin{block}{Correct Answer: A}
			The strategic IT plan is the first policy to create when developing an enterprise's governance model.
		\end{block}
		\begin{block}{Option Analysis}
			\begin{itemize}
				\item \textbf{A - Correct:} Strategic plan comes first
				\item \textbf{B - Incorrect:} Strategic plan comes before data classification
				\item \textbf{C - Incorrect:} Strategic plan comes before information architecture
				\item \textbf{D - Incorrect:} Strategic plan comes before technology infrastructure plan
			\end{itemize}
		\end{block}
	\end{frame}
	
	\begin{frame}{Q34: Greatest Concern for Risk Practitioner}
		\begin{block}{Question}
			Which should be of MOST concern to a risk practitioner?
		\end{block}
		\begin{enumerate}[(A)]
			\item Failure to notify public of intrusion
			\item Failure to notify police of attempted intrusion
			\item Failure to internally report a successful attack
			\item Failure to examine access rights periodically
		\end{enumerate}
	\end{frame}
	
	\begin{frame}{Q34: Explanation}
		\begin{block}{Correct Answer: C}
			Failure to report a successful intrusion is a serious concern to the risk practitioner and could—in some instances—be interpreted as abetting.
		\end{block}
		\begin{block}{Option Analysis}
			\begin{itemize}
				\item \textbf{A - Incorrect:} Reporting to public is not a requirement
				\item \textbf{B - Incorrect:} Attempted intrusion rarely requires police notification
				\item \textbf{C - Correct:} Failure to report successful attack is serious
				\item \textbf{D - Incorrect:} Lack of periodic review is a concern but not as serious
			\end{itemize}
		\end{block}
	\end{frame}
	
	\begin{frame}{Q35: Risk Management Capabilities Insight}
		\begin{block}{Question}
			Which review will provide MOST insight into enterprise's risk management capabilities?
		\end{block}
		\begin{enumerate}[(A)]
			\item Capability maturity model review
			\item Capability comparison with industry standards
			\item Self-assessment of capabilities
			\item Internal audit review of capabilities
		\end{enumerate}
	\end{frame}
	
	\begin{frame}{Q35: Explanation}
		\begin{block}{Correct Answer: A}
			Capability maturity modeling allows an enterprise to understand its level of maturity in its risk capabilities, which is an indicator of operational readiness and effectiveness.
		\end{block}
		\begin{block}{Option Analysis}
			\begin{itemize}
				\item \textbf{A - Correct:} Maturity model shows readiness and effectiveness
				\item \textbf{B - Incorrect:} Only shows existence/non-existence of capabilities
				\item \textbf{C - Incorrect:} Only shows existence/non-existence
				\item \textbf{D - Incorrect:} Only shows existence/non-existence
			\end{itemize}
		\end{block}
	\end{frame}
	
	\begin{frame}{Q36: Leveraging Internal Audit}
		\begin{block}{Question}
			Risk practitioner BEST leverages work performed by internal audit by having it:
		\end{block}
		\begin{enumerate}[(A)]
			\item Design, implement and maintain ERM process
			\item Manage and assess overall risk awareness
			\item Evaluate ongoing changes to enterprise risk factors
			\item Assist in monitoring, evaluating, examining and reporting on controls
		\end{enumerate}
	\end{frame}
	
	\begin{frame}{Q36: Explanation}
		\begin{block}{Correct Answer: D}
			The internal audit function is responsible for assisting management and the board of directors in monitoring, evaluating, examining and reporting on internal controls, regardless of whether an ERM function has been implemented.
		\end{block}
		\begin{block}{Option Analysis}
			\begin{itemize}
				\item \textbf{A - Incorrect:} Design/implementation is management responsibility
				\item \textbf{B - Incorrect:} Risk awareness is risk governance responsibility
				\item \textbf{C - Incorrect:} Evaluating changes is not internal audit responsibility
				\item \textbf{D - Correct:} Internal audit assists in monitoring and reporting on controls
			\end{itemize}
		\end{block}
	\end{frame}
	
	\begin{frame}{Q37: Biggest Concern for CISO}
		\begin{block}{Question}
			Which is BIGGEST concern for CISO regarding interconnections with systems outside the enterprise?
		\end{block}
		\begin{enumerate}[(A)]
			\item Requirements to comply with each other's contractual security obligations
			\item Uncertainty that the other system will be available as needed
			\item The ability to perform risk assessments on the other system
			\item Ensuring communications are encrypted through VPN tunnel
		\end{enumerate}
	\end{frame}
	
	\begin{frame}{Q37: Explanation}
		\begin{block}{Correct Answer: A}
			Ensuring that both systems comply with mutual contractual security obligations should be the primary concern of the risk practitioner. If one system fails to comply, both will likely miss their respective security obligations.
		\end{block}
		\begin{block}{Option Analysis}
			\begin{itemize}
				\item \textbf{A - Correct:} Contractual compliance is primary concern
				\item \textbf{B - Incorrect:} Availability is business owner's concern
				\item \textbf{C - Incorrect:} May or may not be a concern based on agreement
				\item \textbf{D - Incorrect:} Encryption depends on type of data transmitted
			\end{itemize}
		\end{block}
	\end{frame}
	
	\begin{frame}{Q38: IT Organizational Structure Type}
		\begin{block}{Question}
			CIO creates all IT policies managed by IT steering committee making all IT decisions. Which type of IT organizational structure?
		\end{block}
		\begin{enumerate}[(A)]
			\item Project-based
			\item Centralized
			\item Decentralized
			\item Divisional
		\end{enumerate}
	\end{frame}
	
	\begin{frame}{Q38: Explanation}
		\begin{block}{Correct Answer: B}
			Within a centralized IT organizational structure, one group makes all decisions for the entire enterprise.
		\end{block}
		\begin{block}{Option Analysis}
			\begin{itemize}
				\item \textbf{A - Incorrect:} Project-based groups are temporary
				\item \textbf{B - Correct:} Centralized has one decision-making group
				\item \textbf{C - Incorrect:} Decentralized decisions made by each division
				\item \textbf{D - Incorrect:} Divisional has groups by geography/product
			\end{itemize}
		\end{block}
	\end{frame}
	
	\begin{frame}{Q39: Aggregated Monitoring Results}
		\begin{block}{Question}
			Aggregated results of continuous monitoring activities are BEST communicated to:
		\end{block}
		\begin{enumerate}[(A)]
			\item Risk owner
			\item Technical staff
			\item Audit department
			\item Information security manager
		\end{enumerate}
	\end{frame}
	
	\begin{frame}{Q39: Explanation}
		\begin{block}{Correct Answer: A}
			The risk owner is the most suitable target audience for aggregated results of continuous monitoring; the risk owner is accountable for the fact that appropriate risk responses are executed in alignment with the enterprise's risk appetite.
		\end{block}
		\begin{block}{Option Analysis}
			\begin{itemize}
				\item \textbf{A - Correct:} Risk owner is accountable
				\item \textbf{B - Incorrect:} Technical staff need granular information
				\item \textbf{C - Incorrect:} Audit may review but is not primary recipient
				\item \textbf{D - Incorrect:} Security manager is secondary audience
			\end{itemize}
		\end{block}
	\end{frame}
	
	\begin{frame}{Q40: Risk Awareness Program}
		\begin{block}{Question}
			Which is PRIMARY consideration when developing an IT risk awareness program?
		\end{block}
		\begin{enumerate}[(A)]
			\item Why technology risk is owned by IT
			\item How technology risk can affect each attendee's area of business
			\item How business process owners can transfer technology risk
			\item Why technology risk is more difficult to manage than other risk
		\end{enumerate}
	\end{frame}
	
	\begin{frame}{Q40: Explanation}
		\begin{block}{Correct Answer: B}
			Stakeholders must understand how IT-related risk affects the overall business and how potential threats and vulnerabilities can jeopardize the enterprise's people, processes and technology.
		\end{block}
		\begin{block}{Option Analysis}
			\begin{itemize}
				\item \textbf{A - Incorrect:} IT does not own technology risk
				\item \textbf{B - Correct:} Business impact is primary
				\item \textbf{C - Incorrect:} Transferring risk is part of risk response
				\item \textbf{D - Incorrect:} This may or may not be true
			\end{itemize}
		\end{block}
	\end{frame}
	
	\begin{frame}{Q41: IT Success Against Business Requirements}
		\begin{block}{Question}
			Management wants to ensure IT is successful in delivering against business requirements. Which BEST supports that effort?
		\end{block}
		\begin{enumerate}[(A)]
			\item An internal control system or framework
			\item A cost-benefit analysis
			\item A return on investment analysis
			\item A benchmark process
		\end{enumerate}
	\end{frame}
	
	\begin{frame}{Q41: Explanation}
		\begin{block}{Correct Answer: A}
			For IT to be successful in delivering against business requirements, management should develop an internal control system that supports its business requirements.
		\end{block}
		\begin{block}{Option Analysis}
			\begin{itemize}
				\item \textbf{A - Correct:} Internal control system supports business alignment
				\item \textbf{B - Incorrect:} Cost-benefit is useful but not most important
				\item \textbf{C - Incorrect:} ROI is just one metric
				\item \textbf{D - Incorrect:} Benchmarking comes after sound internal control framework
			\end{itemize}
		\end{block}
	\end{frame}
	
	\begin{frame}{Q42: Risk Appetite Alignment}
		\begin{block}{Question}
			It is MOST important that risk appetite be aligned with business objectives to ensure that:
		\end{block}
		\begin{enumerate}[(A)]
			\item Resources are directed toward areas of low risk tolerance
			\item Major risk is identified and eliminated
			\item IT and business goals are aligned
			\item The risk strategy is adequately communicated
		\end{enumerate}
	\end{frame}
	
	\begin{frame}{Q42: Explanation}
		\begin{block}{Correct Answer: A}
			Risk appetite refers to the amount of risk that an enterprise is willing to take on in pursuit of value. Aligning it with business objectives allows an enterprise to evaluate and deploy valuable resources toward those objectives with low risk tolerance for loss.
		\end{block}
		\begin{block}{Option Analysis}
			\begin{itemize}
				\item \textbf{A - Correct:} Resources directed to low tolerance areas
				\item \textbf{B - Incorrect:} Risk cannot be eliminated
				\item \textbf{C - Incorrect:} Alignment affects more than IT and business
				\item \textbf{D - Incorrect:} Communication does not depend on alignment
			\end{itemize}
		\end{block}
	\end{frame}
	
	\begin{frame}{Q43: Outdated Security Policy}
		\begin{block}{Question}
			Which causes GREATEST concern reviewing a corporate information security policy that is out of date?
		\end{block}
		\begin{enumerate}[(A)]
			\item Was not reviewed within the last three years
			\item Is missing newer technologies/platforms
			\item Was not updated for new locations
			\item Does not enforce control monitoring
		\end{enumerate}
	\end{frame}
	
	\begin{frame}{Q43: Explanation}
		\begin{block}{Correct Answer: A}
			Not reviewing the policy for three years and updating it as necessary does not follow good practices and is the greatest concern.
		\end{block}
		\begin{block}{Option Analysis}
			\begin{itemize}
				\item \textbf{A - Correct:} Not reviewing for three years is greatest concern
				\item \textbf{B - Incorrect:} Corporate policies generally don't require modification for specific technologies
				\item \textbf{C - Incorrect:} Well-written policy accommodates multiple locations
				\item \textbf{D - Incorrect:} Lack of control monitoring is a concern but not as great as policy review
			\end{itemize}
		\end{block}
	\end{frame}
	
	\begin{frame}{Q44: Accountability for Business Risk}
		\begin{block}{Question}
			Who is accountable for business risk related to IT?
		\end{block}
		\begin{enumerate}[(A)]
			\item Chief information officer
			\item Chief financial officer
			\item Users of IT services
			\item Chief architect
		\end{enumerate}
	\end{frame}
	
	\begin{frame}{Q44: Explanation}
		\begin{block}{Correct Answer: C}
			Ultimately, the enterprise (i.e., the users of IT services) owns business-related risk, including the risk related to the use of IT. The business should set the mandate for risk management, provide the resources and funding to support a risk management plan.
		\end{block}
		\begin{block}{Option Analysis}
			\begin{itemize}
				\item \textbf{A - Incorrect:} CIO supports business but doesn't own risk
				\item \textbf{B - Incorrect:} CFO tracks costs but doesn't own risk
				\item \textbf{C - Correct:} Enterprise users own business risk
				\item \textbf{D - Incorrect:} Chief architect mitigates risk but doesn't own it
			\end{itemize}
		\end{block}
	\end{frame}
	
	\begin{frame}{Q45: Expansion Risk Reporting}
		\begin{block}{Question}
			Enterprise expanding into new nearby domestic locations. Which is MOST important to report?
		\end{block}
		\begin{enumerate}[(A)]
			\item Competitor analysis
			\item Legal and regulatory requirements
			\item Political issues
			\item Potential of natural disasters
		\end{enumerate}
	\end{frame}
	
	\begin{frame}{Q45: Explanation}
		\begin{block}{Correct Answer: B}
			Legal and regulatory requirements are most likely to change when moving to a nearby location because each municipality may enforce significantly different regulations, including environmental requirements, taxation and others.
		\end{block}
		\begin{block}{Option Analysis}
			\begin{itemize}
				\item \textbf{A - Incorrect:} Competitors likely won't change
				\item \textbf{B - Correct:} Legal requirements vary by municipality
				\item \textbf{C - Incorrect:} Political issues unlikely to change
				\item \textbf{D - Incorrect:} Natural disasters unlikely to change
			\end{itemize}
		\end{block}
	\end{frame}
	
	\begin{frame}{Q46: True Statement about IT Risk}
		\begin{block}{Question}
			Which is true about IT risk?
		\end{block}
		\begin{enumerate}[(A)]
			\item IT risk cannot be assessed and measured quantitatively
			\item IT risk should be calculated separately from business risk
			\item IT risk management is the responsibility of the IT department
			\item IT risk exists regardless of whether it is detected or recognized
		\end{enumerate}
	\end{frame}
	
	\begin{frame}{Q46: Explanation}
		\begin{block}{Correct Answer: D}
			The enterprise must identify, acknowledge and respond to risk; ignorance of risk is not acceptable.
		\end{block}
		\begin{block}{Option Analysis}
			\begin{itemize}
				\item \textbf{A - Incorrect:} IT risk can be assessed both quantitatively and qualitatively
				\item \textbf{B - Incorrect:} IT risk is one type of business risk
				\item \textbf{C - Incorrect:} IT risk is responsibility of senior management
				\item \textbf{D - Correct:} Risk exists regardless of detection
			\end{itemize}
		\end{block}
	\end{frame}
	
	\begin{frame}{Q47: Selecting Risk Management Methodology}
		\begin{block}{Question}
			Which is MOST important when selecting an appropriate risk management methodology?
		\end{block}
		\begin{enumerate}[(A)]
			\item Risk culture
			\item Countermeasure analysis
			\item Cost-benefit analysis
			\item Risk transfer strategy
		\end{enumerate}
	\end{frame}
	
	\begin{frame}{Q47: Explanation}
		\begin{block}{Correct Answer: A}
			Without understanding risk culture—how and why an enterprise makes decisions regarding risk—one cannot select a risk management methodology.
		\end{block}
		\begin{block}{Option Analysis}
			\begin{itemize}
				\item \textbf{A - Correct:} Risk culture informs methodology selection
				\item \textbf{B - Incorrect:} Countermeasure analysis targets specific attacks
				\item \textbf{C - Incorrect:} Cost-benefit does not inform methodology selection
				\item \textbf{D - Incorrect:} Risk transfer is not the basis for methodology selection
			\end{itemize}
		\end{block}
	\end{frame}
	
	\begin{frame}{Q48: Determining Risk Appetite Compliance}
		\begin{block}{Question}
			Which BEST determines compliance with the risk appetite of an enterprise?
		\end{block}
		\begin{enumerate}[(A)]
			\item Balance between preventive and detective controls
			\item Inherent risk and acceptable risk level
			\item Residual risk and acceptable risk level
			\item Balance between countermeasures and preventive controls
		\end{enumerate}
	\end{frame}
	
	\begin{frame}{Q48: Explanation}
		\begin{block}{Correct Answer: C}
			Considering residual risk in terms of acceptable risk yields risk that is appropriately balanced after the application of controls. In this context, management can decide to accept risk or apply additional controls based on current standards of acceptable risk.
		\end{block}
		\begin{block}{Option Analysis}
			\begin{itemize}
				\item \textbf{A - Incorrect:} Controls may have been from earlier risk analysis
				\item \textbf{B - Incorrect:} Inherent risk alone is insufficient
				\item \textbf{C - Correct:} Residual risk vs acceptable risk determines compliance
				\item \textbf{D - Incorrect:} Countermeasures don't help evaluate risk appetite
			\end{itemize}
		\end{block}
	\end{frame}
	
	\begin{frame}{Q49: Monitoring External Requirements}
		\begin{block}{Question}
			A key objective when monitoring IS control effectiveness against external requirements is to:
		\end{block}
		\begin{enumerate}[(A)]
			\item Design applicable security controls for external audits
			\item Create IS policy provisions for third parties
			\item Ensure legal obligations have been satisfied
			\item Identify legal obligations that apply
		\end{enumerate}
	\end{frame}
	
	\begin{frame}{Q49: Explanation}
		\begin{block}{Correct Answer: C}
			Legal obligations are one of the principal external requirements that necessitate compliance monitoring.
		\end{block}
		\begin{block}{Option Analysis}
			\begin{itemize}
				\item \textbf{A - Incorrect:} Control design occurs in risk treatment phase
				\item \textbf{B - Incorrect:} Policy creation occurs before monitoring
				\item \textbf{C - Correct:} Legal obligations must be monitored for compliance
				\item \textbf{D - Incorrect:} Identification occurs before risk treatment
			\end{itemize}
		\end{block}
	\end{frame}
	
	\begin{frame}{Q50: Overall Risk Status}
		\begin{block}{Question}
			Which provides an overall risk status of the enterprise?
		\end{block}
		\begin{enumerate}[(A)]
			\item Risk management
			\item Risk analysis
			\item Risk appetite
			\item Risk profile
		\end{enumerate}
	\end{frame}
	
	\begin{frame}{Q50: Explanation}
		\begin{block}{Correct Answer: D}
			The risk profile provides the current overall portfolio of the identified risk to which the enterprise is exposed. Because the profile is updated with evolving and new risk, it provides the enterprise's current risk status.
		\end{block}
		\begin{block}{Option Analysis}
			\begin{itemize}
				\item \textbf{A - Incorrect:} Risk management encompasses coordinated activities
				\item \textbf{B - Incorrect:} Risk analysis is point-in-time
				\item \textbf{C - Incorrect:} Risk appetite is broad acceptance level
				\item \textbf{D - Correct:} Risk profile provides current status
			\end{itemize}
		\end{block}
	\end{frame}
	
	% ================ SECTION 2: DOMAIN 2 - IT RISK ASSESSMENT ================
	
	\section{Domain 2: IT Risk Assessment}
	
	\begin{frame}{Q51: Firewall Deviation Detection}
		\begin{block}{Question}
			Which activity should a risk professional perform to determine whether firewall deployments are deviating from security policy?
		\end{block}
		\begin{enumerate}[(A)]
			\item Review firewall parameter settings
			\item Review firewall intrusion prevention system logs
			\item Review firewall hardening procedures
			\item Analyze firewall log file for recent attacks
		\end{enumerate}
	\end{frame}
	
	\begin{frame}{Q51: Explanation}
		\begin{block}{Correct Answer: A}
			Firewall parameter settings will tie in with the configurations linked to the governing security policy. If the parameter settings differ from what policy states or requires, then there is a deviation.
		\end{block}
		\begin{block}{Option Analysis}
			\begin{itemize}
				\item \textbf{A - Correct:} Parameters show compliance with policy
				\item \textbf{B - Incorrect:} IPS logs show blocked packets, not configuration
				\item \textbf{C - Incorrect:} Hardening procedures show what was expected, not actual
				\item \textbf{D - Incorrect:} Attacks may not be covered in security policy
			\end{itemize}
		\end{block}
	\end{frame}
	
	\begin{frame}{Q52: Continuous Risk Management Benefit}
		\begin{block}{Question}
			Which is MOST significant benefit of implementing continuous risk management process?
		\end{block}
		\begin{enumerate}[(A)]
			\item Complies with external regulators' requirements
			\item Captures changes to the enterprise's risk profile
			\item Ensures stakeholders are linked to identified risk
			\item Ensures timely action is taken to mitigate risk
		\end{enumerate}
	\end{frame}
	
	\begin{frame}{Q52: Explanation}
		\begin{block}{Correct Answer: B}
			A continuous risk management process ensures changes to the enterprise's risk profile are captured.
		\end{block}
		\begin{block}{Option Analysis}
			\begin{itemize}
				\item \textbf{A - Incorrect:} Compliance is not the aim of continuous process
				\item \textbf{B - Correct:} Captures changes to risk profile
				\item \textbf{C - Incorrect:} Linking stakeholders is not the aim
				\item \textbf{D - Incorrect:} Timely mitigation is not the aim
			\end{itemize}
		\end{block}
	\end{frame}
	
	\begin{frame}{Q53: Risk Response When Risk Exceeds Appetite}
		\begin{block}{Question}
			Senior management defined risk appetite as moderate. Business-critical application poses high risk. What is NEXT action?
		\end{block}
		\begin{enumerate}[(A)]
			\item Remove high-risk application and replace with another system
			\item Recommend management increase acceptable risk level
			\item Assess impact of planned controls to the application
			\item Restrict access to application only to trusted users
		\end{enumerate}
	\end{frame}
	
	\begin{frame}{Q53: Explanation}
		\begin{block}{Correct Answer: C}
			The risk practitioner should determine whether the planned implementation of new controls on the system may lower the risk from high to moderate or low before taking any further action.
		\end{block}
		\begin{block}{Option Analysis}
			\begin{itemize}
				\item \textbf{A - Incorrect:} Removal may cause unacceptable downtime
				\item \textbf{B - Incorrect:} This is task for business owners, not risk practitioner
				\item \textbf{C - Correct:} Assess whether controls can reduce risk
				\item \textbf{D - Incorrect:} Restricting access may not mitigate high risk
			\end{itemize}
		\end{block}
	\end{frame}
	
	\begin{frame}{Q54: Risk Profile Development}
		\begin{block}{Question}
			Which BEST assists in development of the risk profile?
		\end{block}
		\begin{enumerate}[(A)]
			\item Presence of preventive and detective controls
			\item Inherent risk and detection risk
			\item Cost-benefit analysis of controls
			\item Likelihood and impact of risk
		\end{enumerate}
	\end{frame}
	
	\begin{frame}{Q54: Explanation}
		\begin{block}{Correct Answer: D}
			Likelihood and impact of risk help in the development of the risk profile.
		\end{block}
		\begin{block}{Option Analysis}
			\begin{itemize}
				\item \textbf{A - Incorrect:} Controls alone don't help
				\item \textbf{B - Incorrect:} Inherent and detection risk alone don't help
				\item \textbf{C - Incorrect:} Cost-benefit helps select controls, not develop profile
				\item \textbf{D - Correct:} Likelihood and impact are fundamental to risk profile
			\end{itemize}
		\end{block}
	\end{frame}
	
	\begin{frame}{Q55: Multi-Location Handbook Concern}
		\begin{block}{Question}
			Enterprise expanded into Europe, Asia and Latin America. Single-version, multi-language employee handbook last updated three years ago. Which is of MOST concern?
		\end{block}
		\begin{enumerate}[(A)]
			\item Handbook may not have been correctly translated
			\item Newer policies may not be included
			\item Expired policies may be included
			\item Handbook may violate local laws and regulations
		\end{enumerate}
	\end{frame}
	
	\begin{frame}{Q55: Explanation}
		\begin{block}{Correct Answer: D}
			Because customs and laws affect an enterprise's ability to operate in a given location, and because both customs and laws vary by state and by country, it is critical for the employee handbook to acknowledge and account for regional domestic and national differences.
		\end{block}
		\begin{block}{Option Analysis}
			\begin{itemize}
				\item \textbf{A - Incorrect:} Translation errors are less damaging than law violations
				\item \textbf{B - Incorrect:} Law compliance is more important
				\item \textbf{C - Incorrect:} Law compliance is more important
				\item \textbf{D - Correct:} Local law compliance is most critical
			\end{itemize}
		\end{block}
	\end{frame}
	
	\begin{frame}{Q56: Annual Loss Expectancy Formula}
		\begin{block}{Question}
			An asset's annual loss expectancy is calculated as the:
		\end{block}
		\begin{enumerate}[(A)]
			\item EF multiplied by ARO
			\item SLE multiplied by EF
			\item SLE multiplied by ARO
			\item Asset value multiplied by SLE
		\end{enumerate}
	\end{frame}
	
	\begin{frame}{Q56: Explanation}
		\begin{block}{Correct Answer: C}
			ALE is calculated by multiplying the SLE by the ARO (the number of times the enterprise expects the loss to occur).
		\end{block}
		\begin{block}{Option Analysis}
			\begin{itemize}
				\item \textbf{A - Incorrect:} Wrong formula
				\item \textbf{B - Incorrect:} Wrong formula
				\item \textbf{C - Correct:} ALE = SLE × ARO
				\item \textbf{D - Incorrect:} Wrong formula
			\end{itemize}
		\end{block}
	\end{frame}
	
	\begin{frame}{Q57: Effective Information Risk Management}
		\begin{block}{Question}
			Which is BEST indicator of an effective information risk management program?
		\end{block}
		\begin{enumerate}[(A)]
			\item Security policy is made widely available
			\item Risk is considered before all decisions
			\item Security procedures are updated annually
			\item Risk assessments occur on an annual basis
		\end{enumerate}
	\end{frame}
	
	\begin{frame}{Q57: Explanation}
		\begin{block}{Correct Answer: B}
			Defining information risk in advance of business decisions best ensures that risk tolerance remains at approved levels.
		\end{block}
		\begin{block}{Option Analysis}
			\begin{itemize}
				\item \textbf{A - Incorrect:} Making policy available facilitates success but isn't critical
				\item \textbf{B - Correct:} Risk-based decisions are key
				\item \textbf{C - Incorrect:} Procedures only need updating if policy changes
				\item \textbf{D - Incorrect:} Annual assessments facilitate but aren't as critical
			\end{itemize}
		\end{block}
	\end{frame}
	
	\begin{frame}{Q58: Risk Management Strategic Plans}
		\begin{block}{Question}
			Risk management strategic plans are MOST effective when developed for:
		\end{block}
		\begin{enumerate}[(A)]
			\item The enterprise as a whole
			\item Each individual system based on technology
			\item Every location based on geographic threats
			\item End-to-end business processes
		\end{enumerate}
	\end{frame}
	
	\begin{frame}{Q58: Explanation}
		\begin{block}{Correct Answer: A}
			Risk management strategic plans are most effective when they are created and followed by the entire enterprise.
		\end{block}
		\begin{block}{Option Analysis}
			\begin{itemize}
				\item \textbf{A - Correct:} Enterprise-wide plans are most effective
				\item \textbf{B - Incorrect:} Creating per technology creates complexity
				\item \textbf{C - Incorrect:} Geographic plans don't account for other threats
				\item \textbf{D - Incorrect:} Per-process plans may overlap and conflict
			\end{itemize}
		\end{block}
	\end{frame}
	
	\begin{frame}{Q59: Risk Appetite Consideration}
		\begin{block}{Question}
			Which is MOST important when considering the risk appetite of an enterprise?
		\end{block}
		\begin{enumerate}[(A)]
			\item Capacity of enterprise to absorb loss
			\item Definition of responsibilities for risk management
			\item Line of business and typical risk of industry
			\item Culture and predisposition toward risk taking
		\end{enumerate}
	\end{frame}
	
	\begin{frame}{Q59: Explanation}
		\begin{block}{Correct Answer: D}
			When considering risk appetite, two major factors are relevant: the management culture and the predisposition toward risk taking.
		\end{block}
		\begin{block}{Option Analysis}
			\begin{itemize}
				\item \textbf{A - Incorrect:} Capacity is a mitigation factor, not appetite
				\item \textbf{B - Incorrect:} Responsibilities say nothing about appetite
				\item \textbf{C - Incorrect:} Industry risk doesn't determine single enterprise appetite
				\item \textbf{D - Correct:} Culture and predisposition are primary
			\end{itemize}
		\end{block}
	\end{frame}
	
	\begin{frame}{Q60: Enterprise-Wide Risk Framework Purpose}
		\begin{block}{Question}
			PRIMARY purpose of adopting an enterprise-wide risk management framework is to:
		\end{block}
		\begin{enumerate}[(A)]
			\item Allow flexibility to adjust risk response strategy
			\item Centralize responsibility for maintenance
			\item Enable a consistent approach to risk response
			\item Avoid higher costs for risk reduction
		\end{enumerate}
	\end{frame}
	
	\begin{frame}{Q60: Explanation}
		\begin{block}{Correct Answer: C}
			Enabling consistent risk response is a key objective of the risk management framework.
		\end{block}
		\begin{block}{Option Analysis}
			\begin{itemize}
				\item \textbf{A - Incorrect:} Framework enables consistency while allowing flexibility
				\item \textbf{B - Incorrect:} Risk management is everyone's responsibility
				\item \textbf{C - Correct:} Consistent approach is key
				\item \textbf{D - Incorrect:} Cost avoidance is good practice but not primary purpose
			\end{itemize}
		\end{block}
	\end{frame}
	
	\begin{frame}{Q61: Risk Tolerance Example}
		\begin{block}{Question}
			Review finds projects exceed timelines/budget by nearly 10%. Management advises 15% deviation acceptable. This is:
		\end{block}
		\begin{enumerate}[(A)]
			\item Risk avoidance
			\item Risk tolerance
			\item Risk acceptance
			\item Risk mitigation
		\end{enumerate}
	\end{frame}
	
	\begin{frame}{Q61: Explanation}
		\begin{block}{Correct Answer: B}
			Risk tolerance is the permissible deviation from declared risk appetite levels.
		\end{block}
		\begin{block}{Option Analysis}
			\begin{itemize}
				\item \textbf{A - Incorrect:} Avoidance terminates activity
				\item \textbf{B - Correct:} Tolerance is acceptable deviation
				\item \textbf{C - Incorrect:} Acceptance means no action taken
				\item \textbf{D - Incorrect:} Mitigation uses countermeasures
			\end{itemize}
		\end{block}
	\end{frame}
	
	\begin{frame}{Q62: Computing Annual Loss Exposure}
		\begin{block}{Question}
			Which is MOST useful when computing annual loss exposure?
		\end{block}
		\begin{enumerate}[(A)]
			\item Cost of existing controls
			\item Number of vulnerabilities
			\item Net present value of asset
			\item Business value of asset
		\end{enumerate}
	\end{frame}
	
	\begin{frame}{Q62: Explanation}
		\begin{block}{Correct Answer: D}
			Annual loss exposure is a function of the value of the information asset and the impact if a given potential risk should materialize.
		\end{block}
		\begin{block}{Option Analysis}
			\begin{itemize}
				\item \textbf{A - Incorrect:} Cost of controls not used in calculation
				\item \textbf{B - Incorrect:} Number of vulnerabilities doesn't determine exposure
				\item \textbf{C - Incorrect:} NPV may not reflect true risk
				\item \textbf{D - Correct:} Business value is fundamental
			\end{itemize}
		\end{block}
	\end{frame}
	
	\begin{frame}{Q63: GRC Tool Major Risk}
		\begin{block}{Question}
			Which is MAJOR risk associated with use of governance, risk and compliance tools?
		\end{block}
		\begin{enumerate}[(A)]
			\item Misinterpretation of dashboard's output
			\item Incomplete or inaccurate coverage
			\item Obsolescence of content
			\item Complex integration of diverse requirements
		\end{enumerate}
	\end{frame}
	
	\begin{frame}{Q63: Explanation}
		\begin{block}{Correct Answer: C}
			A GRC application has to be updated regularly with current regulations, policies, etc. Obsolete content renders a GRC outdated.
		\end{block}
		\begin{block}{Option Analysis}
			\begin{itemize}
				\item \textbf{A - Incorrect:} Misinterpretation is easily corrected by training
				\item \textbf{B - Incorrect:} Incomplete coverage can be addressed
				\item \textbf{C - Correct:} Obsolescence is the greatest risk
				\item \textbf{D - Incorrect:} Tools are designed to integrate requirements
			\end{itemize}
		\end{block}
	\end{frame}
	
	\begin{frame}{Q64: Need to Review Risk Practices}
		\begin{block}{Question}
			Which signifies the need to review an enterprise's risk practices?
		\end{block}
		\begin{enumerate}[(A)]
			\item Business owners regularly challenge risk assessment findings
			\item Manufacturing assigns its own internal risk management roles
			\item Finance department finds exceptions during yearly risk review
			\item Sales department procedures last reviewed 11 months ago
		\end{enumerate}
	\end{frame}
	
	\begin{frame}{Q64: Explanation}
		\begin{block}{Correct Answer: A}
			An enterprise's risk management practices must be clearly understood and supported by business stakeholders. Business owners who challenge the risk assessment findings either do not support the findings or do not understand them clearly.
		\end{block}
		\begin{block}{Option Analysis}
			\begin{itemize}
				\item \textbf{A - Correct:} Challenge indicates need for review
				\item \textbf{B - Incorrect:} Assigning roles is what departments should do
				\item \textbf{C - Incorrect:} Finding exceptions is normal
				\item \textbf{D - Incorrect:} Annual review is sufficient
			\end{itemize}
		\end{block}
	\end{frame}
	
	\begin{frame}{Q65: Relevance Risk}
		\begin{block}{Question}
			Correct information was not received by necessary recipients in suitable time. This is:
		\end{block}
		\begin{enumerate}[(A)]
			\item Integrity risk
			\item Availability risk
			\item Access risk
			\item Relevance risk
		\end{enumerate}
	\end{frame}
	
	\begin{frame}{Q65: Explanation}
		\begin{block}{Correct Answer: D}
			Relevance risk is a composite form of business risk, requiring both integrity and availability to be addressed for it to be reasonably controlled. Transmitting information to the necessary recipients in a timely manner also creates tension with access (security) risk.
		\end{block}
		\begin{block}{Option Analysis}
			\begin{itemize}
				\item \textbf{A - Incorrect:} Integrity is incomplete/inaccurate data
				\item \textbf{B - Incorrect:} Availability is only one possible reason
				\item \textbf{C - Incorrect:} Access risk is unauthorized disclosure
				\item \textbf{D - Correct:} Relevance risk covers timeliness and correctness
			\end{itemize}
		\end{block}
	\end{frame}
	
	\begin{frame}{Q66: Best Risk Assessment Output}
		\begin{block}{Question}
			Which risk assessment output is MOST suitable to help justify an information security program?
		\end{block}
		\begin{enumerate}[(A)]
			\item Inventory of risk
			\item Documented threats
			\item Evaluation of consequences
			\item List of appropriate controls
		\end{enumerate}
	\end{frame}
	
	\begin{frame}{Q66: Explanation}
		\begin{block}{Correct Answer: D}
			A list of information security controls corresponding to risk scenarios identified during risk assessment is one of the primary deliverables of the risk assessment exercise. The list demonstrates due consideration of risk and applicable controls to address the risk.
		\end{block}
		\begin{block}{Option Analysis}
			\begin{itemize}
				\item \textbf{A - Incorrect:} Doesn't cover how risk will be addressed
				\item \textbf{B - Incorrect:} Doesn't explain how risk will be reduced
				\item \textbf{C - Incorrect:} Doesn't represent controls
				\item \textbf{D - Correct:} List of controls demonstrates risk mitigation
			\end{itemize}
		\end{block}
	\end{frame}
	
	\begin{frame}{Q67: Ultimate Accountability for Risk}
		\begin{block}{Question}
			Accountability for risk ultimately belongs to the:
		\end{block}
		\begin{enumerate}[(A)]
			\item Chief risk officer
			\item Compliance officer
			\item Chief financial officer
			\item Board of directors
		\end{enumerate}
	\end{frame}
	
	\begin{frame}{Q67: Explanation}
		\begin{block}{Correct Answer: D}
			The board of directors of an enterprise has ultimate accountability to shareholders, customers, employees and the general public.
		\end{block}
		\begin{block}{Option Analysis}
			\begin{itemize}
				\item \textbf{A - Incorrect:} CRO has responsibilities but not ultimate accountability
				\item \textbf{B - Incorrect:} Compliance officer has responsibility but not ultimate
				\item \textbf{C - Incorrect:} CFO has major responsibilities but not ultimate
				\item \textbf{D - Correct:} Board has ultimate accountability
			\end{itemize}
		\end{block}
	\end{frame}
	
	\begin{frame}{Q68: Outdated IT Standards}
		\begin{block}{Question}
			Risk assessment reveals many corporate IT standards have not been updated. BEST course of action is:
		\end{block}
		\begin{enumerate}[(A)]
			\item Review standards against current requirements and determine adequacy
			\item Determine standards should be updated annually
			\item Report that IT standards are adequate
			\item Review IT policy to see how frequently standards should be updated
		\end{enumerate}
	\end{frame}
	
	\begin{frame}{Q68: Explanation}
		\begin{block}{Correct Answer: A}
			The risk practitioner should verify that the standards are still adequate. Standards that are lacking should be updated.
		\end{block}
		\begin{block}{Option Analysis}
			\begin{itemize}
				\item \textbf{A - Correct:} Review and determine adequacy
				\item \textbf{B - Incorrect:} Standards may or may not need annual updates
				\item \textbf{C - Incorrect:} Cannot report adequacy without review
				\item \textbf{D - Incorrect:} Policy review won't determine adequacy
			\end{itemize}
		\end{block}
	\end{frame}
	
	\begin{frame}{Q69: Accountability for Third-Party App Risk}
		\begin{block}{Question}
			Marketing procures third-party application. Application poses risk to data privacy regulations in certain regions. Who is accountable if implemented globally?
		\end{block}
		\begin{enumerate}[(A)]
			\item IT department
			\item Data privacy officer
			\item Chief risk officer
			\item Marketing department
		\end{enumerate}
	\end{frame}
	
	\begin{frame}{Q69: Explanation}
		\begin{block}{Correct Answer: D}
			The marketing department is the business owner of the application and, therefore, must be accountable. According to ISACA's COBIT framework, accountability rests with those who own the required resources and have the authority to approve the execution and accept the outcome.
		\end{block}
		\begin{block}{Option Analysis}
			\begin{itemize}
				\item \textbf{A - Incorrect:} IT is responsible, not accountable
				\item \textbf{B - Incorrect:} Privacy officer is advisory
				\item \textbf{C - Incorrect:} CRO is advisory
				\item \textbf{D - Correct:} Marketing owns the application
			\end{itemize}
		\end{block}
	\end{frame}
	
	\begin{frame}{Q70: Risk Treatment Decisions Ownership}
		\begin{block}{Question}
			Who owns the risk treatment decisions?
		\end{block}
		\begin{enumerate}[(A)]
			\item Control owner
			\item Chief information officer
			\item IT security manager
			\item Senior management
		\end{enumerate}
	\end{frame}
	
	\begin{frame}{Q70: Explanation}
		\begin{block}{Correct Answer: D}
			Senior management ultimately owns the risk and would make risk treatment decisions.
		\end{block}
		\begin{block}{Option Analysis}
			\begin{itemize}
				\item \textbf{A - Incorrect:} Control owner would not make risk treatment decisions
				\item \textbf{B - Incorrect:} CIO is part of senior management but not sole decision-maker
				\item \textbf{C - Incorrect:} IT security manager implements, doesn't decide
				\item \textbf{D - Correct:} Senior management owns risk
			\end{itemize}
		\end{block}
	\end{frame}
	
	\begin{frame}{Q71: Responsibility for Third-Party App Risk}
		\begin{block}{Question}
			Marketing procures third-party application with data privacy risk. Which role will be responsible for the risk?
		\end{block}
		\begin{enumerate}[(A)]
			\item Marketing department
			\item IT department
			\item Data privacy officer
			\item Chief risk officer
		\end{enumerate}
	\end{frame}
	
	\begin{frame}{Q71: Explanation}
		\begin{block}{Correct Answer: B}
			The IT department is responsible for the risk posed by this application. The IT department has a policy in place that states that no tool or application can be implemented without a risk assessment and mitigation of all risk to an acceptable level.
		\end{block}
		\begin{block}{Option Analysis}
			\begin{itemize}
				\item \textbf{A - Incorrect:} Marketing is accountable, not responsible
				\item \textbf{B - Correct:} IT is responsible for risk assessment
				\item \textbf{C - Incorrect:} Privacy officer is advisory
				\item \textbf{D - Incorrect:} CRO is advisory
			\end{itemize}
		\end{block}
	\end{frame}
	
	\begin{frame}{Q72: Justifying Encryption}
		\begin{block}{Question}
			Which provides GREATEST support to risk practitioner recommending encryption as risk mitigation?
		\end{block}
		\begin{enumerate}[(A)]
			\item Benchmarking with peers
			\item Evaluating public reports on encryption
			\item Developing a business case
			\item Scanning unencrypted systems for vulnerabilities
		\end{enumerate}
	\end{frame}
	
	\begin{frame}{Q72: Explanation}
		\begin{block}{Correct Answer: C}
			A business case documents how encryption addresses risk and illustrates that risk losses can be reduced by recommended controls.
		\end{block}
		\begin{block}{Option Analysis}
			\begin{itemize}
				\item \textbf{A - Incorrect:} Peers have different risk environments
				\item \textbf{B - Incorrect:} Need to analyze in enterprise context
				\item \textbf{C - Correct:} Business case justifies the control
				\item \textbf{D - Incorrect:} Vulnerability assessment without business justification doesn't help
			\end{itemize}
		\end{block}
	\end{frame}
	
	\begin{frame}{Q73: Additional Stakeholder for Risk Reporting}
		\begin{block}{Question}
			Marketing procures third-party application with data privacy risk. Which additional stakeholder should be included in reporting?
		\end{block}
		\begin{enumerate}[(A)]
			\item Chief privacy officer
			\item Third-party vendor
			\item Business continuity manager
			\item Audit manager
		\end{enumerate}
	\end{frame}
	
	\begin{frame}{Q73: Explanation}
		\begin{block}{Correct Answer: A}
			The chief privacy officer will offer support in defining controls that will be implemented to address and mitigate the data transfer and data privacy requirements. If the application must be used by the business, the chief privacy officer ultimately becomes a critical stakeholder.
		\end{block}
		\begin{block}{Option Analysis}
			\begin{itemize}
				\item \textbf{A - Correct:} Privacy officer is critical stakeholder
				\item \textbf{B - Incorrect:} No need to inform vendor
				\item \textbf{C - Incorrect:} No need to inform BC manager
				\item \textbf{D - Incorrect:} No need to inform audit manager
			\end{itemize}
		\end{block}
	\end{frame}
	
	\begin{frame}{Q74: Cloud Email Retention}
		\begin{block}{Question}
			Enterprise learned cloud email provider was never contracted to back up messages despite retention policy requiring three-year saving. Which would have BEST safeguarded?
		\end{block}
		\begin{enumerate}[(A)]
			\item Providing contractor with retention policy upfront
			\item Validating company policies to provider's contract
			\item Providing contractor with email retention policy upfront
			\item Backing up data on internal network nightly
		\end{enumerate}
	\end{frame}
	
	\begin{frame}{Q74: Explanation}
		\begin{block}{Correct Answer: B}
			The initial review of third-party services should confirm that vendors are contractually required to enforce all internal policies—including the policy on record retention if the enterprise's record retention policy specifically covers data that will be managed by a third party.
		\end{block}
		\begin{block}{Option Analysis}
			\begin{itemize}
				\item \textbf{A - Incorrect:} Providing policy doesn't legally obligate
				\item \textbf{B - Correct:} Contractual validation is essential
				\item \textbf{C - Incorrect:} Providing policy doesn't legally obligate
				\item \textbf{D - Incorrect:} Internal backup doesn't solve policy variance issue
			\end{itemize}
		\end{block}
	\end{frame}
	
	\begin{frame}{Q75: Personal Device Policy}
		\begin{block}{Question}
			If enterprise does not have formal policy regarding personal devices in workplace, what should risk practitioner recommend?
		\end{block}
		\begin{enumerate}[(A)]
			\item Introduce remote wipe functionality
			\item Implement exception process based on approvals
			\item Update incident response procedures
			\item Create and maintain inventory of personal devices
		\end{enumerate}
	\end{frame}
	
	\begin{frame}{Q75: Explanation}
		\begin{block}{Correct Answer: B}
			If employees are allowed to attach personal devices to a company network but there is no formal policy in place, employees should be required to gain approval from management, and this approval should be documented and reviewed in a security exception form based on the exception process.
		\end{block}
		\begin{block}{Option Analysis}
			\begin{itemize}
				\item \textbf{A - Incorrect:} Implementation would be part of formal policy
				\item \textbf{B - Correct:} Exception process is needed without formal policy
				\item \textbf{C - Incorrect:} Procedures can't be updated without policy
				\item \textbf{D - Incorrect:} Inventorying doesn't reduce risk
			\end{itemize}
		\end{block}
	\end{frame}
	
	\begin{frame}{Q76: Risk Management Application}
		\begin{block}{Question}
			To be effective, risk management should be applied to:
		\end{block}
		\begin{enumerate}[(A)]
			\item Elements identified by risk assessment
			\item Any area that exceeds acceptable risk levels
			\item All enterprise activities
			\item Only areas that have potential impact
		\end{enumerate}
	\end{frame}
	
	\begin{frame}{Q76: Explanation}
		\begin{block}{Correct Answer: C}
			Although not all enterprise activities entail unacceptable risk, the practice of risk management is ideally applied to all enterprise activities.
		\end{block}
		\begin{block}{Option Analysis}
			\begin{itemize}
				\item \textbf{A - Incorrect:} All activities subject to risk management oversight
				\item \textbf{B - Incorrect:} Must be holistic, not limited to unacceptable risk
				\item \textbf{C - Correct:} Applied to all enterprise activities
				\item \textbf{D - Incorrect:} Determining impact is part of risk management
			\end{itemize}
		\end{block}
	\end{frame}
	
	\begin{frame}{Q77: Corporate Security Policy Basis}
		\begin{block}{Question}
			Development of corporate information security policy should PRIMARILY be based on:
		\end{block}
		\begin{enumerate}[(A)]
			\item Vulnerabilities
			\item Threats
			\item Assets
			\item Impacts
		\end{enumerate}
	\end{frame}
	
	\begin{frame}{Q77: Explanation}
		\begin{block}{Correct Answer: C}
			The corporate information security policy is based on management's commitment to protect the assets of the enterprise (and relevant information of its business partners) from threats, risk and exposures that could occur.
		\end{block}
		\begin{block}{Option Analysis}
			\begin{itemize}
				\item \textbf{A - Incorrect:} Vulnerabilities don't pose risk without threats
				\item \textbf{B - Incorrect:} Policy addresses protection of assets from threats
				\item \textbf{C - Correct:} Assets are the basis for protection
				\item \textbf{D - Incorrect:} Impact doesn't drive policy development
			\end{itemize}
		\end{block}
	\end{frame}
	
	\begin{frame}{Q78: Enabling Risk-Aware Decisions}
		\begin{block}{Question}
			Which MOST enables risk-aware business decisions?
		\end{block}
		\begin{enumerate}[(A)]
			\item Robust information security policies
			\item Exchange of accurate and timely information
			\item Skilled risk management personnel
			\item Robust process controls
		\end{enumerate}
	\end{frame}
	
	\begin{frame}{Q78: Explanation}
		\begin{block}{Correct Answer: B}
			Robust exchange of information enables management to optimize risk-related decisions. Accuracy and timeliness of information are critical success factors.
		\end{block}
		\begin{block}{Option Analysis}
			\begin{itemize}
				\item \textbf{A - Incorrect:} Policies focus on protecting business
				\item \textbf{B - Correct:} Accurate/timely information enables decisions
				\item \textbf{C - Incorrect:} Personnel enable decisions but bidirectional flow is needed
				\item \textbf{D - Incorrect:} Process controls exist for known threats
			\end{itemize}
		\end{block}
	\end{frame}
	
	\begin{frame}{Q79: Initial Risk Management Program Requirements}
		\begin{block}{Question}
			Which MUST be met during initial stages of developing a risk management program?
		\end{block}
		\begin{enumerate}[(A)]
			\item Management establishes ownership of identified risk
			\item Information security policies and standards are established
			\item Management committee exists to provide program oversight
			\item Context and purpose of program are defined
		\end{enumerate}
	\end{frame}
	
	\begin{frame}{Q79: Explanation}
		\begin{block}{Correct Answer: D}
			Initial requirements to determine the enterprise's purpose for creating an information security risk management program include determining the desired outcomes and defining objectives.
		\end{block}
		\begin{block}{Option Analysis}
			\begin{itemize}
				\item \textbf{A - Incorrect:} Ownership occurs later
				\item \textbf{B - Incorrect:} Policies are based on decisions in planning phase
				\item \textbf{C - Incorrect:} Oversight is established later
				\item \textbf{D - Correct:} Context and purpose come first
			\end{itemize}
		\end{block}
	\end{frame}
	
	\begin{frame}{Q80: BYOD Policy Effectiveness}
		\begin{block}{Question}
			Which approach MOST enables BYOD policy to be effective?
		\end{block}
		\begin{enumerate}[(A)]
			\item Obtaining signed acceptance from users
			\item Educating users on acceptable and unacceptable practices
			\item Requiring users to read policy and updates
			\item Clearly stating disciplinary action for noncompliance
		\end{enumerate}
	\end{frame}
	
	\begin{frame}{Q80: Explanation}
		\begin{block}{Correct Answer: B}
			A BYOD policy that is approved and communicated will not be effective without proper training of all users.
		\end{block}
		\begin{block}{Option Analysis}
			\begin{itemize}
				\item \textbf{A - Incorrect:} Sign-off doesn't imply understanding
				\item \textbf{B - Correct:} Training is essential for effectiveness
				\item \textbf{C - Incorrect:} Reading alone won't ensure understanding
				\item \textbf{D - Incorrect:} Disciplinary action is important but not best option
			\end{itemize}
		\end{block}
	\end{frame}
	
	\begin{frame}{Q81: New Regulation First Step}
		\begin{block}{Question}
			New regulation for safeguarding information comes to IT manager's attention. Manager should FIRST:
		\end{block}
		\begin{enumerate}[(A)]
			\item Meet with stakeholders to decide how to comply
			\item Analyze key risk in compliance process
			\item Update existing security/privacy policy
			\item Assess whether existing controls meet the regulation
		\end{enumerate}
	\end{frame}
	
	\begin{frame}{Q81: Explanation}
		\begin{block}{Correct Answer: D}
			The first step is to understand the impact and requirements of the new regulation, which includes assessing how the enterprise will comply with the regulation and to what extent existing controls support compliance. The risk practitioner should then assess any existing gaps.
		\end{block}
		\begin{block}{Option Analysis}
			\begin{itemize}
				\item \textbf{A - Incorrect:} Meeting is subsequent
				\item \textbf{B - Incorrect:} Analysis is subsequent
				\item \textbf{C - Incorrect:} Updating policy is subsequent
				\item \textbf{D - Correct:} Assessment comes first
			\end{itemize}
		\end{block}
	\end{frame}
	
	\begin{frame}{Q82: Compliance Risk and Corrective Actions}
		\begin{block}{Question}
			Which document BEST identifies compliance risk and corrective actions in progress?
		\end{block}
		\begin{enumerate}[(A)]
			\item Internal audit report
			\item Risk register
			\item External audit report
			\item Risk assessment report
		\end{enumerate}
	\end{frame}
	
	\begin{frame}{Q82: Explanation}
		\begin{block}{Correct Answer: B}
			A risk register provides a report of all current identified risk within an enterprise (including compliance risk) and tracks the status of corrective actions or exceptions.
		\end{block}
		\begin{block}{Option Analysis}
			\begin{itemize}
				\item \textbf{A - Incorrect:} Audit reports track findings but not necessarily compliance risk
				\item \textbf{B - Correct:} Risk register tracks all risk and corrective actions
				\item \textbf{C - Incorrect:} External audits focus on one requirement at a time
				\item \textbf{D - Incorrect:} Risk assessment may include compliance risk but doesn't track corrective actions
			\end{itemize}
		\end{block}
	\end{frame}
	
	\begin{frame}{Q83: Impact on Security Governance Models}
		\begin{block}{Question}
			Which will have MOST significant impact on standard information security governance models?
		\end{block}
		\begin{enumerate}[(A)]
			\item Number of employees and consulting staff
			\item Cultural differences between physical locations
			\item Complexity of organizational structure
			\item Uncertainty in legislative requirements
		\end{enumerate}
	\end{frame}
	
	\begin{frame}{Q83: Explanation}
		\begin{block}{Correct Answer: C}
			Information security governance models are highly dependent on the complexity of the organizational structure. Elements that affect organizational structure include multiple business units, dispersion of multiple functions, multiple leadership hierarchies and multiple lines of communication.
		\end{block}
		\begin{block}{Option Analysis}
			\begin{itemize}
				\item \textbf{A - Incorrect:} Well-defined processes provide governance
				\item \textbf{B - Incorrect:} Well-defined processes provide governance
				\item \textbf{C - Correct:} Organizational structure is primary
				\item \textbf{D - Incorrect:} Good governance models handle legislative requirements
			\end{itemize}
		\end{block}
	\end{frame}
	
	\begin{frame}{Q84: Highest Tolerance for Public Cloud}
		\begin{block}{Question}
			Senior management will MOST likely have highest tolerance for moving which to public cloud?
		\end{block}
		\begin{enumerate}[(A)]
			\item Credit-card processing
			\item Research and development
			\item Legacy financial system
			\item Corporate email system
		\end{enumerate}
	\end{frame}
	
	\begin{frame}{Q84: Explanation}
		\begin{block}{Correct Answer: D}
			Consideration for moving processes and information to the cloud should include the criticality, complexity and classification of the data. Of the options offered, the corporate email system has the least competitive distinction, complexity, and sensitive/highly classified information.
		\end{block}
		\begin{block}{Option Analysis}
			\begin{itemize}
				\item \textbf{A - Incorrect:} Credit-card processing is high security
				\item \textbf{B - Incorrect:} R\&D contains confidential proprietary information
				\item \textbf{C - Incorrect:} Legacy financial system is sensitive and complex
				\item \textbf{D - Correct:} Email has least competitive distinction
			\end{itemize}
		\end{block}
	\end{frame}
	
	\begin{frame}{Q85: Risk Profile Consideration}
		\begin{block}{Question}
			Which is MOST important to consider in relation to a risk profile?
		\end{block}
		\begin{enumerate}[(A)]
			\item Summary of regional loss events
			\item Aggregated risk to the enterprise
			\item Description of critical risk
			\item Analysis of historical loss events
		\end{enumerate}
	\end{frame}
	
	\begin{frame}{Q85: Explanation}
		\begin{block}{Correct Answer: B}
			The risk profile is based on the aggregated risk to the enterprise, including historical risk, critical risk and emerging risk.
		\end{block}
		\begin{block}{Option Analysis}
			\begin{itemize}
				\item \textbf{A - Incorrect:} Regional events are one component
				\item \textbf{B - Correct:} Aggregated risk is the basis
				\item \textbf{C - Incorrect:} All risk is considered, not just critical
				\item \textbf{D - Incorrect:} Historical analysis assists but doesn't substantively inform
			\end{itemize}
		\end{block}
	\end{frame}
	
	\begin{frame}{Q86: Collaborative Research Cloud Model}
		\begin{block}{Question}
			Which cloud computing deployment model is MOST appropriate for collaborative research program between universities?
		\end{block}
		\begin{enumerate}[(A)]
			\item Private cloud
			\item Community cloud
			\item Public cloud
			\item Hybrid cloud
		\end{enumerate}
	\end{frame}
	
	\begin{frame}{Q86: Explanation}
		\begin{block}{Correct Answer: B}
			A community cloud deployment model is appropriate because it is shared between several enterprises and supports a specific community that has a common mission or interest.
		\end{block}
		\begin{block}{Option Analysis}
			\begin{itemize}
				\item \textbf{A - Incorrect:} Private cloud is for single enterprise
				\item \textbf{B - Correct:} Community cloud is for shared mission
				\item \textbf{C - Incorrect:} Public cloud is for general public
				\item \textbf{D - Incorrect:} Hybrid is combination
			\end{itemize}
		\end{block}
	\end{frame}
	
	\begin{frame}{Q87: Control Environment Review}
		\begin{block}{Question}
			Which line of defense is MOST likely involved when conducting a review of adequacy of control environment?
		\end{block}
		\begin{enumerate}[(A)]
			\item First line of defense
			\item Second line of defense
			\item Third line of defense
			\item Third-party line of defense
		\end{enumerate}
	\end{frame}
	
	\begin{frame}{Q87: Explanation}
		\begin{block}{Correct Answer: A}
			A control environment review is a first line of defense activity, where control owners are responsible for adequacy of the control environment.
		\end{block}
		\begin{block}{Option Analysis}
			\begin{itemize}
				\item \textbf{A - Correct:} First line reviews control environment
				\item \textbf{B - Incorrect:} Second line provides risk-based consultation
				\item \textbf{C - Incorrect:} Third line provides independent attestation
				\item \textbf{D - Incorrect:} Not part of 3LoD model
			\end{itemize}
		\end{block}
	\end{frame}
	
	\begin{frame}{Q88: Board Concern About Risk Management}
		\begin{block}{Question}
			Which issue would be of MOST concern to the board when assessing risk management capability?
		\end{block}
		\begin{enumerate}[(A)]
			\item Third line of defense is acting independently
			\item Second line is checking and challenging first line
			\item Internal audit is preparing risk management strategy
			\item Each line of defense is planning independently
		\end{enumerate}
	\end{frame}
	
	\begin{frame}{Q88: Explanation}
		\begin{block}{Correct Answer: D}
			Joint planning across the three lines of defense is key to achieving an effective risk management capability. If each line is planning independently, it is an indicator that the lines of defense are not working as expected.
		\end{block}
		\begin{block}{Option Analysis}
			\begin{itemize}
				\item \textbf{A - Incorrect:} Independence is expected
				\item \textbf{B - Incorrect:} Checking is a role of second line
				\item \textbf{C - Incorrect:} Preparing strategy is expected audit activity
				\item \textbf{D - Correct:} Independent planning indicates dysfunction
			\end{itemize}
		\end{block}
	\end{frame}
	
	\begin{frame}{Q89: Code of Ethics Adherence}
		\begin{block}{Question}
			Which MOST effectively supports adherence to enterprise's code of ethics?
		\end{block}
		\begin{enumerate}[(A)]
			\item Ensuring periodic training, evaluation and attestation
			\item Performing background checks at hire
			\item Providing clear policies and standards
			\item Providing continuous awareness throughout the year
		\end{enumerate}
	\end{frame}
	
	\begin{frame}{Q89: Explanation}
		\begin{block}{Correct Answer: A}
			Attestation to comply with the enterprise's code of conduct during ethics training ensures that employees have a clear idea of what is expected of them in terms of aligning with the enterprise's ethics code. Training will also provide evaluation to measure employees' understanding.
		\end{block}
		\begin{block}{Option Analysis}
			\begin{itemize}
				\item \textbf{A - Correct:} Training and attestation ensure understanding
				\item \textbf{B - Incorrect:} Background checks don't check current behavior
				\item \textbf{C - Incorrect:} Policies without attestation don't ensure understanding
				\item \textbf{D - Incorrect:} Continuous training may not be feasible
			\end{itemize}
		\end{block}
	\end{frame}
	
	\begin{frame}{Q90: Explaining Zero-Day Exploit}
		\begin{block}{Question}
			Who is responsible for explaining ramifications of new zero-day exploit to senior management?
		\end{block}
		\begin{enumerate}[(A)]
			\item Chief operating officer
			\item Chief risk officer
			\item Chief information security officer
			\item Chief information officer
		\end{enumerate}
	\end{frame}
	
	\begin{frame}{Q90: Explanation}
		\begin{block}{Correct Answer: B}
			The chief risk officer is the most senior official accountable for all aspects of risk management across the enterprise, including explaining risk to senior management.
		\end{block}
		\begin{block}{Option Analysis}
			\begin{itemize}
				\item \textbf{A - Incorrect:} COO operates the enterprise
				\item \textbf{B - Correct:} CRO explains risk to management
				\item \textbf{C - Incorrect:} CISO leads security program but doesn't explain to senior management
				\item \textbf{D - Incorrect:} CIO aligns IT and business
			\end{itemize}
		\end{block}
	\end{frame}
	
	\begin{frame}{Q91: Information Classification Responsibility}
		\begin{block}{Question}
			Application of information classification is the responsibility of the:
		\end{block}
		\begin{enumerate}[(A)]
			\item Information security officer
			\item Information owner
			\item Information systems auditor
			\item Information custodian
		\end{enumerate}
	\end{frame}
	
	\begin{frame}{Q91: Explanation}
		\begin{block}{Correct Answer: B}
			The information owner determines classification based on the criticality and sensitivity of information.
		\end{block}
		\begin{block}{Option Analysis}
			\begin{itemize}
				\item \textbf{A - Incorrect:} Security officer has functional responsibility
				\item \textbf{B - Correct:} Information owner classifies data
				\item \textbf{C - Incorrect:} Auditor examines security
				\item \textbf{D - Incorrect:} Custodian preserves CIA
			\end{itemize}
		\end{block}
	\end{frame}
	
	\begin{frame}{Q92: 3LoD Model Benefit}
		\begin{block}{Question}
			MOST significant benefit of using three lines of defense model in risk management framework is that it:
		\end{block}
		\begin{enumerate}[(A)]
			\item Ensures ongoing success of risk management initiative
			\item Enhances communication between stakeholders
			\item Clarifies essential roles of key stakeholders
			\item Helps risk owners select appropriate control owners
		\end{enumerate}
	\end{frame}
	
	\begin{frame}{Q92: Explanation}
		\begin{block}{Correct Answer: C}
			The three lines of defense model is designed to assign roles related to risk, controls and security while better defining and understanding the relationships between them.
		\end{block}
		\begin{block}{Option Analysis}
			\begin{itemize}
				\item \textbf{A - Incorrect:} Model supports allocation, doesn't ensure success
				\item \textbf{B - Incorrect:} Not the most significant benefit
				\item \textbf{C - Correct:} Clarifies roles is the main benefit
				\item \textbf{D - Incorrect:} Not the most significant benefit
			\end{itemize}
		\end{block}
	\end{frame}
	
	\begin{frame}{Q93: First Line of Defense Role}
		\begin{block}{Question}
			Which risk management role is part of first line of defense?
		\end{block}
		\begin{enumerate}[(A)]
			\item Chief risk officer
			\item Risk steering committee
			\item Risk owner
			\item Board of directors
		\end{enumerate}
	\end{frame}
	
	\begin{frame}{Q93: Explanation}
		\begin{block}{Correct Answer: C}
			The first line of defense is operational management, and risk owners are part of operational management.
		\end{block}
		\begin{block}{Option Analysis}
			\begin{itemize}
				\item \textbf{A - Incorrect:} CRO is second line
				\item \textbf{B - Incorrect:} Steering committee is second line
				\item \textbf{C - Correct:} Risk owner is operational management
				\item \textbf{D - Incorrect:} Board provides direction, not part of 3LoD
			\end{itemize}
		\end{block}
	\end{frame}
	
	\begin{frame}{Q94: Ultimate Objective of Risk Governance}
		\begin{block}{Question}
			What is the ULTIMATE objective of risk governance?
		\end{block}
		\begin{enumerate}[(A)]
			\item Benefits realization
			\item Risk optimization
			\item Resource optimization
			\item Value creation
		\end{enumerate}
	\end{frame}
	
	\begin{frame}{Q94: Explanation}
		\begin{block}{Correct Answer: D}
			Value creation is the main objective of risk governance and is achieved when the three underlying objectives (benefits realization, risk optimization and resource optimization) are balanced.
		\end{block}
		\begin{block}{Option Analysis}
			\begin{itemize}
				\item \textbf{A - Incorrect:} Benefits realization is one objective
				\item \textbf{B - Incorrect:} Risk optimization is one objective
				\item \textbf{C - Incorrect:} Resource optimization is one objective
				\item \textbf{D - Correct:} Value creation is ultimate objective
			\end{itemize}
		\end{block}
	\end{frame}
	
	\begin{frame}{Q95: Second Line of Defense Responsibility}
		\begin{block}{Question}
			Which activity would be a direct responsibility of the second line of defense?
		\end{block}
		\begin{enumerate}[(A)]
			\item Establish a risk management framework
			\item Identify, assess and respond to risk
			\item Implement information security controls
			\item Test and report control effectiveness
		\end{enumerate}
	\end{frame}
	
	\begin{frame}{Q95: Explanation}
		\begin{block}{Correct Answer: A}
			Establishing a risk management framework, providing awareness training, and supervising overall risk management are responsibilities of the second line of defense.
		\end{block}
		\begin{block}{Option Analysis}
			\begin{itemize}
				\item \textbf{A - Correct:} Framework establishment is second line
				\item \textbf{B - Incorrect:} This is first line operational management
				\item \textbf{C - Incorrect:} Implementation is first line
				\item \textbf{D - Incorrect:} Testing is third line
			\end{itemize}
		\end{block}
	\end{frame}
	
	\begin{frame}{Q96: Developing Ethical Behavior}
		\begin{block}{Question}
			Developing and practicing ethical behavior within an enterprise contributes the most to:
		\end{block}
		\begin{enumerate}[(A)]
			\item Determining risk appetite
			\item Building risk culture
			\item Enhancing risk management
			\item Increasing risk tolerance
		\end{enumerate}
	\end{frame}
	
	\begin{frame}{Q96: Explanation}
		\begin{block}{Correct Answer: B}
			Practicing professional ethics within an enterprise helps stakeholders in reviewing risk associated with unethical behavior before making decisions. This establishes the risk culture of the enterprise.
		\end{block}
		\begin{block}{Option Analysis}
			\begin{itemize}
				\item \textbf{A - Incorrect:} Ethics doesn't impact appetite
				\item \textbf{B - Correct:} Ethics builds risk culture
				\item \textbf{C - Incorrect:} Ethics doesn't enhance risk management
				\item \textbf{D - Incorrect:} Ethics doesn't impact tolerance
			\end{itemize}
		\end{block}
	\end{frame}
	
	\begin{frame}{Q97: Most Important Benefit of Professional Ethics}
		\begin{block}{Question}
			Which is MOST important benefit when stakeholders promote and follow professional ethics?
		\end{block}
		\begin{enumerate}[(A)]
			\item Increases reliability in risk management processes
			\item Increases risk management effectiveness within enterprise
			\item Improves compliance with regulatory requirements
			\item Improves enterprise's reputation and business growth
		\end{enumerate}
	\end{frame}
	
	\begin{frame}{Q97: Explanation}
		\begin{block}{Correct Answer: B}
			When stakeholders follow ethical practices, actions and decisions are based on reviewing the results and impact of these actions and decisions. This promotes a risk-based approach to decision-making and enhances overall risk management.
		\end{block}
		\begin{block}{Option Analysis}
			\begin{itemize}
				\item \textbf{A - Incorrect:} Ethics may reduce some vulnerabilities
				\item \textbf{B - Correct:} Enhances overall risk management
				\item \textbf{C - Incorrect:} Compliance is a result of enhanced risk management
				\item \textbf{D - Incorrect:} Reputation is a secondary benefit
			\end{itemize}
		\end{block}
	\end{frame}
	
	\begin{frame}{Q98: Proactive Ethics Approach}
		\begin{block}{Question}
			Which is BEST proactive approach for practicing professional ethics?
		\end{block}
		\begin{enumerate}[(A)]
			\item Develop an ethics policy
			\item Implement processes to report unethical behavior
			\item Define disciplinary action procedures
			\item Provide ethics awareness training
		\end{enumerate}
	\end{frame}
	
	\begin{frame}{Q98: Explanation}
		\begin{block}{Correct Answer: D}
			Awareness training is the most proactive approach to improve professional ethics within an enterprise; training will inform employees on the enterprise's code of ethics and ethics policies, and enable employees to make decisions and take appropriate action.
		\end{block}
		\begin{block}{Option Analysis}
			\begin{itemize}
				\item \textbf{A - Incorrect:} Policy without awareness is ineffective
				\item \textbf{B - Incorrect:} Reporting is reactive
				\item \textbf{C - Incorrect:} Discipline is reactive
				\item \textbf{D - Correct:} Awareness training is proactive
			\end{itemize}
		\end{block}
	\end{frame}
	
	\begin{frame}{Q99: Critical for 3LoD Effectiveness}
		\begin{block}{Question}
			Which is MOST critical to ensure three lines of defense work together effectively?
		\end{block}
		\begin{enumerate}[(A)]
			\item Documented disaster recovery plan
			\item Automated operational processes
			\item Independent external audits
			\item Clear, detailed roles and responsibilities
		\end{enumerate}
	\end{frame}
	
	\begin{frame}{Q99: Explanation}
		\begin{block}{Correct Answer: D}
			Without clear, detailed job descriptions for employees working in each line of defense, there can be confusion on what is expected at each level, which may hinder the effectiveness of each line of defense.
		\end{block}
		\begin{block}{Option Analysis}
			\begin{itemize}
				\item \textbf{A - Incorrect:} DR plan doesn't directly affect 3LoD
				\item \textbf{B - Incorrect:} Automation may or may not improve effectiveness
				\item \textbf{C - Incorrect:} Audits uncover issues but don't ensure effectiveness
				\item \textbf{D - Correct:} Clear roles are essential
			\end{itemize}
		\end{block}
	\end{frame}
	
	\begin{frame}{Q100: IT Plan Driver}
		\begin{block}{Question}
			Which choice should drive the IT plan?
		\end{block}
		\begin{enumerate}[(A)]
			\item Strategic planning and business requirements
			\item Technology and operational procedures
			\item Compliance with laws and regulations
			\item Project plans and stakeholder requirements
		\end{enumerate}
	\end{frame}
	
	\begin{frame}{Q100: Explanation}
		\begin{block}{Correct Answer: A}
			IT exists to support business objectives. Management of enterprise IT should align the IT plan closely with the business.
		\end{block}
		\begin{block}{Option Analysis}
			\begin{itemize}
				\item \textbf{A - Correct:} Business requirements drive IT plan
				\item \textbf{B - Incorrect:} Technology should not eclipse business strategy
				\item \textbf{C - Incorrect:} Compliance is evaluated like any other risk
				\item \textbf{D - Incorrect:} Project-based approach leads to duplication
			\end{itemize}
		\end{block}
	\end{frame}
	
	% ================ SECTION 3: DOMAIN 3 - RISK RESPONSE ================
	
	\section{Domain 3: Risk Response}
	
	\begin{frame}{Q101: Risk Response Report}
		\begin{block}{Question}
			A risk response report includes recommendations for:
		\end{block}
		\begin{enumerate}[(A)]
			\item Acceptance
			\item Assessment
			\item Evaluation
			\item Quantification
		\end{enumerate}
	\end{frame}
	
	\begin{frame}{Q101: Explanation}
		\begin{block}{Correct Answer: A}
			Acceptance of a risk is an alternative to be considered in the risk response process.
		\end{block}
		\begin{block}{Option Analysis}
			\begin{itemize}
				\item \textbf{A - Correct:} Acceptance is a risk response option
				\item \textbf{B - Incorrect:} Assessment is completed prior to risk response
				\item \textbf{C - Incorrect:} Evaluation is part of assessment
				\item \textbf{D - Incorrect:} Quantification is an input to risk response
			\end{itemize}
		\end{block}
	\end{frame}
	
	\begin{frame}{Q102: Reducing to Acceptable Risk}
		\begin{block}{Question}
			Which is MOST likely to be reduced to achieve acceptable risk?
		\end{block}
		\begin{enumerate}[(A)]
			\item Risk appetite
			\item Control risk
			\item Residual risk
			\item Inherent risk
		\end{enumerate}
	\end{frame}
	
	\begin{frame}{Q102: Explanation}
		\begin{block}{Correct Answer: C}
			Residual risk is the remaining risk after management has implemented a risk response. Acceptable risk is achieved when the residual risk is reduced to the levels within the enterprise's risk appetite.
		\end{block}
		\begin{block}{Option Analysis}
			\begin{itemize}
				\item \textbf{A - Incorrect:} Risk appetite doesn't change with mitigation
				\item \textbf{B - Incorrect:} Control risk is not necessarily related
				\item \textbf{C - Correct:} Residual risk must be reduced to acceptable level
				\item \textbf{D - Incorrect:} Inherent risk cannot be eliminated
			\end{itemize}
		\end{block}
	\end{frame}
	
	\begin{frame}{Q103: Risk Acceptance Reason}
		\begin{block}{Question}
			Global financial institution decided not to take further action on a denial-of-service vulnerability. MOST likely reason is:
		\end{block}
		\begin{enumerate}[(A)]
			\item Needed countermeasure is too complicated to deploy
			\item Sufficient safeguards are in place
			\item Likelihood of risk occurring is unknown
			\item Cost of countermeasure outweighs value of asset and potential loss
		\end{enumerate}
	\end{frame}
	
	\begin{frame}{Q103: Explanation}
		\begin{block}{Correct Answer: D}
			An enterprise may decide to accept a specific risk because the protection would cost more than the potential loss.
		\end{block}
		\begin{block}{Option Analysis}
			\begin{itemize}
				\item \textbf{A - Incorrect:} Complicated doesn't mean cost prohibitive
				\item \textbf{B - Incorrect:} Safeguards must match risk impact
				\item \textbf{C - Incorrect:} Frequency can't be predicted but not reason for acceptance
				\item \textbf{D - Correct:} Cost-benefit analysis justifies acceptance
			\end{itemize}
		\end{block}
	\end{frame}
	
	\begin{frame}{Q104: Cost-Benefit Analysis Component}
		\begin{block}{Question}
			Which is MOST relevant to include in a cost-benefit analysis of two-factor authentication system?
		\end{block}
		\begin{enumerate}[(A)]
			\item Approved budget of project
			\item Frequency of incidents
			\item Annual loss expectancy of security incidents
			\item Total cost of ownership
		\end{enumerate}
	\end{frame}
	
	\begin{frame}{Q104: Explanation}
		\begin{block}{Correct Answer: D}
			Total cost of ownership is the most relevant piece of information to be included in the cost-benefit analysis because it establishes a cost baseline that must be considered for the full life cycle of the control.
		\end{block}
		\begin{block}{Option Analysis}
			\begin{itemize}
				\item \textbf{A - Incorrect:} Budget may not reflect actual cost
				\item \textbf{B - Incorrect:} Indirect relationship to benefit
				\item \textbf{C - Incorrect:} Indirect relationship
				\item \textbf{D - Correct:} TCO establishes full life cycle cost baseline
			\end{itemize}
		\end{block}
	\end{frame}
	
	\begin{frame}{Q105: Outsourcing Contract Priority}
		\begin{block}{Question}
			Which is MOST important part of any outsourcing contract?
		\end{block}
		\begin{enumerate}[(A)]
			\item Right to audit outsourcing provider
			\item Provisions to assess compliance of provider
			\item Procedures for incident notification
			\item Requirements to encrypt hosted data
		\end{enumerate}
	\end{frame}
	
	\begin{frame}{Q105: Explanation}
		\begin{block}{Correct Answer: B}
			If a contract contains no provision to monitor and hold a supplier accountable for security, then the outsourcing enterprise cannot ensure compliance or proper handling of its data.
		\end{block}
		\begin{block}{Option Analysis}
			\begin{itemize}
				\item \textbf{A - Incorrect:} Provider may provide proof of compliance
				\item \textbf{B - Correct:} Compliance assessment is essential
				\item \textbf{C - Incorrect:} Contract usually doesn't contain procedure details
				\item \textbf{D - Incorrect:} Encryption may not be required for all data
			\end{itemize}
		\end{block}
	\end{frame}
	
	\begin{frame}{Q106: Cloud Provider Greatest Risk}
		\begin{block}{Question}
			Which poses GREATEST risk to enterprise that recently engaged services of cloud provider?
		\end{block}
		\begin{enumerate}[(A)]
			\item Cloud provider's facility is in same vicinity as subscriber
			\item Service level agreement is ambiguous
			\item References from other customers were not obtained
			\item Auditing vendor requires dependence on third-party audit firm
		\end{enumerate}
	\end{frame}
	
	\begin{frame}{Q106: Explanation}
		\begin{block}{Correct Answer: B}
			If the service level agreement is ambiguous, it will be difficult to determine whether the provider complies.
		\end{block}
		\begin{block}{Option Analysis}
			\begin{itemize}
				\item \textbf{A - Incorrect:} No direct impact if continuity plan is adequate
				\item \textbf{B - Correct:} Ambiguous SLA makes compliance verification difficult
				\item \textbf{C - Incorrect:} References are important but can't assure delivery
				\item \textbf{D - Incorrect:} Inability to audit is less than desirable but allowed under SSAE 18
			\end{itemize}
		\end{block}
	\end{frame}
	
	\begin{frame}{Q107: Risk Cannot Be Sufficiently Mitigated}
		\begin{block}{Question}
			When risk cannot be sufficiently mitigated through controls, which option BEST protects enterprise from potential financial impact?
		\end{block}
		\begin{enumerate}[(A)]
			\item Insuring against the risk
			\item Updating IT risk registry
			\item Improving staff training in risk area
			\item Outsourcing related process to third party
		\end{enumerate}
	\end{frame}
	
	\begin{frame}{Q107: Explanation}
		\begin{block}{Correct Answer: A}
			An insurance policy can compensate the enterprise monetarily for the impact of the risk by transferring the risk to the insurance company.
		\end{block}
		\begin{block}{Option Analysis}
			\begin{itemize}
				\item \textbf{A - Correct:} Insurance transfers financial risk
				\item \textbf{B - Incorrect:} Updating registry doesn't change risk
				\item \textbf{C - Incorrect:} Training may reduce but not completely mitigate
				\item \textbf{D - Incorrect:} Outsourcing doesn't necessarily remove risk
			\end{itemize}
		\end{block}
	\end{frame}
	
	\begin{frame}{Q108: Cost Exceeds Benefit}
		\begin{block}{Question}
			Cost to mitigate risk is much greater than benefit. BEST risk response is:
		\end{block}
		\begin{enumerate}[(A)]
			\item Treated
			\item Terminated
			\item Accepted
			\item Transferred
		\end{enumerate}
	\end{frame}
	
	\begin{frame}{Q108: Explanation}
		\begin{block}{Correct Answer: C}
			When the cost of control is more than the cost of the potential impact, the risk should be accepted.
		\end{block}
		\begin{block}{Option Analysis}
			\begin{itemize}
				\item \textbf{A - Incorrect:} Treatment costs more than benefit
				\item \textbf{B - Incorrect:} Risk can be avoided, not terminated
				\item \textbf{C - Correct:} Acceptance is appropriate when cost > benefit
				\item \textbf{D - Incorrect:} Transfer is limited benefit if cost > benefit
			\end{itemize}
		\end{block}
	\end{frame}
	
	\begin{frame}{Q109: Risk Scenario for Service Disruptions}
		\begin{block}{Question}
			Senior management at data center provider wants to improve continuity of operations. Which risk management activity will MOST benefit?
		\end{block}
		\begin{enumerate}[(A)]
			\item Develop key risk indicators
			\item Determine risk and control ownership
			\item Update risk register
			\item Develop IT risk scenarios
		\end{enumerate}
	\end{frame}
	
	\begin{frame}{Q109: Explanation}
		\begin{block}{Correct Answer: D}
			The relationship between IT risk scenarios and business impact must be established to understand the effects of possible adverse events on enterprise objectives.
		\end{block}
		\begin{block}{Option Analysis}
			\begin{itemize}
				\item \textbf{A - Incorrect:} KRIs help understand potential changes
				\item \textbf{B - Incorrect:} Ownership is important but scenarios identify circumstances
				\item \textbf{C - Incorrect:} Scenarios specifically correlate adverse events to business outcomes
				\item \textbf{D - Correct:} Risk scenarios connect IT risk to business impact
			\end{itemize}
		\end{block}
	\end{frame}
	
	\begin{frame}{Q110: Reciprocal Agreement}
		\begin{block}{Question}
			Critical IT equipment delivered several days late due to flooding. Reciprocal agreement exists with another company for replacement. This is risk:
		\end{block}
		\begin{enumerate}[(A)]
			\item Transfer
			\item Avoidance
			\item Acceptance
			\item Mitigation
		\end{enumerate}
	\end{frame}
	
	\begin{frame}{Q110: Explanation}
		\begin{block}{Correct Answer: D}
			Risk mitigation attempts to reduce the impact when a risk event occurs. Making plans such as a reciprocal arrangement with another company reduces the consequence of the risk event.
		\end{block}
		\begin{block}{Option Analysis}
			\begin{itemize}
				\item \textbf{A - Incorrect:} Not transferred using insurance or other strategy
				\item \textbf{B - Incorrect:} Arranging standby is mitigation, not avoidance
				\item \textbf{C - Incorrect:} Risk is not accepted
				\item \textbf{D - Correct:} Reciprocal arrangement reduces impact
			\end{itemize}
		\end{block}
	\end{frame}
	
	\begin{frame}{Q111: Selecting Risk Response}
		\begin{block}{Question}
			Which BEST helps enterprise select an appropriate risk response?
		\end{block}
		\begin{enumerate}[(A)]
			\item Analysis of change in risk environment
			\item Analysis of risk that can be transferred
			\item Analysis of likelihood and impact of risk scenarios
			\item Analysis of control costs and benefits
		\end{enumerate}
	\end{frame}
	
	\begin{frame}{Q111: Explanation}
		\begin{block}{Correct Answer: D}
			An analysis of costs and benefits for controls helps an enterprise understand if it can mitigate the risk to an acceptable level.
		\end{block}
		\begin{block}{Option Analysis}
			\begin{itemize}
				\item \textbf{A - Incorrect:} Change in environment doesn't provide cost/benefit
				\item \textbf{B - Incorrect:} Risk can never be eliminated
				\item \textbf{C - Incorrect:} Likelihood/impact help establish risk level, not response
				\item \textbf{D - Correct:} Cost-benefit analysis helps select response
			\end{itemize}
		\end{block}
	\end{frame}
	
	\begin{frame}{Q112: Optimal Return on Security Investment}
		\begin{block}{Question}
			Which leads to BEST optimal return on security investment?
		\end{block}
		\begin{enumerate}[(A)]
			\item Deploying maximum security protection across all assets
			\item Focusing on most important information assets and determining their protection
			\item Deploying minimum protection across all assets
			\item Investing only after major security incident
		\end{enumerate}
	\end{frame}
	
	\begin{frame}{Q112: Explanation}
		\begin{block}{Correct Answer: B}
			To optimize return on security investment, the primary focus should be identifying the important information assets and protecting them appropriately to optimize investment (i.e., important information assets get more protection than less important or critical assets).
		\end{block}
		\begin{block}{Option Analysis}
			\begin{itemize}
				\item \textbf{A - Incorrect:} Overprotects less critical assets
				\item \textbf{B - Correct:} Focus on important assets optimizes investment
				\item \textbf{C - Incorrect:} Compromises security of critical assets
				\item \textbf{D - Incorrect:} Reactive approach may severely compromise operations
			\end{itemize}
		\end{block}
	\end{frame}
	
	\begin{frame}{Q113: Reducing Impact of Compromised Admin Account}
		\begin{block}{Question}
			Which control can reduce potential impact of malicious hacker who gains access to administrator account?
		\end{block}
		\begin{enumerate}[(A)]
			\item Multifactor authentication
			\item Audit logging
			\item Least privilege
			\item Password policy
		\end{enumerate}
	\end{frame}
	
	\begin{frame}{Q113: Explanation}
		\begin{block}{Correct Answer: C}
			Least privilege can reduce the impact of a compromised account within the scope of its intended use and limit impact to the enterprise as a whole, without restricting the performance of employees in administrative roles.
		\end{block}
		\begin{block}{Option Analysis}
			\begin{itemize}
				\item \textbf{A - Incorrect:} Control has already been bypassed
				\item \textbf{B - Incorrect:} Audit logging is lagging indicator
				\item \textbf{C - Correct:} Least privilege limits scope of damage
				\item \textbf{D - Incorrect:} Password changes may not be frequent enough
			\end{itemize}
		\end{block}
	\end{frame}
	
	\begin{frame}{Q114: Critical for Effective Risk Management}
		\begin{block}{Question}
			Which is critical to risk practitioner for effective risk management program?
		\end{block}
		\begin{enumerate}[(A)]
			\item Risk response strategy
			\item Risk register content
			\item Risk profile content
			\item Risk owners and accountability
		\end{enumerate}
	\end{frame}
	
	\begin{frame}{Q114: Explanation}
		\begin{block}{Correct Answer: D}
			The identification of risk owners is critical because risk owners must make informed and cost-effective business decisions regarding appropriate controls to mitigate their owned risk. Risk response strategy, risk registers and risk profiles are tools that require ownership by users who apply them proactively and with accountability.
		\end{block}
		\begin{block}{Option Analysis}
			\begin{itemize}
				\item \textbf{A - Incorrect:} Strategy provides options but owners are critical
				\item \textbf{B - Incorrect:} Register identifies risk but owners are critical
				\item \textbf{C - Incorrect:} Profile describes risk but owners are critical
				\item \textbf{D - Correct:} Risk owners are essential for accountability
			\end{itemize}
		\end{block}
	\end{frame}
	
	\begin{frame}{Q115: Anti-Malware Cost Exceeds Loss Expectancy}
		\begin{block}{Question}
			Cost of anti-malware exceeds loss expectancy of malware threats. What is MOST viable risk response?
		\end{block}
		\begin{enumerate}[(A)]
			\item Risk elimination
			\item Risk acceptance
			\item Risk transfer
			\item Risk mitigation
		\end{enumerate}
	\end{frame}
	
	\begin{frame}{Q115: Explanation}
		\begin{block}{Correct Answer: B}
			When the cost of a risk response exceeds the loss expectancy, the most viable risk response is risk acceptance.
		\end{block}
		\begin{block}{Option Analysis}
			\begin{itemize}
				\item \textbf{A - Incorrect:} Risk elimination is not possible
				\item \textbf{B - Correct:} Acceptance when cost exceeds expected loss
				\item \textbf{C - Incorrect:} Transfer also based on cost-benefit analysis
				\item \textbf{D - Incorrect:} Mitigation also based on cost-benefit analysis
			\end{itemize}
		\end{block}
	\end{frame}
	
	\begin{frame}{Q116: Risk Avoidance Example}
		\begin{block}{Question}
			Which is BEST example of risk avoidance behavior?
		\end{block}
		\begin{enumerate}[(A)]
			\item Taking no action against the risk
			\item Outsourcing the related process
			\item Insuring against a specific event
			\item Exiting the process that gives rise to risk
		\end{enumerate}
	\end{frame}
	
	\begin{frame}{Q116: Explanation}
		\begin{block}{Correct Answer: D}
			Avoidance means exiting the activities or conditions that give rise to risk.
		\end{block}
		\begin{block}{Option Analysis}
			\begin{itemize}
				\item \textbf{A - Incorrect:} This is acceptance
				\item \textbf{B - Incorrect:} This is transfer/sharing
				\item \textbf{C - Incorrect:} This is transfer/sharing
				\item \textbf{D - Correct:} Exiting process is avoidance
			\end{itemize}
		\end{block}
	\end{frame}
	
	\begin{frame}{Q117: Critical Dependencies in Risk Treatment}
		\begin{block}{Question}
			Which would ensure that critical dependencies are addressed in the risk treatment plan?
		\end{block}
		\begin{enumerate}[(A)]
			\item Implement risk treatment strategy for all possible risk
			\item Verify accomplishment of business objectives through top-down process review
			\item Treat each risk independently
			\item Verify accomplishment through bottom-up process review
		\end{enumerate}
	\end{frame}
	
	\begin{frame}{Q117: Explanation}
		\begin{block}{Correct Answer: B}
			Development of an overall risk treatment strategy should be a top-down process, driven jointly by the need to achieve business objectives and to apply economically feasible constraints while controlling uncertainty within the enterprise risk appetite.
		\end{block}
		\begin{block}{Option Analysis}
			\begin{itemize}
				\item \textbf{A - Incorrect:} Not economically feasible for all risk
				\item \textbf{B - Correct:} Top-down ensures business objectives drive treatment
				\item \textbf{C - Incorrect:} Independent treatment creates gaps
				\item \textbf{D - Incorrect:} Bottom-up doesn't ensure dependencies addressed
			\end{itemize}
		\end{block}
	\end{frame}
	
	\begin{frame}{Q118: Anti-Malware as Risk Response}
		\begin{block}{Question}
			CISO recommends controls like anti-malware to protect IS. Which risk handling approach is this?
		\end{block}
		\begin{enumerate}[(A)]
			\item Risk transference
			\item Risk mitigation
			\item Risk acceptance
			\item Risk avoidance
		\end{enumerate}
	\end{frame}
	
	\begin{frame}{Q118: Explanation}
		\begin{block}{Correct Answer: B}
			By implementing controls, the company is trying to reduce the risk to an acceptable level, thereby mitigating risk.
		\end{block}
		\begin{block}{Option Analysis}
			\begin{itemize}
				\item \textbf{A - Incorrect:} Transfer involves another entity
				\item \textbf{B - Correct:} Controls reduce risk to acceptable level
				\item \textbf{C - Incorrect:} Acceptance is taking no action
				\item \textbf{D - Incorrect:} Avoidance stops activity causing risk
			\end{itemize}
		\end{block}
	\end{frame}
	
	\begin{frame}{Q119: Outsourcing Risk Mitigation}
		\begin{block}{Question}
			Which practice BEST mitigates risk associated with outsourcing a business function?
		\end{block}
		\begin{enumerate}[(A)]
			\item Performing audits to verify compliance with contract
			\item Requiring all vendor staff to complete annual awareness training
			\item Retaining copies of all sensitive data on internal systems
			\item Reviewing financial records of vendor to verify financial soundness
		\end{enumerate}
	\end{frame}
	
	\begin{frame}{Q119: Explanation}
		\begin{block}{Correct Answer: A}
			When an outsourcing relationship is established, the risk of noncompliance with the agreement must be met through review, monitoring and enforcement of the contract terms. Therefore, conducting regular audits to verify that the vendor is compliant with contract requirements is an important practice.
		\end{block}
		\begin{block}{Option Analysis}
			\begin{itemize}
				\item \textbf{A - Correct:} Audits ensure contract compliance
				\item \textbf{B - Incorrect:} Training doesn't provide same level of mitigation
				\item \textbf{C - Incorrect:} Duplicating data increases liability
				\item \textbf{D - Incorrect:} Financial review is important but not best practice
			\end{itemize}
		\end{block}
	\end{frame}
	
	\begin{frame}{Q120: Risk Transfer Best Addressed}
		\begin{block}{Question}
			Which is BEST addressed by transferring risk?
		\end{block}
		\begin{enumerate}[(A)]
			\item Antiquated fire suppression system in computer room
			\item Threat of disgruntled employee sabotage
			\item Threat of disgruntled employee theft
			\item Building located in 100-year flood plain
		\end{enumerate}
	\end{frame}
	
	\begin{frame}{Q120: Explanation}
		\begin{block}{Correct Answer: D}
			Purchasing an insurance policy transfers the risk of a flood. Risk transfer is the process of assigning risk to another entity, usually through the purchase of an insurance policy or through outsourcing the service.
		\end{block}
		\begin{block}{Option Analysis}
			\begin{itemize}
				\item \textbf{A - Incorrect:} Should update fire suppression system
				\item \textbf{B - Incorrect:} Not readily transferable
				\item \textbf{C - Incorrect:} Not readily transferable
				\item \textbf{D - Correct:} Flood risk can be transferred through insurance
			\end{itemize}
		\end{block}
	\end{frame}
	
	\begin{frame}{Q121: Finalizing Risk Treatment Plan}
		\begin{block}{Question}
			Which would BEST help finalize the risk treatment plan?
		\end{block}
		\begin{enumerate}[(A)]
			\item Vulnerability analysis
			\item Impact analysis
			\item Cost-benefit analysis
			\item SWOT analysis
		\end{enumerate}
	\end{frame}
	
	\begin{frame}{Q121: Explanation}
		\begin{block}{Correct Answer: C}
			A cost-benefit analysis helps determine if the benefit of a control outweighs the cost of implementing the control.
		\end{block}
		\begin{block}{Option Analysis}
			\begin{itemize}
				\item \textbf{A - Incorrect:} Vulnerability analysis identifies risk to treat
				\item \textbf{B - Incorrect:} Impact analysis is part of risk assessment
				\item \textbf{C - Correct:} Cost-benefit helps evaluate treatment options
				\item \textbf{D - Incorrect:} SWOT results must be translated into risk terms
			\end{itemize}
		\end{block}
	\end{frame}
	
	\begin{frame}{Q122: Risk Treatment Plan Greatest Benefit}
		\begin{block}{Question}
			GREATEST benefit of implementing a risk treatment plan is:
		\end{block}
		\begin{enumerate}[(A)]
			\item To reduce impact and likelihood of risk occurrence
			\item To identify unmitigated risk that can be transferred
			\item To exploit risk to test organizational preparedness
			\item To enhance overall risk appetite
		\end{enumerate}
	\end{frame}
	
	\begin{frame}{Q122: Explanation}
		\begin{block}{Correct Answer: A}
			Implementing the risk treatment plan reduces the negative impact and likelihood of a risk occurrence.
		\end{block}
		\begin{block}{Option Analysis}
			\begin{itemize}
				\item \textbf{A - Correct:} Reduces impact and likelihood
				\item \textbf{B - Incorrect:} Transfer is not the only option
				\item \textbf{C - Incorrect:} Exploiting risk creates more risk
				\item \textbf{D - Incorrect:} Risk appetite is established, not influenced
			\end{itemize}
		\end{block}
	\end{frame}
	
	\begin{frame}{Q123: Validating Patching Program}
		\begin{block}{Question}
			What is BEST method to validate effectiveness of enterprise's patching program?
		\end{block}
		\begin{enumerate}[(A)]
			\item Conduct penetration testing
			\item Conduct risk identification initiative
			\item Carry out vulnerability scans
			\item Review requests for change
		\end{enumerate}
	\end{frame}
	
	\begin{frame}{Q123: Explanation}
		\begin{block}{Correct Answer: C}
			Performing vulnerability scans will enable the IT risk practitioner to determine if patches are being installed on a timely basis.
		\end{block}
		\begin{block}{Option Analysis}
			\begin{itemize}
				\item \textbf{A - Incorrect:} Penetration testing could elevate risk
				\item \textbf{B - Incorrect:} Risk identification identifies new risk, not patching effectiveness
				\item \textbf{C - Correct:} Vulnerability scans show patch status
				\item \textbf{D - Incorrect:} RFCs don't mean patches are applied
			\end{itemize}
		\end{block}
	\end{frame}
	
	\begin{frame}{Q124: Designing IS Controls}
		\begin{block}{Question}
			Which is MOST important factor when designing information systems controls in complex environment?
		\end{block}
		\begin{enumerate}[(A)]
			\item Development methodologies
			\item Scalability of solution
			\item Technical platform interfaces
			\item Stakeholder requirements
		\end{enumerate}
	\end{frame}
	
	\begin{frame}{Q124: Explanation}
		\begin{block}{Correct Answer: D}
			The most important factor when designing information systems controls is their ability to advance the interests of the business by addressing stakeholder requirements.
		\end{block}
		\begin{block}{Option Analysis}
			\begin{itemize}
				\item \textbf{A - Incorrect:} Important but not as important as stakeholder requirements
				\item \textbf{B - Incorrect:} Important but not as important
				\item \textbf{C - Incorrect:} Important but not as important
				\item \textbf{D - Correct:} Stakeholder requirements drive control design
			\end{itemize}
		\end{block}
	\end{frame}
	
	\begin{frame}{Q125: EFT System Testing}
		\begin{block}{Question}
			Risk practitioner would MOST benefit from a test that:
		\end{block}
		\begin{enumerate}[(A)]
			\item Identifies introduction of potential new gaps in security
			\item Verifies adequate system recovery in case of failure
			\item Ensures system performs to expectation
			\item Ensures system can support volume of transactions
		\end{enumerate}
	\end{frame}
	
	\begin{frame}{Q125: Explanation}
		\begin{block}{Correct Answer: A}
			The security of an electronic funds transfer (EFT) system would be of the most interest to a risk practitioner. Testing should identify the introduction of potential new gaps in information security.
		\end{block}
		\begin{block}{Option Analysis}
			\begin{itemize}
				\item \textbf{A - Correct:} Security gaps are most important to risk practitioner
				\item \textbf{B - Incorrect:} Recovery is not as high priority
				\item \textbf{C - Incorrect:} Performance is not as high priority
				\item \textbf{D - Incorrect:} Stress testing is not as high priority
			\end{itemize}
		\end{block}
	\end{frame}
	
	\begin{frame}{Q126: Creating KRIs for Reporting}
		\begin{block}{Question}
			What is BEST approach for creating KRIs for quarterly reporting to senior leadership?
		\end{block}
		\begin{enumerate}[(A)]
			\item Survey senior leaders about primary risk concerns
			\item Identify list of most common network vulnerabilities
			\item Determine which KRIs are used in similar industry
			\item Identify enterprise risk appetite and metrics of current risk
		\end{enumerate}
	\end{frame}
	
	\begin{frame}{Q126: Explanation}
		\begin{block}{Correct Answer: D}
			Key risk indicators (KRIs) are measurable and quantifiable indicators of risk to assist in the monitoring and appropriate risk treatment proactively.
		\end{block}
		\begin{block}{Option Analysis}
			\begin{itemize}
				\item \textbf{A - Incorrect:} Feedback may be subjective and immeasurable
				\item \textbf{B - Incorrect:} Vulnerabilities lack likelihood and impact
				\item \textbf{C - Incorrect:} KRIs should be based on enterprise's risk appetite
				\item \textbf{D - Correct:} Risk appetite and current risk measures drive KRIs
			\end{itemize}
		\end{block}
	\end{frame}
	
	\begin{frame}{Q127: Gauging Risk for Stakeholders}
		\begin{block}{Question}
			Which can serve to gauge risk that may trigger stakeholder concern?
		\end{block}
		\begin{enumerate}[(A)]
			\item Indicators with approved thresholds
			\item Approved status report of completed milestones
			\item List of controls to be implemented by supplier
			\item Number of findings in external audit reports
		\end{enumerate}
	\end{frame}
	
	\begin{frame}{Q127: Explanation}
		\begin{block}{Correct Answer: A}
			Indicators with approved thresholds demonstrate the acceptable risk levels stakeholders are willing to tolerate, and any risk above those approved levels will likely trigger stakeholder concern.
		\end{block}
		\begin{block}{Option Analysis}
			\begin{itemize}
				\item \textbf{A - Correct:} Thresholds show acceptable risk levels
				\item \textbf{B - Incorrect:} Status report shows completed milestones only
				\item \textbf{C - Incorrect:} Controls list is requirements, not thresholds
				\item \textbf{D - Incorrect:} Audit findings are not indicators of tolerance
			\end{itemize}
		\end{block}
	\end{frame}
	
	\begin{frame}{Q128: Ensuring Access Revocation}
		\begin{block}{Question}
			Which control within user provision process BEST ensures revocation of system access for contractors when no longer required?
		\end{block}
		\begin{enumerate}[(A)]
			\item Log all account usage and send to managers
			\item Establish predetermined, automatic expiration dates
			\item Ensure each user signs security acknowledgment
			\item Require managers to email security when user leaves
		\end{enumerate}
	\end{frame}
	
	\begin{frame}{Q128: Explanation}
		\begin{block}{Correct Answer: B}
			Predetermined expiration dates are the most effective means of removing systems access for temporary users.
		\end{block}
		\begin{block}{Option Analysis}
			\begin{itemize}
				\item \textbf{A - Incorrect:} Logging is detective, not as effective as preventive
				\item \textbf{B - Correct:} Automatic expiration is most effective
				\item \textbf{C - Incorrect:} Signing acknowledgment has little effect
				\item \textbf{D - Incorrect:} Managers often don't submit termination notices promptly
			\end{itemize}
		\end{block}
	\end{frame}
	
	\begin{frame}{Q129: Presenting Risk Profile to Board}
		\begin{block}{Question}
			Which BEST helps while presenting current risk profile to executive management and board?
		\end{block}
		\begin{enumerate}[(A)]
			\item Risk response dashboard
			\item Emerging risk report
			\item Risk register dashboard
			\item Key risk indicators report
		\end{enumerate}
	\end{frame}
	
	\begin{frame}{Q129: Explanation}
		\begin{block}{Correct Answer: C}
			A risk register dashboard would provide a comprehensive overview of the risk profile of the enterprise.
		\end{block}
		\begin{block}{Option Analysis}
			\begin{itemize}
				\item \textbf{A - Incorrect:} Risk response is one component only
				\item \textbf{B - Incorrect:} Emerging risk report would not be complete
				\item \textbf{C - Correct:} Risk register dashboard provides comprehensive view
				\item \textbf{D - Incorrect:} KRI report is only one component
			\end{itemize}
		\end{block}
	\end{frame}
	
	\begin{frame}{Q130: Main Reason for KCI Trends}
		\begin{block}{Question}
			Which is MAIN reason senior management monitors and analyzes trends in key control indicators?
		\end{block}
		\begin{enumerate}[(A)]
			\item Provides feedback on overall control environment
			\item Helps in identifying redundant controls
			\item Proactively identifies impacts to risk profile
			\item Helps determine if additional controls required
		\end{enumerate}
	\end{frame}
	
	\begin{frame}{Q130: Explanation}
		\begin{block}{Correct Answer: C}
			The primary objective of key control indicators (KCIs) is to ensure that controls actually mitigate risk at an effective level. Analysis of KCI trends provides information on the overall effectiveness of controls and provides management information on the status of risk management.
		\end{block}
		\begin{block}{Option Analysis}
			\begin{itemize}
				\item \textbf{A - Incorrect:} Control environment is operations responsibility
				\item \textbf{B - Incorrect:} Redundancy is design/operations responsibility
				\item \textbf{C - Correct:} KCIs proactively identify risk profile impacts
				\item \textbf{D - Incorrect:} Determining additional controls is operations responsibility
			\end{itemize}
		\end{block}
	\end{frame}
	
	\begin{frame}{Q131: Risk Treatment Plans Purpose}
		\begin{block}{Question}
			Risk treatment plans are necessary to describe how the:
		\end{block}
		\begin{enumerate}[(A)]
			\item Identified risk is further analyzed
			\item Chosen treatment options will be implemented
			\item Accepted risk is treated
			\item Risk indicators will monitor the risk
		\end{enumerate}
	\end{frame}
	
	\begin{frame}{Q131: Explanation}
		\begin{block}{Correct Answer: B}
			A risk treatment plan includes the plan of action for the chosen treatment, how it will be implemented, who will implement it, key dates and resource requirements.
		\end{block}
		\begin{block}{Option Analysis}
			\begin{itemize}
				\item \textbf{A - Incorrect:} Treatment plan doesn't further analyze risk
				\item \textbf{B - Correct:} Plan describes implementation
				\item \textbf{C - Incorrect:} Accepted risk doesn't need treatment plan
				\item \textbf{D - Incorrect:} Risk indicators are not part of treatment plan
			\end{itemize}
		\end{block}
	\end{frame}
	
	\begin{frame}{Q132: Cost-Benefit Analysis Timing}
		\begin{block}{Question}
			In risk management process, cost-benefit analysis is MAINLY performed:
		\end{block}
		\begin{enumerate}[(A)]
			\item As part of initial risk assessment
			\item As part of risk-response planning
			\item During information asset valuation
			\item When insurance is calculated for risk transfer
		\end{enumerate}
	\end{frame}
	
	\begin{frame}{Q132: Explanation}
		\begin{block}{Correct Answer: B}
			In risk response, a range of controls that can mitigate risk will be identified; however, a cost-benefit analysis in this process will help identify the right controls to address the risk at acceptable levels within the budget.
		\end{block}
		\begin{block}{Option Analysis}
			\begin{itemize}
				\item \textbf{A - Incorrect:} Performed not just once but every time controls needed
				\item \textbf{B - Correct:} During risk-response planning
				\item \textbf{C - Incorrect:} Value determined by business importance
				\item \textbf{D - Incorrect:} Not the stage where cost-benefit is performed
			\end{itemize}
		\end{block}
	\end{frame}
	
	\begin{frame}{Q133: Insufficient Risk Treatment Resources}
		\begin{block}{Question}
			In event available resources for risk treatment are not sufficient, the risk treatment plan should:
		\end{block}
		\begin{enumerate}[(A)]
			\item Define priorities across all treatments to assist in resource allocation
			\item Recommend postponing treatment until resources available
			\item Recommend reassessing risk treatment options
			\item Suggest increasing priority to ensure resource availability
		\end{enumerate}
	\end{frame}
	
	\begin{frame}{Q133: Explanation}
		\begin{block}{Correct Answer: A}
			When the available resources for treating the risk are not sufficient, the risk treatment plan should include prioritization of the risk so resources can be properly allocated.
		\end{block}
		\begin{block}{Option Analysis}
			\begin{itemize}
				\item \textbf{A - Correct:} Prioritization drives resource allocation
				\item \textbf{B - Incorrect:} May not be most appropriate option
				\item \textbf{C - Incorrect:} Prioritization already determined
				\item \textbf{D - Incorrect:} Doesn't consider other priorities
			\end{itemize}
		\end{block}
	\end{frame}
	
	\begin{frame}{Q134: Backup/Restore Control Classification}
		\begin{block}{Question}
			System backup and restore procedures can BEST be classified as:
		\end{block}
		\begin{enumerate}[(A)]
			\item Technical controls
			\item Detective controls
			\item Corrective controls
			\item Deterrent controls
		\end{enumerate}
	\end{frame}
	
	\begin{frame}{Q134: Explanation}
		\begin{block}{Correct Answer: C}
			Corrective controls remediate the impact from negative events. If a system suffers harm so extensive that processing cannot continue, backup restore procedures enable that system to be recovered.
		\end{block}
		\begin{block}{Option Analysis}
			\begin{itemize}
				\item \textbf{A - Incorrect:} Technical controls are hardware/software safeguards
				\item \textbf{B - Incorrect:} Detective controls identify violations
				\item \textbf{C - Correct:} Backup/restore corrects impact
				\item \textbf{D - Incorrect:} Deterrent controls discourage compromise
			\end{itemize}
		\end{block}
	\end{frame}
	
	\begin{frame}{Q135: Policy Exception Process Reason}
		\begin{block}{Question}
			PRIMARY reason for initiating a policy-exception process is:
		\end{block}
		\begin{enumerate}[(A)]
			\item Risk is justified by the benefit
			\item Policy compliance is difficult to enforce
			\item Operations are too busy to comply
			\item Users may initially be inconvenienced
		\end{enumerate}
	\end{frame}
	
	\begin{frame}{Q135: Explanation}
		\begin{block}{Correct Answer: A}
			Exceptions to policies are warranted if the benefits outweigh the costs of policy compliance; however, the enterprise needs to assess both the tangible and intangible risk and evaluate both in the context of existing risk.
		\end{block}
		\begin{block}{Option Analysis}
			\begin{itemize}
				\item \textbf{A - Correct:} Benefit justifies exception
				\item \textbf{B - Incorrect:} Difficulty doesn't justify exception
				\item \textbf{C - Incorrect:} Lack of resources doesn't justify exception
				\item \textbf{D - Incorrect:} Inconvenience doesn't justify exception
			\end{itemize}
		\end{block}
	\end{frame}
	
	\begin{frame}{Q136: Measuring Operational Control Effectiveness}
		\begin{block}{Question}
			Which would BEST measure effectiveness of operational controls?
		\end{block}
		\begin{enumerate}[(A)]
			\item Control matrix
			\item Key performance indicator
			\item Statement of applicability
			\item Key control indicator
		\end{enumerate}
	\end{frame}
	
	\begin{frame}{Q136: Explanation}
		\begin{block}{Correct Answer: D}
			Key control indicators, also referred to as control effectiveness indicators, are metrics that provide information on the degree to which a control is working.
		\end{block}
		\begin{block}{Option Analysis}
			\begin{itemize}
				\item \textbf{A - Incorrect:} Control matrix is analysis tool
				\item \textbf{B - Incorrect:} KPIs don't measure control effectiveness
				\item \textbf{C - Incorrect:} SOA is ISO 27001 specific
				\item \textbf{D - Correct:} KCIs measure control effectiveness
			\end{itemize}
		\end{block}
	\end{frame}
	
	\begin{frame}{Q137: Security Technology Selection Basis}
		\begin{block}{Question}
			Security technologies should be selected PRIMARILY on basis of:
		\end{block}
		\begin{enumerate}[(A)]
			\item Evaluation in security publications
			\item Compliance with industry standards
			\item Ability to mitigate risk to organizational objectives
			\item Cost compared to IT budget
		\end{enumerate}
	\end{frame}
	
	\begin{frame}{Q137: Explanation}
		\begin{block}{Correct Answer: C}
			The most fundamental criterion for selecting security technology is the capacity to reduce risk for organizational objectives.
		\end{block}
		\begin{block}{Option Analysis}
			\begin{itemize}
				\item \textbf{A - Incorrect:} Publication evaluation is valuable but secondary
				\item \textbf{B - Incorrect:} Industry compliance is secondary
				\item \textbf{C - Correct:} Risk mitigation is primary
				\item \textbf{D - Incorrect:} Cost is an important but secondary consideration
			\end{itemize}
		\end{block}
	\end{frame}
	
	\begin{frame}{Q138: Black Box Penetration Test Preparation}
		\begin{block}{Question}
			Which should be in place before a black box penetration test begins?
		\end{block}
		\begin{enumerate}[(A)]
			\item Clearly stated definition of scope
			\item Previous test results
			\item Proper communication and awareness training
			\item Incident response plan
		\end{enumerate}
	\end{frame}
	
	\begin{frame}{Q138: Explanation}
		\begin{block}{Correct Answer: A}
			A clearly stated definition of scope ensures a proper understanding of risk and success criteria.
		\end{block}
		\begin{block}{Option Analysis}
			\begin{itemize}
				\item \textbf{A - Correct:} Scope definition is essential
				\item \textbf{B - Incorrect:} Previous results may help define scope
				\item \textbf{C - Incorrect:} Communication/training not necessary
				\item \textbf{D - Incorrect:} IR plan not necessary
			\end{itemize}
		\end{block}
	\end{frame}
	
	\begin{frame}{Q139: Securing Network Against External Attack}
		\begin{block}{Question}
			Which is BEST way to ensure corporate network is adequately secured against external attack?
		\end{block}
		\begin{enumerate}[(A)]
			\item Use intrusion detection system
			\item Establish minimum security baselines
			\item Implement vendor recommended settings
			\item Perform periodic penetration testing
		\end{enumerate}
	\end{frame}
	
	\begin{frame}{Q139: Explanation}
		\begin{block}{Correct Answer: D}
			Penetration testing is the best way to ensure that perimeter security is adequate.
		\end{block}
		\begin{block}{Option Analysis}
			\begin{itemize}
				\item \textbf{A - Incorrect:} IDS detects but doesn't confirm security
				\item \textbf{B - Incorrect:} Baselines are beneficial but not sufficient
				\item \textbf{C - Incorrect:} Vendor settings are beneficial but not sufficient
				\item \textbf{D - Correct:} Penetration testing provides highest assurance
			\end{itemize}
		\end{block}
	\end{frame}
	
	\begin{frame}{Q140: Establishing Key Control Indicators}
		\begin{block}{Question}
			Which is PRIMARILY defined before establishing key control indicators?
		\end{block}
		\begin{enumerate}[(A)]
			\item Desired capacity
			\item Desired risk threshold
			\item Desired tolerances
			\item Desired risk appetite
		\end{enumerate}
	\end{frame}
	
	\begin{frame}{Q140: Explanation}
		\begin{block}{Correct Answer: C}
			The goal of KCIs is to track performance of control actions relative to tolerances, providing insight into the ongoing adequacy of a given control in keeping risk within acceptable levels; for this reason, KCIs are sometimes called control effectiveness indicators. Therefore, tolerances would need to be defined before a KCI could be established.
		\end{block}
		\begin{block}{Option Analysis}
			\begin{itemize}
				\item \textbf{A - Incorrect:} Capacity management is separate
				\item \textbf{B - Incorrect:} KCIs are not related to risk thresholds
				\item \textbf{C - Correct:} Tolerances must be defined first
				\item \textbf{D - Incorrect:} KCIs are not related to risk appetite
			\end{itemize}
		\end{block}
	\end{frame}
	
	\begin{frame}{Q141: Substantive Test for Tape Inventory}
		\begin{block}{Question}
			A substantive test to verify tape library inventory records are accurate involves:
		\end{block}
		\begin{enumerate}[(A)]
			\item Determining whether bar code readers are installed
			\item Conducting physical count of tape inventory
			\item Checking whether receipts and issues are accurately recorded
			\item Determining whether movement of tapes is authorized
		\end{enumerate}
	\end{frame}
	
	\begin{frame}{Q141: Explanation}
		\begin{block}{Correct Answer: B}
			A substantive test involves gathering evidence to evaluate the integrity of individual transactions, data or other information. Conducting a physical count of the tape inventory is a substantive test.
		\end{block}
		\begin{block}{Option Analysis}
			\begin{itemize}
				\item \textbf{A - Incorrect:} Testing bar code readers is compliance test
				\item \textbf{B - Correct:} Physical count is substantive test
				\item \textbf{C - Incorrect:} Accuracy check is compliance test
				\item \textbf{D - Incorrect:} Authorization check is compliance test
			\end{itemize}
		\end{block}
	\end{frame}
	
	\begin{frame}{Q142: Ensuring Vulnerability Mitigation}
		\begin{block}{Question}
			Which BEST ensures appropriate mitigation occurs on identified IS vulnerabilities?
		\end{block}
		\begin{enumerate}[(A)]
			\item Presenting root cause analysis to management
			\item Implementing software to input action points
			\item Incorporating findings into annual report
			\item Assigning action plans with deadlines to responsible personnel
		\end{enumerate}
	\end{frame}
	
	\begin{frame}{Q142: Explanation}
		\begin{block}{Correct Answer: D}
			Assigning mitigation to personnel establishes responsibility for its completion within the deadline.
		\end{block}
		\begin{block}{Option Analysis}
			\begin{itemize}
				\item \textbf{A - Incorrect:} Management awareness doesn't ensure action
				\item \textbf{B - Incorrect:} Software helps monitor, doesn't ensure completion
				\item \textbf{C - Incorrect:} Reporting to shareholders doesn't ensure completion
				\item \textbf{D - Correct:} Assigning to personnel ensures accountability
			\end{itemize}
		\end{block}
	\end{frame}
	
	\begin{frame}{Q143: Audit Report Supplementary Controls}
		\begin{block}{Question}
			Which is BEST action if audit report recommends supplemental controls for uncritical functions?
		\end{block}
		\begin{enumerate}[(A)]
			\item Appoint project lead to start implementation
			\item Eliminate risk by purchasing risk insurance
			\item Request IT department implement all supplemental controls
			\item Request cost-benefit analysis to determine business need
		\end{enumerate}
	\end{frame}
	
	\begin{frame}{Q143: Explanation}
		\begin{block}{Correct Answer: D}
			A cost-benefit analysis may help determine if supplemental controls are needed, making this the best action to take. This is to ensure the cost of the control does not exceed the impact from the risk being realized.
		\end{block}
		\begin{block}{Option Analysis}
			\begin{itemize}
				\item \textbf{A - Incorrect:} Need must be justified first
				\item \textbf{B - Incorrect:} Insurance decision comes after cost-benefit
				\item \textbf{C - Incorrect:} Need must be justified first
				\item \textbf{D - Correct:} Cost-benefit determines business need
			\end{itemize}
		\end{block}
	\end{frame}
	
	\begin{frame}{Q144: Peer Review Enablement}
		\begin{block}{Question}
			Which BEST enables a peer review of an enterprise's risk management process?
		\end{block}
		\begin{enumerate}[(A)]
			\item Balanced scorecard
			\item Industry survey
			\item Capability maturity model
			\item Framework
		\end{enumerate}
	\end{frame}
	
	\begin{frame}{Q144: Explanation}
		\begin{block}{Correct Answer: C}
			A capability maturity model describes essential elements and criteria for effective processes for one or more disciplines. It also outlines an evolutionary improvement path from ad hoc, immature processes to disciplined, mature processes with improved quality and effectiveness.
		\end{block}
		\begin{block}{Option Analysis}
			\begin{itemize}
				\item \textbf{A - Incorrect:} Balanced scorecard is performance measurement
				\item \textbf{B - Incorrect:} Survey results are aggregated
				\item \textbf{C - Correct:} Maturity model enables peer comparison
				\item \textbf{D - Incorrect:} Framework defines concepts, not maturity
			\end{itemize}
		\end{block}
	\end{frame}
	
	\begin{frame}{Q145: Risk Treatment Plan Information}
		\begin{block}{Question}
			Which is MOST important information to include in risk treatment plan with resolution and date for completion?
		\end{block}
		\begin{enumerate}[(A)]
			\item Responsible personnel
			\item Mitigating factors
			\item Likelihood of occurrence
			\item Cost of completion
		\end{enumerate}
	\end{frame}
	
	\begin{frame}{Q145: Explanation}
		\begin{block}{Correct Answer: A}
			Risk response activities must be assigned to a responsible person or group; if this assignment is not included, it will be unclear who will implement the countermeasure.
		\end{block}
		\begin{block}{Option Analysis}
			\begin{itemize}
				\item \textbf{A - Correct:} Responsible personnel is essential
				\item \textbf{B - Incorrect:} Can be included but not as important
				\item \textbf{C - Incorrect:} Can be included but not as important
				\item \textbf{D - Incorrect:} Cost is optional
			\end{itemize}
		\end{block}
	\end{frame}
	
	\begin{frame}{Q146: Ensuring Acceptable Risk Level}
		\begin{block}{Question}
			Which BEST ensures identified risk remains at acceptable level?
		\end{block}
		\begin{enumerate}[(A)]
			\item Reviewing controls periodically, according to risk treatment plan
			\item Listing each risk as separate entry in risk register
			\item Creating separate risk register for every department
			\item Maintaining KRI for assets in risk register
		\end{enumerate}
	\end{frame}
	
	\begin{frame}{Q146: Explanation}
		\begin{block}{Correct Answer: A}
			Controls deployed according to the risk treatment plan should provide the desired results, because the risk treatment plan is based on management's acceptance of residual risk and management's approval of deployment steps in the plan.
		\end{block}
		\begin{block}{Option Analysis}
			\begin{itemize}
				\item \textbf{A - Correct:} Periodic review ensures continued effectiveness
				\item \textbf{B - Incorrect:} Register in itself doesn't ensure management
				\item \textbf{C - Incorrect:} Separate registers don't ensure management
				\item \textbf{D - Incorrect:} KRI in register doesn't ensure management
			\end{itemize}
		\end{block}
	\end{frame}
	
	\begin{frame}{Q147: Risk Response Increasing Liability}
		\begin{block}{Question}
			Which risk response option is MOST likely to increase liability of enterprise?
		\end{block}
		\begin{enumerate}[(A)]
			\item Risk acceptance
			\item Risk reduction
			\item Risk transfer
			\item Risk avoidance
		\end{enumerate}
	\end{frame}
	
	\begin{frame}{Q147: Explanation}
		\begin{block}{Correct Answer: A}
			An enterprise may choose to accept risk without knowing the correct level of risk that is being accepted; this may result in accusations of negligence.
		\end{block}
		\begin{block}{Option Analysis}
			\begin{itemize}
				\item \textbf{A - Correct:} Acceptance without proper assessment may be negligence
				\item \textbf{B - Incorrect:} Reduction attempts to reduce risk
				\item \textbf{C - Incorrect:} Transfer allocates risk to another party
				\item \textbf{D - Incorrect:} Avoidance terminates process
			\end{itemize}
		\end{block}
	\end{frame}
	
	\begin{frame}{Q148: Reason Not to Reduce Risk}
		\begin{block}{Question}
			Which is BEST reason for enterprise to decide not to reduce an identified risk?
		\end{block}
		\begin{enumerate}[(A)]
			\item No regulatory requirement
			\item There are mitigating controls in place
			\item Cost of mitigation exceeds the risk
			\item Budget for mitigation is limited
		\end{enumerate}
	\end{frame}
	
	\begin{frame}{Q148: Explanation}
		\begin{block}{Correct Answer: C}
			Enterprises will accept the risk when the cost of mitigation exceeds the risk.
		\end{block}
		\begin{block}{Option Analysis}
			\begin{itemize}
				\item \textbf{A - Incorrect:} Regulatory isn't the only factor
				\item \textbf{B - Incorrect:} Residual risk may still be above acceptable levels
				\item \textbf{C - Correct:} Cost-benefit justifies acceptance
				\item \textbf{D - Incorrect:} May reduce even when budget exceeded
			\end{itemize}
		\end{block}
	\end{frame}
	
	\begin{frame}{Q149: Outsourcing IT Activities}
		\begin{block}{Question}
			Enterprise addresses risk associated with IT project by outsourcing part of IT activities to third party with specialized skill set. This is:
		\end{block}
		\begin{enumerate}[(A)]
			\item Risk transfer
			\item Risk avoidance
			\item Risk acceptance
			\item Risk mitigation
		\end{enumerate}
	\end{frame}
	
	\begin{frame}{Q149: Explanation}
		\begin{block}{Correct Answer: D}
			When specific activities are outsourced to an entity with a specialized skill set, the inherent risk of the activity is reduced. The outsourcing process works as a control in this case.
		\end{block}
		\begin{block}{Option Analysis}
			\begin{itemize}
				\item \textbf{A - Incorrect:} Risk remains with enterprise
				\item \textbf{B - Incorrect:} Not avoidance
				\item \textbf{C - Incorrect:} Not acceptance
				\item \textbf{D - Correct:} Outsourcing reduces inherent risk
			\end{itemize}
		\end{block}
	\end{frame}
	
	\begin{frame}{Q150: Cost-Effective Risk Response}
		\begin{block}{Question}
			Which BEST facilitates cost-effective risk response?
		\end{block}
		\begin{enumerate}[(A)]
			\item Prioritizing and addressing risk according to risk management strategy
			\item Mitigating risk on basis of likelihood and impact
			\item Performing countermeasure analysis for each control
			\item Selecting controls that are at zero or near-zero costs
		\end{enumerate}
	\end{frame}
	
	\begin{frame}{Q150: Explanation}
		\begin{block}{Correct Answer: A}
			Prioritizing and addressing risk in line with the risk treatment strategy balances the costs and benefits of managing the IT risk.
		\end{block}
		\begin{block}{Option Analysis}
			\begin{itemize}
				\item \textbf{A - Correct:} Prioritization balances costs and benefits
				\item \textbf{B - Incorrect:} Multiple occurrences with similar products
				\item \textbf{C - Incorrect:} May not be a countermeasure for every control
				\item \textbf{D - Incorrect:} Zero-cost controls may not be effective
			\end{itemize}
		\end{block}
	\end{frame}
	
	% ================ SECTION 4: DOMAIN 4 - IT AND SECURITY ================
	
	\section{Domain 4: Information Technology and Security}
	
	\begin{frame}{Q151: Enterprise Architecture Purpose}
		\begin{block}{Question}
			PRIMARY reason for developing an enterprise security architecture is to:
		\end{block}
		\begin{enumerate}[(A)]
			\item Align security strategies among functional areas and external entities
			\item Build barrier between IT systems and outside world
			\item Foster understanding of technologies and their interactions
			\item Protect enterprise from external threats and monitor proactively
		\end{enumerate}
	\end{frame}
	
	\begin{frame}{Q151: Explanation}
		\begin{block}{Correct Answer: A}
			The enterprise security architecture must align strategies and objectives of diverse functional areas within the enterprise, optimize the flow of information within an enterprise, and support all required communication with external partners, customers and suppliers.
		\end{block}
		\begin{block}{Option Analysis}
			\begin{itemize}
				\item \textbf{A - Correct:} Alignment is primary purpose
				\item \textbf{B - Incorrect:} Building barrier may interfere with business
				\item \textbf{C - Incorrect:} Architecture documents interactions in relation to business
				\item \textbf{D - Incorrect:} Architecture establishes blueprint, doesn't directly protect
			\end{itemize}
		\end{block}
	\end{frame}
	
	\begin{frame}{Q152: Approving Risk Appetite}
		\begin{block}{Question}
			Who is responsible for approving enterprise's risk appetite and risk tolerance related to information security?
		\end{block}
		\begin{enumerate}[(A)]
			\item Business unit manager
			\item Information security officer
			\item Senior management
			\item Risk manager
		\end{enumerate}
	\end{frame}
	
	\begin{frame}{Q152: Explanation}
		\begin{block}{Correct Answer: C}
			Senior management determines enterprise risk appetite and risk tolerance.
		\end{block}
		\begin{block}{Option Analysis}
			\begin{itemize}
				\item \textbf{A - Incorrect:} BU manager doesn't determine appetite
				\item \textbf{B - Incorrect:} ISO doesn't determine appetite
				\item \textbf{C - Correct:} Senior management approves appetite/tolerance
				\item \textbf{D - Incorrect:} Risk manager doesn't determine appetite
			\end{itemize}
		\end{block}
	\end{frame}
	
	\begin{frame}{Q153: Influencing Risk Management Approach}
		\begin{block}{Question}
			Which would be MOST influential in determining an enterprise's approach to risk management?
		\end{block}
		\begin{enumerate}[(A)]
			\item Recent audit report findings
			\item Chief information security officer approval
			\item Key performance indicators
			\item Enterprise policies
		\end{enumerate}
	\end{frame}
	
	\begin{frame}{Q153: Explanation}
		\begin{block}{Correct Answer: D}
			Enterprise policies state the guidance for risk management.
		\end{block}
		\begin{block}{Option Analysis}
			\begin{itemize}
				\item \textbf{A - Incorrect:} Audit findings focus on specific areas
				\item \textbf{B - Incorrect:} Approach cannot be determined only from CISO
				\item \textbf{C - Incorrect:} KPIs monitor attainment, don't determine approach
				\item \textbf{D - Correct:} Enterprise policies guide risk management
			\end{itemize}
		\end{block}
	\end{frame}
	
	\begin{frame}{Q154: New Data Protection Regulation}
		\begin{block}{Question}
			What information should risk practitioner gather to BEST ensure compliance with new data protection regulation?
		\end{block}
		\begin{enumerate}[(A)]
			\item List of controls to implement
			\item Gaps associated with existing controls
			\item Risk scenarios with potential impact on compliance
			\item Enterprise's risk appetite
		\end{enumerate}
	\end{frame}
	
	\begin{frame}{Q154: Explanation}
		\begin{block}{Correct Answer: C}
			Risk scenarios should indicate potential effects of noncompliance with the new regulation and guide management in evaluating whether the cost of compliance outweighs the cost of noncompliance and aligns with the enterprise's risk tolerance.
		\end{block}
		\begin{block}{Option Analysis}
			\begin{itemize}
				\item \textbf{A - Incorrect:} Control list comes after business decision
				\item \textbf{B - Incorrect:} Control gaps useful after business decision
				\item \textbf{C - Correct:} Risk scenarios inform compliance decision
				\item \textbf{D - Incorrect:} Risk appetite suggests scenarios
			\end{itemize}
		\end{block}
	\end{frame}
	
	\begin{frame}{Q155: Supply Chain Risk}
		\begin{block}{Question}
			Which will BEST assist risk practitioner when addressing risk within supply chain life cycle?
		\end{block}
		\begin{enumerate}[(A)]
			\item Understanding supplier's organizational systems
			\item Understanding business case in support of life cycle
			\item Understanding relevant jurisdictional legal requirements
			\item Understanding enterprise's risk approach toward supply chain
		\end{enumerate}
	\end{frame}
	
	\begin{frame}{Q155: Explanation}
		\begin{block}{Correct Answer: C}
			Identification and understanding of the legal requirements relevant to the supply chain will assist the risk practitioner to identify, assess and monitor risk on an ongoing basis.
		\end{block}
		\begin{block}{Option Analysis}
			\begin{itemize}
				\item \textbf{A - Incorrect:} Important but not best
				\item \textbf{B - Incorrect:} Business case created prior to agreement
				\item \textbf{C - Correct:} Legal requirements are essential across regions
				\item \textbf{D - Incorrect:} Important but legal requirements are more important
			\end{itemize}
		\end{block}
	\end{frame}
	
	\begin{frame}{Q156: First Management Priority}
		\begin{block}{Question}
			To ensure risk program effectively mitigates risk, which should FIRST be addressed by management?
		\end{block}
		\begin{enumerate}[(A)]
			\item Risk findings
			\item Incomplete new hire checks
			\item Low rate of risk management training completion
			\item Blame culture
		\end{enumerate}
	\end{frame}
	
	\begin{frame}{Q156: Explanation}
		\begin{block}{Correct Answer: D}
			A blame culture should be addressed and remedied as soon as possible. It hinders effective risk mitigation because it defeats accountability.
		\end{block}
		\begin{block}{Option Analysis}
			\begin{itemize}
				\item \textbf{A - Incorrect:} Addressing findings is harder in blame culture
				\item \textbf{B - Incorrect:} Consideration but not as important
				\item \textbf{C - Incorrect:} Consideration but not as important
				\item \textbf{D - Correct:} Blame culture defeats accountability
			\end{itemize}
		\end{block}
	\end{frame}
	
	\begin{frame}{Q157: Key Element of Risk Culture}
		\begin{block}{Question}
			Which is a KEY element of risk culture?
		\end{block}
		\begin{enumerate}[(A)]
			\item Tolerance
			\item Policies
			\item Behavior
			\item Controls
		\end{enumerate}
	\end{frame}
	
	\begin{frame}{Q157: Explanation}
		\begin{block}{Correct Answer: C}
			Behavior toward taking risk, negative outcomes, and policy compliance are all elements of an enterprise's risk culture.
		\end{block}
		\begin{block}{Option Analysis}
			\begin{itemize}
				\item \textbf{A - Incorrect:} Tolerance is one element of behavior
				\item \textbf{B - Incorrect:} Policies are influenced by culture
				\item \textbf{C - Correct:} Behavior determines risk culture
				\item \textbf{D - Incorrect:} Controls are influenced by culture
			\end{itemize}
		\end{block}
	\end{frame}
	
	\begin{frame}{Q158: Accountable for Risk Governance Strategy}
		\begin{block}{Question}
			Who is accountable for overall enterprise strategy for risk governance?
		\end{block}
		\begin{enumerate}[(A)]
			\item Senior management
			\item Business unit management
			\item Chief risk officer
			\item Board of directors
		\end{enumerate}
	\end{frame}
	
	\begin{frame}{Q158: Explanation}
		\begin{block}{Correct Answer: D}
			The board of directors is accountable for the overall enterprise risk governance strategy as it states the enterprise strategy.
		\end{block}
		\begin{block}{Option Analysis}
			\begin{itemize}
				\item \textbf{A - Incorrect:} Strategy comes from board, not senior management
				\item \textbf{B - Incorrect:} Strategy comes from board
				\item \textbf{C - Incorrect:} Strategy comes from board
				\item \textbf{D - Correct:} Board is accountable for strategy
			\end{itemize}
		\end{block}
	\end{frame}
	
	\begin{frame}{Q159: Primary Goal of IT Risk Management}
		\begin{block}{Question}
			What is PRIMARY goal of enterprise's IT risk management process?
		\end{block}
		\begin{enumerate}[(A)]
			\item Protect IT assets and environment
			\item Protect enterprise and its ability to perform mission
			\item Report on status of IT risk register
			\item Ensure risk is identified and properly managed
		\end{enumerate}
	\end{frame}
	
	\begin{frame}{Q159: Explanation}
		\begin{block}{Correct Answer: B}
			The primary goal of an enterprise's risk management process is to protect the enterprise and its ability to perform its mission.
		\end{block}
		\begin{block}{Option Analysis}
			\begin{itemize}
				\item \textbf{A - Incorrect:} Should protect enterprise, not just IT assets
				\item \textbf{B - Correct:} Protect enterprise and mission
				\item \textbf{C - Incorrect:} Reporting is one tool, not the goal
				\item \textbf{D - Incorrect:} Identification is one component
			\end{itemize}
		\end{block}
	\end{frame}
	
	\begin{frame}{Q160: Risk Management Programs Goal}
		\begin{block}{Question}
			Risk management programs are designed to reduce risk to:
		\end{block}
		\begin{enumerate}[(A)]
			\item Point where benefit exceeds expense
			\item Level too small to be significant
			\item Rate of return equals cost of capital
			\item Level enterprise is willing to accept
		\end{enumerate}
	\end{frame}
	
	\begin{frame}{Q160: Explanation}
		\begin{block}{Correct Answer: D}
			Risk should be reduced to a level that an enterprise (i.e., management) is willing to accept; therefore, risk management primarily seeks to facilitate business judgment and support enterprise business goals.
		\end{block}
		\begin{block}{Option Analysis}
			\begin{itemize}
				\item \textbf{A - Incorrect:} Depends on risk appetite
				\item \textbf{B - Incorrect:} Not practical and cost prohibitive
				\item \textbf{C - Incorrect:} Not practical to calculate for each action
				\item \textbf{D - Correct:} Acceptable level determined by management
			\end{itemize}
		\end{block}
	\end{frame}
	
	\begin{frame}{Q161: Enterprise Security Policy Control Type}
		\begin{block}{Question}
			An enterprise security policy is an example of which control?
		\end{block}
		\begin{enumerate}[(A)]
			\item Operational control
			\item Management control
			\item Technical control
			\item Corrective control
		\end{enumerate}
	\end{frame}
	
	\begin{frame}{Q161: Explanation}
		\begin{block}{Correct Answer: B}
			There are two control methods: technical and nontechnical. Enterprise security policies are nontechnical management controls.
		\end{block}
		\begin{block}{Option Analysis}
			\begin{itemize}
				\item \textbf{A - Incorrect:} Manufacturing procedures are operational
				\item \textbf{B - Correct:} Policy is management control
				\item \textbf{C - Incorrect:} Encryption/IDS are technical
				\item \textbf{D - Incorrect:} Corrective controls respond to incidents
			\end{itemize}
		\end{block}
	\end{frame}
	
	\begin{frame}{Q162: Most Essential for Risk Program}
		\begin{block}{Question}
			Which is MOST essential for a risk management program to be effective?
		\end{block}
		\begin{enumerate}[(A)]
			\item New risk detection
			\item Sound risk baseline
			\item Accurate risk reporting
			\item Flexible security budget
		\end{enumerate}
	\end{frame}
	
	\begin{frame}{Q162: Explanation}
		\begin{block}{Correct Answer: A}
			Without identifying new risk, other measures will succeed only for a limited period.
		\end{block}
		\begin{block}{Option Analysis}
			\begin{itemize}
				\item \textbf{A - Correct:} New risk detection is essential
				\item \textbf{B - Incorrect:} Baseline is essential but detection is more essential
				\item \textbf{C - Incorrect:} Reporting is essential but detection is more essential
				\item \textbf{D - Incorrect:} Flexible budget is not available to most enterprises
			\end{itemize}
		\end{block}
	\end{frame}
	
	\begin{frame}{Q163: Most Important for Risk Mitigation}
		\begin{block}{Question}
			Which is MOST important when mitigating or managing risk?
		\end{block}
		\begin{enumerate}[(A)]
			\item Vulnerability assessment results
			\item Business impact analysis
			\item Risk appetite and tolerance levels
			\item Security controls framework
		\end{enumerate}
	\end{frame}
	
	\begin{frame}{Q163: Explanation}
		\begin{block}{Correct Answer: C}
			The risk tolerance level (along with risk appetite) determines what kind of risk response an enterprise selects, and it needs to be defined in order for an enterprise to appropriately address risk.
		\end{block}
		\begin{block}{Option Analysis}
			\begin{itemize}
				\item \textbf{A - Incorrect:} Provides view of control environment
				\item \textbf{B - Incorrect:} Concerned with prioritizing, not how to mitigate
				\item \textbf{C - Correct:} Tolerance/appetite determine response
				\item \textbf{D - Incorrect:} Framework describes what, not actionable recommendations
			\end{itemize}
		\end{block}
	\end{frame}
	
	\begin{frame}{Q164: Reporting IT Control Status}
		\begin{block}{Question}
			Which is MOST important component when reporting status of IT control environment to management?
		\end{block}
		\begin{enumerate}[(A)]
			\item Amendments to risk register
			\item Lessons learned from loss events
			\item Technical details on vulnerabilities
			\item Risk profile of enterprise
		\end{enumerate}
	\end{frame}
	
	\begin{frame}{Q164: Explanation}
		\begin{block}{Correct Answer: D}
			The risk profile of the enterprise is the most important component because it provides the overall portfolio of identified risk to which the enterprise is exposed. This is most important for management to know.
		\end{block}
		\begin{block}{Option Analysis}
			\begin{itemize}
				\item \textbf{A - Incorrect:} Amendments are important but not as important
				\item \textbf{B - Incorrect:} Lessons learned are not sufficient
				\item \textbf{C - Incorrect:} Technical details are not expected at management level
				\item \textbf{D - Correct:} Risk profile is most important
			\end{itemize}
		\end{block}
	\end{frame}
	
	\begin{frame}{Q165: Risk Scenarios in Assessments}
		\begin{block}{Question}
			Which uses risk scenarios when estimating likelihood and impact of significant risk?
		\end{block}
		\begin{enumerate}[(A)]
			\item IT audit
			\item Security gap analysis
			\item Threat and vulnerability assessment
			\item IT security assessment
		\end{enumerate}
	\end{frame}
	
	\begin{frame}{Q165: Explanation}
		\begin{block}{Correct Answer: C}
			A threat and vulnerability assessment typically evaluates all elements of a business process for threats and vulnerabilities and identifies the likelihood of occurrence and the business impact if the threats were realized.
		\end{block}
		\begin{block}{Option Analysis}
			\begin{itemize}
				\item \textbf{A - Incorrect:} IT audit uses technical evaluation
				\item \textbf{B - Incorrect:} Gap analysis doesn't use risk scenarios
				\item \textbf{C - Correct:} Threat/vulnerability assessment uses risk scenarios
				\item \textbf{D - Incorrect:} Security assessment uses technical evaluation
			\end{itemize}
		\end{block}
	\end{frame}
	
	\begin{frame}{Q166: Deviation from Policy}
		\begin{block}{Question}
			Enterprise wants to quickly implement technical solution deviating from policy. Risk practitioner should:
		\end{block}
		\begin{enumerate}[(A)]
			\item Recommend against implementation
			\item Recommend revision of current policy
			\item Conduct risk assessment and allow/disallow based on outcome
			\item Recommend risk assessment and implementation only if residual risk accepted
		\end{enumerate}
	\end{frame}
	
	\begin{frame}{Q166: Explanation}
		\begin{block}{Correct Answer: D}
			A risk assessment should be conducted to clarify the risk whenever the enterprise's policies cannot be followed. The solution should be implemented only if the related risk is formally accepted by the enterprise.
		\end{block}
		\begin{block}{Option Analysis}
			\begin{itemize}
				\item \textbf{A - Incorrect:} Business decisions driven by cost/benefit
				\item \textbf{B - Incorrect:} Should not be triggered by single request
				\item \textbf{C - Incorrect:} Management makes final decision
				\item \textbf{D - Correct:} Risk assessment and formal acceptance required
			\end{itemize}
		\end{block}
	\end{frame}
	
	\begin{frame}{Q167: Qualitative Risk Analysis}
		\begin{block}{Question}
			Which will produce comprehensive results when performing a qualitative risk analysis?
		\end{block}
		\begin{enumerate}[(A)]
			\item Vulnerability assessment
			\item Scenarios with threats and impacts
			\item Value of information assets
			\item Estimated productivity losses
		\end{enumerate}
	\end{frame}
	
	\begin{frame}{Q167: Explanation}
		\begin{block}{Correct Answer: B}
			Using a list of possible scenarios with threats and impacts will better frame the range of risk and facilitate a more informed discussion and decision.
		\end{block}
		\begin{block}{Option Analysis}
			\begin{itemize}
				\item \textbf{A - Incorrect:} One-sided view unless linked to scenarios
				\item \textbf{B - Correct:} Scenarios frame risk range
				\item \textbf{C - Incorrect:} Value without threats/impacts is insufficient
				\item \textbf{D - Incorrect:} Productivity losses are insufficient
			\end{itemize}
		\end{block}
	\end{frame}
	
	\begin{frame}{Q168: Removing Risk from Register}
		\begin{block}{Question}
			Risk may be removed from risk register when risk:
		\end{block}
		\begin{enumerate}[(A)]
			\item Has been eliminated
			\item Is transferred to vendor
			\item Threshold is exceeded
			\item Is no longer relevant
		\end{enumerate}
	\end{frame}
	
	\begin{frame}{Q168: Explanation}
		\begin{block}{Correct Answer: D}
			If a risk is relevant, it will be listed in the risk register. Once the risk is no longer relevant, it may be removed.
		\end{block}
		\begin{block}{Option Analysis}
			\begin{itemize}
				\item \textbf{A - Incorrect:} Risk cannot be eliminated
				\item \textbf{B - Incorrect:} Enterprise still accountable for risk
				\item \textbf{C - Incorrect:} If threshold exceeded, risk remains relevant
				\item \textbf{D - Correct:} No longer relevant means can be removed
			\end{itemize}
		\end{block}
	\end{frame}
	
	\begin{frame}{Q169: Accurate Risk Register Maintenance}
		\begin{block}{Question}
			Which is BEST way to ensure accurate risk register is maintained over time?
		\end{block}
		\begin{enumerate}[(A)]
			\item Monitor KRIs and record findings
			\item Publish risk register centrally with workflow features
			\item Distribute register to business process owners for review
			\item Use audit personnel to perform regular audits
		\end{enumerate}
	\end{frame}
	
	\begin{frame}{Q169: Explanation}
		\begin{block}{Correct Answer: B}
			Centrally publishing the risk register and enabling periodic polling of risk assessors through workflow features will ensure accuracy of content. A knowledge management platform with workflow and polling features will automate the process of maintaining the risk register.
		\end{block}
		\begin{block}{Option Analysis}
			\begin{itemize}
				\item \textbf{A - Incorrect:} Only provides insights to known risk
				\item \textbf{B - Correct:} Workflow automation ensures accuracy
				\item \textbf{C - Incorrect:} Business owners may not be unbiased
				\item \textbf{D - Incorrect:} Audit personnel may not have business knowledge
			\end{itemize}
		\end{block}
	\end{frame}
	
	\begin{frame}{Q170: Adding Emerging Risk}
		\begin{block}{Question}
			An emerging risk should be added to risk register by risk practitioner when:
		\end{block}
		\begin{enumerate}[(A)]
			\item Impact can be quantified
			\item Probability of occurrence is high
			\item Competitor has added risk
			\item Activity that triggers risk initiates
		\end{enumerate}
	\end{frame}
	
	\begin{frame}{Q170: Explanation}
		\begin{block}{Correct Answer: D}
			Risk identification usually starts when planning an activity, and risk identified at planning needs to be added to the risk register to ensure effective risk management.
		\end{block}
		\begin{block}{Option Analysis}
			\begin{itemize}
				\item \textbf{A - Incorrect:} Quantification is important but not decision factor
				\item \textbf{B - Incorrect:} High impact alone not the only criterion
				\item \textbf{C - Incorrect:} Industry benchmark shouldn't be used
				\item \textbf{D - Correct:} Risk identified at planning should be added
			\end{itemize}
		\end{block}
	\end{frame}
	
	\begin{frame}{Q171: Primary Objective of Risk Management}
		\begin{block}{Question}
			Which is PRIMARY objective of a risk management program?
		\end{block}
		\begin{enumerate}[(A)]
			\item Maintain residual risk at acceptable level
			\item Implement preventive controls for every threat
			\item Remove all identified risk
			\item Reduce inherent risk to zero
		\end{enumerate}
	\end{frame}
	
	\begin{frame}{Q171: Explanation}
		\begin{block}{Correct Answer: A}
			Ensuring that all residual risk is maintained at a level acceptable to the business is the objective of a risk management program.
		\end{block}
		\begin{block}{Option Analysis}
			\begin{itemize}
				\item \textbf{A - Correct:} Maintain residual risk at acceptable level
				\item \textbf{B - Incorrect:} Not the objective of risk management
				\item \textbf{C - Incorrect:} Not intended to remove every identified risk
				\item \textbf{D - Incorrect:} Inherent risk is always greater than zero
			\end{itemize}
		\end{block}
	\end{frame}
	
	\begin{frame}{Q172: Risk Assessment Techniques Use}
		\begin{block}{Question}
			Risk assessment techniques should be used by risk practitioner to:
		\end{block}
		\begin{enumerate}[(A)]
			\item Maximize return on investment
			\item Provide documentation for auditors
			\item Justify selection of risk mitigation strategies
			\item Quantify risk that would otherwise be subjective
		\end{enumerate}
	\end{frame}
	
	\begin{frame}{Q172: Explanation}
		\begin{block}{Correct Answer: C}
			A risk practitioner should use risk assessment techniques to justify and implement a risk mitigation strategy as efficiently as possible.
		\end{block}
		\begin{block}{Option Analysis}
			\begin{itemize}
				\item \textbf{A - Incorrect:} Maximizing ROI is not part of risk assessment
				\item \textbf{B - Incorrect:} Doesn't focus on auditors as primary recipients
				\item \textbf{C - Correct:} Justifies risk mitigation strategies
				\item \textbf{D - Incorrect:} Risk assessment is high-level; analysis can be quantitative or qualitative
			\end{itemize}
		\end{block}
	\end{frame}
	
	\begin{frame}{Q173: New Printer Risk}
		\begin{block}{Question}
			Procurement employee notices new printer models keep copy of all printed documents on internal hard disk. Employee should:
		\end{block}
		\begin{enumerate}[(A)]
			\item Proceed with order and configure printers to auto-wipe data
			\item Notify security manager to conduct risk assessment
			\item Seek vendor without built-in hard disks
			\item Procure with hard disks and notify staff to wipe when decommissioning
		\end{enumerate}
	\end{frame}
	
	\begin{frame}{Q173: Explanation}
		\begin{block}{Correct Answer: B}
			Risk assessment is most appropriate because it yields risk mitigation techniques that are appropriate for enterprise risk context and appetite.
		\end{block}
		\begin{block}{Option Analysis}
			\begin{itemize}
				\item \textbf{A - Incorrect:} Not appropriate without prior risk assessment
				\item \textbf{B - Correct:} Risk assessment identifies appropriate controls
				\item \textbf{C - Incorrect:} Focusing solely on risk ignores opportunity
				\item \textbf{D - Incorrect:} Notifying staff is not sufficient control
			\end{itemize}
		\end{block}
	\end{frame}
	
	\begin{frame}{Q174: Risk Assessment Recurrence}
		\begin{block}{Question}
			Risk assessments should be repeated at regular intervals because:
		\end{block}
		\begin{enumerate}[(A)]
			\item Omissions in earlier assessments can be addressed
			\item Periodic assessments allow various methodologies
			\item Business threats are constantly changing
			\item They help raise risk awareness among staff
		\end{enumerate}
	\end{frame}
	
	\begin{frame}{Q174: Explanation}
		\begin{block}{Correct Answer: C}
			As business objectives and methods change, the nature and relevance of threats also change.
		\end{block}
		\begin{block}{Option Analysis}
			\begin{itemize}
				\item \textbf{A - Incorrect:} Omissions don't justify regular reassessments
				\item \textbf{B - Incorrect:} Same methodologies should be used
				\item \textbf{C - Correct:} Threats constantly change
				\item \textbf{D - Incorrect:} Better ways to raise awareness exist
			\end{itemize}
		\end{block}
	\end{frame}
	
	\begin{frame}{Q175: Assessing IS Risk}
		\begin{block}{Question}
			Assessing information systems risk is BEST achieved by:
		\end{block}
		\begin{enumerate}[(A)]
			\item Using enterprise's past actual loss experience
			\item Reviewing published loss statistics
			\item Evaluating threats associated with existing IS assets
			\item Reviewing IS control weaknesses in audit reports
		\end{enumerate}
	\end{frame}
	
	\begin{frame}{Q175: Explanation}
		\begin{block}{Correct Answer: C}
			To assess IT risk, threats and vulnerabilities need to be evaluated using qualitative or quantitative risk assessment approaches.
		\end{block}
		\begin{block}{Option Analysis}
			\begin{itemize}
				\item \textbf{A - Incorrect:} Doesn't address realistic risk scenarios
				\item \textbf{B - Incorrect:} Doesn't address enterprise-specific scenarios
				\item \textbf{C - Correct:} Evaluating threats is fundamental
				\item \textbf{D - Incorrect:} Control weaknesses alone are not useful
			\end{itemize}
		\end{block}
	\end{frame}
	
	\begin{frame}{Q176: Lack of Adequate Controls}
		\begin{block}{Question}
			A lack of adequate controls represents:
		\end{block}
		\begin{enumerate}[(A)]
			\item A vulnerability
			\item An impact
			\item An asset
			\item A threat
		\end{enumerate}
	\end{frame}
	
	\begin{frame}{Q176: Explanation}
		\begin{block}{Correct Answer: A}
			Lack of adequate controls represents a vulnerability, exposing sensitive information and data to the risk of malicious damage, attack, or unauthorized access.
		\end{block}
		\begin{block}{Option Analysis}
			\begin{itemize}
				\item \textbf{A - Correct:} Vulnerability = lack of controls
				\item \textbf{B - Incorrect:} Impact is financial loss measure
				\item \textbf{C - Incorrect:} Assets have tangible/intangible value
				\item \textbf{D - Incorrect:} Threat is potential cause
			\end{itemize}
		\end{block}
	\end{frame}
	
	\begin{frame}{Q177: Security Breach at Another Entity}
		\begin{block}{Question}
			Enterprise learns of security breach at another entity using similar network technology. MOST important action is:
		\end{block}
		\begin{enumerate}[(A)]
			\item Assess likelihood of incident occurring
			\item Discontinue use of vulnerable technology
			\item Report to senior management that enterprise is not affected
			\item Remind staff no similar breaches have taken place
		\end{enumerate}
	\end{frame}
	
	\begin{frame}{Q177: Explanation}
		\begin{block}{Correct Answer: A}
			The likelihood of a similar incident occurring at the risk practitioner's enterprise should be assessed first, based on available information.
		\end{block}
		\begin{block}{Option Analysis}
			\begin{itemize}
				\item \textbf{A - Correct:} Assess likelihood first
				\item \textbf{B - Incorrect:} Discontinuation may not be required
				\item \textbf{C - Incorrect:} Premature report
				\item \textbf{D - Incorrect:} Cannot be certain
			\end{itemize}
		\end{block}
	\end{frame}
	
	\begin{frame}{Q178: Comprehensive Risk Assessment Trigger}
		\begin{block}{Question}
			MOST likely trigger for conducting comprehensive risk assessment is detection of changes in:
		\end{block}
		\begin{enumerate}[(A)]
			\item Asset inventory
			\item Asset classification levels
			\item Business environment
			\item Information security policies
		\end{enumerate}
	\end{frame}
	
	\begin{frame}{Q178: Explanation}
		\begin{block}{Correct Answer: C}
			Changes in the business environment, including new threats, vulnerabilities, or changes to information asset deployment, will trigger comprehensive periodic risk assessment.
		\end{block}
		\begin{block}{Option Analysis}
			\begin{itemize}
				\item \textbf{A - Incorrect:} Ongoing process, doesn't trigger
				\item \textbf{B - Incorrect:} Classification changes don't trigger
				\item \textbf{C - Correct:} Business environment changes trigger assessment
				\item \textbf{D - Incorrect:} Policy changes don't trigger assessment
			\end{itemize}
		\end{block}
	\end{frame}
	
	\begin{frame}{Q179: Anonymous Risk Identification}
		\begin{block}{Question}
			Which is BEST risk identification technique to support enterprise that allows employees to identify risk anonymously?
		\end{block}
		\begin{enumerate}[(A)]
			\item Delphi technique
			\item Isolated pilot groups
			\item SWOT analysis
			\item Root cause analysis
		\end{enumerate}
	\end{frame}
	
	\begin{frame}{Q179: Explanation}
		\begin{block}{Correct Answer: A}
			With the Delphi technique, polling or information gathering is done either anonymously or privately between the interviewer and interviewee.
		\end{block}
		\begin{block}{Option Analysis}
			\begin{itemize}
				\item \textbf{A - Correct:} Delphi uses anonymous polling
				\item \textbf{B - Incorrect:} Participants don't identify anonymously
				\item \textbf{C - Incorrect:} Participants don't identify anonymously
				\item \textbf{D - Incorrect:} Participants don't identify anonymously
			\end{itemize}
		\end{block}
	\end{frame}
	
	\begin{frame}{Q180: Risk Register Information}
		\begin{block}{Question}
			Which BEST describes information needed for each risk in a risk register?
		\end{block}
		\begin{enumerate}[(A)]
			\item Date, description, impact, probability, risk score, mitigation action, owner
			\item Date, description, risk score, cost to remediate, communication plan, owner
			\item Date, description, impact, cost to remediate, owner
			\item Various activities leading to risk management planning
		\end{enumerate}
	\end{frame}
	
	\begin{frame}{Q180: Explanation}
		\begin{block}{Correct Answer: A}
			Information required for each risk in a risk register includes date, description, impact, probability, risk score, mitigation action and owner.
		\end{block}
		\begin{block}{Option Analysis}
			\begin{itemize}
				\item \textbf{A - Correct:} All required elements included
				\item \textbf{B - Incorrect:} Missing impact, probability, mitigation action
				\item \textbf{C - Incorrect:} Missing probability, risk score, mitigation action
				\item \textbf{D - Incorrect:} Risk register results from planning, not vice versa
			\end{itemize}
		\end{block}
	\end{frame}
	
	\begin{frame}{Q181: Business Impact Analysis Advantage}
		\begin{block}{Question}
			GREATEST advantage in performing a business impact analysis is that it:
		\end{block}
		\begin{enumerate}[(A)]
			\item Does not have to be updated
			\item Promotes continuity awareness in enterprise
			\item Requires only qualitative estimates
			\item Eliminates need for risk analysis
		\end{enumerate}
	\end{frame}
	
	\begin{frame}{Q181: Explanation}
		\begin{block}{Correct Answer: B}
			A BIA raises enterprise-wide awareness of risk to business recovery and continuity.
		\end{block}
		\begin{block}{Option Analysis}
			\begin{itemize}
				\item \textbf{A - Incorrect:} Should be updated periodically
				\item \textbf{B - Correct:} Promotes continuity awareness
				\item \textbf{C - Incorrect:} Should use both qualitative and quantitative
				\item \textbf{D - Incorrect:} Cannot eliminate risk analysis
			\end{itemize}
		\end{block}
	\end{frame}
	
	\begin{frame}{Q182: Risk Register Primary Advantage}
		\begin{block}{Question}
			PRIMARY advantage of creating and updating a risk register is to:
		\end{block}
		\begin{enumerate}[(A)]
			\item Ensure inventory of identified risk is maintained
			\item Record all risk scenarios considered
			\item Collect similar data on all risk
			\item Run reports based on various risk scenarios
		\end{enumerate}
	\end{frame}
	
	\begin{frame}{Q182: Explanation}
		\begin{block}{Correct Answer: A}
			Once assets and risk are identified, the risk register is used as an inventory of that risk. The risk register can accelerate risk decision-making and establish accountability for specific risk.
		\end{block}
		\begin{block}{Option Analysis}
			\begin{itemize}
				\item \textbf{A - Correct:} Maintaining inventory is primary advantage
				\item \textbf{B - Incorrect:} Recording is good practice but inventory is primary
				\item \textbf{C - Incorrect:} Data collection is benefit but inventory is primary
				\item \textbf{D - Incorrect:} Running reports is a benefit but not primary
			\end{itemize}
		\end{block}
	\end{frame}
	
	\begin{frame}{Q183: Measuring Risk Management Development}
		\begin{block}{Question}
			Which BEST assists in measuring existing level of development of risk management processes against desired state?
		\end{block}
		\begin{enumerate}[(A)]
			\item Capability maturity model
			\item Risk management audit reports
			\item Balanced scorecard
			\item Enterprise security architecture
		\end{enumerate}
	\end{frame}
	
	\begin{frame}{Q183: Explanation}
		\begin{block}{Correct Answer: A}
			The capability maturity model grades processes on a scale of 0 to 5, based on their maturity, and is commonly used by entities to measure their existing state and then to determine the desired one.
		\end{block}
		\begin{block}{Option Analysis}
			\begin{itemize}
				\item \textbf{A - Correct:} Maturity model measures current vs desired
				\item \textbf{B - Incorrect:} Offers limited view
				\item \textbf{C - Incorrect:} Measures implementation of strategy
				\item \textbf{D - Incorrect:} Explains security architecture
			\end{itemize}
		\end{block}
	\end{frame}
	
	\begin{frame}{Q184: Assessing Business Risk}
		\begin{block}{Question}
			Which is MOST effective in assessing business risk?
		\end{block}
		\begin{enumerate}[(A)]
			\item Use case analysis
			\item Business case analysis
			\item Risk scenarios
			\item Risk plan
		\end{enumerate}
	\end{frame}
	
	\begin{frame}{Q184: Explanation}
		\begin{block}{Correct Answer: C}
			Risk scenarios are the most effective technique in assessing business risk. Scenarios help determine the likelihood and impact of an identified risk.
		\end{block}
		\begin{block}{Option Analysis}
			\begin{itemize}
				\item \textbf{A - Incorrect:} Identifies business requirements
				\item \textbf{B - Incorrect:} Part of project charter
				\item \textbf{C - Correct:} Scenarios determine likelihood and impact
				\item \textbf{D - Incorrect:} Output of risk assessment
			\end{itemize}
		\end{block}
	\end{frame}
	
	\begin{frame}{Q185: Risk Register Preparation}
		\begin{block}{Question}
			Preparation of a risk register begins in which risk management process?
		\end{block}
		\begin{enumerate}[(A)]
			\item Risk response planning
			\item Risk monitoring and control
			\item Risk management planning
			\item Risk identification
		\end{enumerate}
	\end{frame}
	
	\begin{frame}{Q185: Explanation}
		\begin{block}{Correct Answer: D}
			The risk register details all identified risk, including description, category, cause, probability of occurring, impact on objectives, proposed responses, owners and current status. The primary outputs of risk identification are the initial entries into the risk register.
		\end{block}
		\begin{block}{Option Analysis}
			\begin{itemize}
				\item \textbf{A - Incorrect:} Responses determined and included in register
				\item \textbf{B - Incorrect:} New risk identified triggers updates
				\item \textbf{C - Incorrect:} Describes how risk management will be performed
				\item \textbf{D - Correct:} Risk identification creates initial entries
			\end{itemize}
		\end{block}
	\end{frame}
	
	\begin{frame}{Q186: BIA Primary Use}
		\begin{block}{Question}
			A business impact analysis is PRIMARILY used to:
		\end{block}
		\begin{enumerate}[(A)]
			\item Estimate resources to resume normal operations
			\item Evaluate impact of disruption on ability to operate over time
			\item Calculate likelihood and impact of known threats
			\item Evaluate high-level business requirements
		\end{enumerate}
	\end{frame}
	
	\begin{frame}{Q186: Explanation}
		\begin{block}{Correct Answer: B}
			A BIA is primarily intended to evaluate the impact of disruption over time to an enterprise's ability to operate. It determines the urgency of each business activity. Key deliverables include recovery time objectives and recovery point objectives.
		\end{block}
		\begin{block}{Option Analysis}
			\begin{itemize}
				\item \textbf{A - Incorrect:} This is business continuity planning
				\item \textbf{B - Correct:} Evaluates impact over time
				\item \textbf{C - Incorrect:} Calculated during risk analysis
				\item \textbf{D - Incorrect:} Defined during early SDLC phases
			\end{itemize}
		\end{block}
	\end{frame}
	
	\begin{frame}{Q187: Risk Scenarios Enable Assessment}
		\begin{block}{Question}
			Risk scenarios enable the risk assessment process because they:
		\end{block}
		\begin{enumerate}[(A)]
			\item Cover a wide range of potential risk
			\item Minimize need for quantitative risk analysis
			\item Segregate IT risk from business risk
			\item Help estimate frequency and impact of risk
		\end{enumerate}
	\end{frame}
	
	\begin{frame}{Q187: Explanation}
		\begin{block}{Correct Answer: D}
			While risk scenarios may address a wide range of risk, risk scenarios help to estimate the frequency and impact of risk—two key elements of the risk assessment process. These elements aid subsequent steps in risk management by making risk relevant to business process owners.
		\end{block}
		\begin{block}{Option Analysis}
			\begin{itemize}
				\item \textbf{A - Incorrect:} Not always the result
				\item \textbf{B - Incorrect:} Don't necessarily minimize quantitative analysis
				\item \textbf{C - Incorrect:} Can be applied to both
				\item \textbf{D - Correct:} Help estimate frequency and impact
			\end{itemize}
		\end{block}
	\end{frame}
	
	\begin{frame}{Q188: Best Help for Risk Scenarios}
		\begin{block}{Question}
			Which information in risk register BEST helps in developing proper risk scenarios?
		\end{block}
		\begin{enumerate}[(A)]
			\item List of potential threats to assets
			\item List of residual risk on individual assets
			\item List of accepted risk
			\item List of security incidents
		\end{enumerate}
	\end{frame}
	
	\begin{frame}{Q188: Explanation}
		\begin{block}{Correct Answer: A}
			Identifying potential threats to business assets will help isolate vulnerabilities and associated risk, all of which contribute to developing proper risk scenarios.
		\end{block}
		\begin{block}{Option Analysis}
			\begin{itemize}
				\item \textbf{A - Correct:} Potential threats help develop scenarios
				\item \textbf{B - Incorrect:} Doesn't help develop scenarios
				\item \textbf{C - Incorrect:} Generally small subset
				\item \textbf{D - Incorrect:} Previous incidents may inspire scenarios but not best approach
			\end{itemize}
		\end{block}
	\end{frame}
	
	\begin{frame}{Q189: Identifying Control Deficiencies}
		\begin{block}{Question}
			Which BEST helps identify information systems control deficiencies?
		\end{block}
		\begin{enumerate}[(A)]
			\item Gap analysis
			\item Current IT risk profile
			\item IT controls framework
			\item Countermeasure analysis
		\end{enumerate}
	\end{frame}
	
	\begin{frame}{Q189: Explanation}
		\begin{block}{Correct Answer: A}
			Controls are deployed to achieve control objectives based on risk assessments and business requirements. The gap between desired control objectives and actual control design and operational effectiveness identifies control deficiencies in information systems.
		\end{block}
		\begin{block}{Option Analysis}
			\begin{itemize}
				\item \textbf{A - Correct:} Gap analysis identifies deficiencies
				\item \textbf{B - Incorrect:} Doesn't expose gap to desired state
				\item \textbf{C - Incorrect:} Generic document with no current state info
				\item \textbf{D - Incorrect:} Only helps identify countermeasure deficiencies
			\end{itemize}
		\end{block}
	\end{frame}
	
	\begin{frame}{Q190: Assessing Application Server Performance}
		\begin{block}{Question}
			When assessing performance of critical application server, MOST reliable results may be obtained from:
		\end{block}
		\begin{enumerate}[(A)]
			\item Activation of native database auditing
			\item Documentation of performance objectives
			\item Continuous monitoring
			\item Documentation of security modules
		\end{enumerate}
	\end{frame}
	
	\begin{frame}{Q190: Explanation}
		\begin{block}{Correct Answer: C}
			With continuous monitoring it is possible to track key performance metrics and also possible to address critical issues with minimum impact.
		\end{block}
		\begin{block}{Option Analysis}
			\begin{itemize}
				\item \textbf{A - Incorrect:} Database logs don't provide performance info
				\item \textbf{B - Incorrect:} Documentation is important but doesn't provide performance info
				\item \textbf{C - Correct:} Continuous monitoring tracks performance
				\item \textbf{D - Incorrect:} Security modules don't provide performance info
			\end{itemize}
		\end{block}
	\end{frame}
	
	\begin{frame}{Q191: Greatest Risk When Updating Register}
		\begin{block}{Question}
			Which presents GREATEST risk when updating the risk register?
		\end{block}
		\begin{enumerate}[(A)]
			\item Updates carried out jointly with other functions
			\item Updates carried out following incidents
			\item Updates carried out annually
			\item Updates subject to CISO approval
		\end{enumerate}
	\end{frame}
	
	\begin{frame}{Q191: Explanation}
		\begin{block}{Correct Answer: C}
			Updating the risk register only annually means the risk register does not reflect the true status of IT risk in the enterprise.
		\end{block}
		\begin{block}{Option Analysis}
			\begin{itemize}
				\item \textbf{A - Incorrect:} May be managed by multiple functions
				\item \textbf{B - Incorrect:} This presents a risk but not greatest
				\item \textbf{C - Correct:} Annual updates don't reflect true status
				\item \textbf{D - Incorrect:} Problematic but not as great as annual updates
			\end{itemize}
		\end{block}
	\end{frame}
	
	\begin{frame}{Q192: IT Risk Measurement}
		\begin{block}{Question}
			IT risk is measured by its:
		\end{block}
		\begin{enumerate}[(A)]
			\item Level of damage to IT systems
			\item Impact on business operations
			\item Cost of countermeasures
			\item Annual loss expectancy
		\end{enumerate}
	\end{frame}
	
	\begin{frame}{Q192: Explanation}
		\begin{block}{Correct Answer: B}
			IT risk includes information and communication technology risk but is primarily measured by its impact on the business. IT risk is the business risk associated with the use, ownership, operation, involvement, influence and adoption of IT within an enterprise.
		\end{block}
		\begin{block}{Option Analysis}
			\begin{itemize}
				\item \textbf{A - Incorrect:} IT damage alone is not comprehensive
				\item \textbf{B - Correct:} Impact on business is primary
				\item \textbf{C - Incorrect:} Cost/benefit is risk response, not assessment
				\item \textbf{D - Incorrect:} ALE must be used with qualitative measures
			\end{itemize}
		\end{block}
	\end{frame}
	
	\begin{frame}{Q193: Quantitative Risk Analysis}
		\begin{block}{Question}
			Deriving likelihood and impact of risk scenarios through statistical methods is BEST described as:
		\end{block}
		\begin{enumerate}[(A)]
			\item Quantitative risk analysis
			\item Risk scenario analysis
			\item Qualitative risk analysis
			\item Probabilistic risk assessment
		\end{enumerate}
	\end{frame}
	
	\begin{frame}{Q193: Explanation}
		\begin{block}{Correct Answer: A}
			Quantitative risk analysis derives the probability and impact of risk scenarios from statistical methods and data.
		\end{block}
		\begin{block}{Option Analysis}
			\begin{itemize}
				\item \textbf{A - Correct:} Uses statistical methods
				\item \textbf{B - Incorrect:} Includes several methods
				\item \textbf{C - Incorrect:} Uses non-quantitative measures
				\item \textbf{D - Incorrect:} Applied to complex engineered technology
			\end{itemize}
		\end{block}
	\end{frame}
	
	\begin{frame}{Q194: Local Mitigation Responsibility}
		\begin{block}{Question}
			During internal risk assessment, local management proactively mitigated high-level global purchasing process risk. This means:
		\end{block}
		\begin{enumerate}[(A)]
			\item Local management is now responsible for risk
			\item Risk owner is corporate chief risk officer
			\item Risk owner is local purchasing manager
			\item Corporate management remains responsible for risk
		\end{enumerate}
	\end{frame}
	
	\begin{frame}{Q194: Explanation}
		\begin{block}{Correct Answer: D}
			Corporate management remains responsible for the risk, even when the risk response is executed at a lower organizational level.
		\end{block}
		\begin{block}{Option Analysis}
			\begin{itemize}
				\item \textbf{A - Incorrect:} Corporate management remains responsible
				\item \textbf{B - Incorrect:} CRO is responsible for program, not purchasing risk
				\item \textbf{C - Incorrect:} Risk owner is global purchasing manager
				\item \textbf{D - Correct:} Corporate management remains responsible
			\end{itemize}
		\end{block}
	\end{frame}
	
	\begin{frame}{Q195: Estimating Likelihood of Events}
		\begin{block}{Question}
			Which BEST estimates likelihood of significant events affecting enterprise?
		\end{block}
		\begin{enumerate}[(A)]
			\item Threat analysis
			\item Cost-benefit analysis
			\item Scenario analysis
			\item Countermeasure analysis
		\end{enumerate}
	\end{frame}
	
	\begin{frame}{Q195: Explanation}
		\begin{block}{Correct Answer: C}
			Scenario analysis, along with vulnerability analysis, best determines whether a particular risk is relevant to the enterprise, and helps estimate the likelihood that significant events will affect the enterprise.
		\end{block}
		\begin{block}{Option Analysis}
			\begin{itemize}
				\item \textbf{A - Incorrect:} Doesn't provide sufficient info
				\item \textbf{B - Incorrect:} Used in selecting controls
				\item \textbf{C - Correct:} Scenario analysis estimates likelihood
				\item \textbf{D - Incorrect:} Assesses controls during attacks
			\end{itemize}
		\end{block}
	\end{frame}
	
	\begin{frame}{Q196: Improving Risk Decision-Making}
		\begin{block}{Question}
			Which BEST improves decision-making related to risk?
		\end{block}
		\begin{enumerate}[(A)]
			\item Maintaining documented risk register of all possible risk
			\item Risk awareness training in line with risk culture
			\item Maintaining updated security policies and procedures
			\item Allocating accountability of risk to department as a whole
		\end{enumerate}
	\end{frame}
	
	\begin{frame}{Q196: Explanation}
		\begin{block}{Correct Answer: A}
			Maintaining a documented risk register improves decision-making related to risk response because a risk register captures the population of relevant risk scenarios and provides a basis for prioritization of risk responses.
		\end{block}
		\begin{block}{Option Analysis}
			\begin{itemize}
				\item \textbf{A - Correct:} Register captures and prioritizes risk
				\item \textbf{B - Incorrect:} Training enhances accountability but less useful for emerging threats
				\item \textbf{C - Incorrect:} Policies don't necessarily improve decisions
				\item \textbf{D - Incorrect:} Department accountability dilutes ownership
			\end{itemize}
		\end{block}
	\end{frame}
	
	\begin{frame}{Q197: Independent Risk Management Review}
		\begin{block}{Question}
			PRIMARY reason to have risk management process reviewed by independent risk management professional is to:
		\end{block}
		\begin{enumerate}[(A)]
			\item Validate cost-effective solutions
			\item Validate control weaknesses
			\item Assess validity of end-to-end process
			\item Assess whether risk profile and factors are properly defined
		\end{enumerate}
	\end{frame}
	
	\begin{frame}{Q197: Explanation}
		\begin{block}{Correct Answer: C}
			The independent risk professional will be unbiased to review the risk management process end to end. The independent reviewer will not have any involvement in any stage of the risk management process and will be unaffected by all internal factors.
		\end{block}
		\begin{block}{Option Analysis}
			\begin{itemize}
				\item \textbf{A - Incorrect:} Internal teams can provide cost-effective solutions
				\item \textbf{B - Incorrect:} Internal team can find weaknesses
				\item \textbf{C - Correct:} Independent review of end-to-end process
				\item \textbf{D - Incorrect:} Risk profile is properly defined when assessment is correct
			\end{itemize}
		\end{block}
	\end{frame}
	
	\begin{frame}{Q198: Reviewing IT Risk Analysis Results}
		\begin{block}{Question}
			Which is BEST suited for review of IT risk analysis results before sending to management?
		\end{block}
		\begin{enumerate}[(A)]
			\item Internal audit review
			\item Peer review
			\item Compliance review
			\item Risk policy review
		\end{enumerate}
	\end{frame}
	
	\begin{frame}{Q198: Explanation}
		\begin{block}{Correct Answer: B}
			It is effective, efficient and good practice to perform a peer review of IT risk analysis results before sending them to management.
		\end{block}
		\begin{block}{Option Analysis}
			\begin{itemize}
				\item \textbf{A - Incorrect:} Internal audit is independent assurance
				\item \textbf{B - Correct:} Peer review is good practice
				\item \textbf{C - Incorrect:} Compliance reviews measure conformance
				\item \textbf{D - Incorrect:} Policy review may change content but isn't review method
			\end{itemize}
		\end{block}
	\end{frame}
	
	\begin{frame}{Q199: First Step in Identifying IT Risk}
		\begin{block}{Question}
			FIRST step in identifying and assessing IT risk is to:
		\end{block}
		\begin{enumerate}[(A)]
			\item Confirm risk tolerance level
			\item Identify threats and vulnerabilities
			\item Gather information on current and future environment
			\item Review past incident reports
		\end{enumerate}
	\end{frame}
	
	\begin{frame}{Q199: Explanation}
		\begin{block}{Correct Answer: C}
			The first step in any risk assessment is to gather information about the current state and pending internal and external changes to the enterprise's environment (scope, technology, incidents, modifications, etc.).
		\end{block}
		\begin{block}{Option Analysis}
			\begin{itemize}
				\item \textbf{A - Incorrect:} Tolerance informs response, not identification
				\item \textbf{B - Incorrect:} Must be supplemented by environmental changes
				\item \textbf{C - Correct:} Gather environmental information first
				\item \textbf{D - Incorrect:} Past incidents are one input but not adequate
			\end{itemize}
		\end{block}
	\end{frame}
	
	\begin{frame}{Q200: Risk Scenarios Creation Basis}
		\begin{block}{Question}
			Risk scenarios should be created PRIMARILY based on:
		\end{block}
		\begin{enumerate}[(A)]
			\item Input from senior management
			\item Previous security incidents
			\item Threats that the enterprise faces
			\item Results of risk analysis
		\end{enumerate}
	\end{frame}
	
	\begin{frame}{Q200: Explanation}
		\begin{block}{Correct Answer: C}
			When creating risk scenarios, the most important factor to consider is the likelihood of threats or threat actions occurring due to the risk.
		\end{block}
		\begin{block}{Option Analysis}
			\begin{itemize}
				\item \textbf{A - Incorrect:} Not as critical as threats
				\item \textbf{B - Incorrect:} Not as critical as threats
				\item \textbf{C - Correct:} Threats are most important
				\item \textbf{D - Incorrect:} Scenarios should input to risk analysis
			\end{itemize}
		\end{block}
	\end{frame}
	
	% ================ FINAL SLIDES ================
	
	\begin{frame}{CRISC Key Takeaways}
		\begin{block}{Essential Concepts from QAE}
			\begin{itemize}
				\item Risk identification is the first step
				\item Risk appetite = willingness to accept risk
				\item Risk tolerance = acceptable deviation from appetite
				\item Risk capacity = maximum before existence threatened
				\item Residual risk = risk remaining after controls
				\item Three lines of defense: Operations, Oversight, Audit
			\end{itemize}
		\end{block}
	\end{frame}
	
	\begin{frame}{CRISC Key Takeaways (continued)}
		\begin{block}{More Essential Concepts}
			\begin{itemize}
				\item AL E = SLE × ARO
				\item SLE = Asset Value × Exposure Factor
				\item RPO is backward-looking; RTO is forward-looking
				\item Risk response: Mitigate, Accept, Transfer, Avoid
				\item Security policy = management control
				\item KRIs = early warning signals
			\end{itemize}
		\end{block}
	\end{frame}
	
	\begin{frame}{Final Advice}
		\begin{block}{Exam Preparation}
			\centering
			\vspace{0.3cm}
			\textbf{Remember:}
			\begin{itemize}
				\item Understand the \textit{why} behind each concept
				\item Know control categories (Preventive, Detective, Corrective, Deterrent, Compensating)
				\item Understand recovery objectives (RPO, RTO, MTD)
				\item Know risk appetite vs tolerance vs capacity
				\item Practice scenario-based questions
			\end{itemize}
			\vspace{0.5cm}
			\textbf{Good Luck on Your CRISC Exam!}
		\end{block}
	\end{frame}
	
\end{document}